\documentclass[11pt,a4paper,oneside,openany]{ltjsbook}

% 空白ページの完全排除（デジタル閲覧・PDF最適化）
\let\cleardoublepage\clearpage

% 数学・物理記号パッケージ
\usepackage{amsmath,amssymb,amsfonts,amsthm,mathtools}
\usepackage{bm}

% ページレイアウト（教科書標準マージン）
\usepackage{geometry}
\geometry{
    top=25mm,
    bottom=25mm,
    left=23mm,
    right=23mm,
    headheight=28pt,
    headsep=8mm,
    footskip=13mm
}

% 表・グラフィクス
\usepackage{booktabs}
\usepackage{tabularx}
\usepackage{array}
\usepackage{xcolor}
\usepackage{graphicx}
\usepackage{enumitem}

% ハイパーリンク（控えめな学術標準：すべて黒で統一）
\usepackage{hyperref}
\hypersetup{
    colorlinks=true,
    linkcolor=black,
    citecolor=black,
    urlcolor=black,
    pdfauthor={量子コンピューティング研究グループ},
    pdftitle={量子近似最適化アルゴリズム（QAOA）体系的教科書}
}

% 装飾ボックス（tcolorbox: 質素で端正な白黒学術スタイル）
\usepackage{tcolorbox}
\tcbuselibrary{skins,breakable}

% TikZ
\usepackage{tikz}
\usetikzlibrary{shapes,arrows.meta,positioning,calc,decorations.pathreplacing,fit}

% 見出しスタイリング（titlesec: 過剰な装飾を排した標準的学術見出し）
\usepackage{titlesec}

\titleformat{\chapter}[hang]
  {\normalfont\LARGE\gtfamily\bfseries}
  {第\thechapter 章\hspace{1em}}
  {0pt}
  {}

\titleformat{name=\chapter,numberless}[hang]
  {\normalfont\LARGE\gtfamily\bfseries}
  {}
  {0pt}
  {}

\titlespacing*{\chapter}{0pt}{15pt}{25pt}

\titleformat{\section}[hang]
  {\normalfont\Large\gtfamily\bfseries}
  {\thesection\hspace{0.8em}}
  {0pt}
  {}
\titlespacing*{\section}{0pt}{14pt}{8pt}

\titleformat{\subsection}[hang]
  {\normalfont\large\gtfamily\bfseries}
  {\thesubsection\hspace{0.8em}}
  {0pt}
  {}
\titlespacing*{\subsection}{0pt}{10pt}{6pt}

\titleformat{\subsubsection}[hang]
  {\normalfont\normalsize\gtfamily\bfseries}
  {\thesubsubsection\hspace{0.7em}}
  {0pt}
  {}
\titlespacing*{\subsubsection}{0pt}{8pt}{4pt}

% ヘッダー・フッター（標準モノクロ学術スタイル：被り防止）
\usepackage{fancyhdr}
\pagestyle{fancy}
\fancyhf{}
\fancyhead[L]{\small\gtfamily\nouppercase{\leftmark}}
\fancyhead[R]{\small\thepage}
\fancyfoot[C]{}
\renewcommand{\headrulewidth}{0.4pt}
\renewcommand{\footrulewidth}{0pt}

% 章マークのフォーマット（ピリオドを排し，自然な日本語学術スタイルに統一）
\renewcommand{\chaptermark}[1]{\markboth{第\thechapter 章\hspace{1em}#1}{}}
\renewcommand{\sectionmark}[1]{}

\fancypagestyle{plain}{
  \fancyhf{}
  \fancyfoot[C]{\small\thepage}
  \renewcommand{\headrulewidth}{0pt}
}

% 目次スタイリング（章番号とタイトルの被りを防ぎつつ，装飾を排したスタイル）
\usepackage{titletoc}
\titlecontents{chapter}[0em]
  {\addvspace{0.8em}\gtfamily\bfseries}
  {\contentslabel[\thecontentslabel]{4.2em}}
  {\hspace*{-0em}}
  {\titlerule*[0.7pc]{.}\contentspage}
  [\addvspace{0.2em}]

\titlecontents{section}[4.2em]
  {\addvspace{0.15em}\normalfont}
  {\contentslabel[\thecontentslabel]{2.6em}}
  {\hspace*{-0em}}
  {\titlerule*[0.7pc]{.}\contentspage}

\titlecontents{subsection}[7.2em]
  {\addvspace{0.1em}\small}
  {\contentslabel[\thecontentslabel]{3.2em}}
  {\hspace*{-0em}}
  {\titlerule*[0.7pc]{.}\contentspage}

% -----------------------------------------------------------------------------
% 定理・定義環境
% -----------------------------------------------------------------------------
\theoremstyle{definition}
\newtheorem{definition}{定義}[chapter]
\newtheorem{example}[definition]{例}
\theoremstyle{plain}
\newtheorem{theorem}{定理}[chapter]
\newtheorem{lemma}[theorem]{補題}
\newtheorem{proposition}[theorem]{命題}
\newtheorem{corollary}[theorem]{系}
\theoremstyle{remark}
\newtheorem{remark}{注釈}[chapter]

% -----------------------------------------------------------------------------
% コラム・補足ボックス（質素・端正な白黒標準スタイル）
% -----------------------------------------------------------------------------
\tcbset{
    breakable,
    sharp corners,
    colback=white,
    colframe=black!55,
    boxrule=0.5pt,
    top=3mm, bottom=3mm, left=3.5mm, right=3.5mm,
    before skip=3.5mm, after skip=3.5mm
}

% 解説コラム
\newtcolorbox{insightbox}[1][]{
    title={\gtfamily\bfseries 【解説】 #1},
    coltitle=black,
    colbacktitle=white,
    attach title to upper,
    after title={\par\vspace{1.2mm}\hrule\vspace{2mm}}
}
\newenvironment{pointbox}[1][]{\begin{insightbox}[#1]}{\end{insightbox}}

% 原著論文データ
\newtcolorbox{paperbox}[1][]{
    title={\gtfamily\bfseries 【原論文データ】 #1},
    coltitle=black,
    colbacktitle=white,
    attach title to upper,
    after title={\par\vspace{1.2mm}\hrule\vspace{2mm}}
}

% 注意コラム
\newtcolorbox{cautionbox}[1][]{
    title={\gtfamily\bfseries 【注意】 #1},
    coltitle=black,
    colbacktitle=white,
    attach title to upper,
    after title={\par\vspace{1.2mm}\hrule\vspace{2mm}}
}
\newenvironment{alertbox}[1][]{\begin{cautionbox}[#1]}{\end{cautionbox}}

% ブラケット記法マクロ
\newcommand{\ket}[1]{|#1\rangle}
\newcommand{\bra}[1]{\langle#1|}
\newcommand{\braket}[2]{\langle#1|#2\rangle}
\newcommand{\mel}[3]{\langle#1|#2|#3\rangle}

% 参考文献タイトルの統一
\renewcommand{\bibname}{参考文献（References）}

\begin{document}

% -----------------------------------------------------------------------------
% 教科書 表紙（シンプル・標準学術仕様）
% -----------------------------------------------------------------------------
\begin{titlepage}
\centering
\vspace*{35mm}

{\huge\gtfamily\bfseries
量子近似最適化アルゴリズム（QAOA）\\[3mm]体系的教科書
}\par
\vspace{15mm}

{\Large\gtfamily
基礎数理・回路図解から原論文完全解読，XYミキサー，古典角度最適化まで
}\par

\vspace{35mm}

{\normalsize\gtfamily
対象読者：量子力学を学ぶ理工系学部生・大学院生・研究者
}\par

\vfill

{\large\gtfamily\bfseries
量子コンピューティング・最適化アルゴリズム研究グループ 編
}\par
\vspace{6mm}

{\small\color{black!70}
2026年9月（第1版） $\mid$ 学術完全収録版
}\par

\vspace{25mm}
\end{titlepage}

% -----------------------------------------------------------------------------
% 前付（Frontmatter）
% -----------------------------------------------------------------------------
\frontmatter

\chapter*{まえがき（Preface）}
\markboth{まえがき}{}
\addcontentsline{toc}{chapter}{まえがき（Preface）}

本書は，\textbf{大学3年生で量子力学の基礎講義（シュレーディンガー方程式，ブラケット記法，角運動量・スピン$1/2$，エルミート演算子，時間発展演算子，摂動論・断熱近似など）を学んでいる，あるいは履修し終えた理工系学部生}を主たる対象として執筆された，量子近似最適化アルゴリズム（Quantum Approximate Optimization Algorithm: QAOA\cite{farhi2014quantum}）の本格的な学術教科書である．

大学の量子力学の講義において，多くの学生は「水素原子のエネルギー準位」や「調和振動子の昇降演算子」「磁場中のスピン歳差運動」といった基礎物理を学ぶ．しかし，それらの物理法則が，いかにして現代最先端の「量子コンピュータ」へと姿を変え，実社会の超難問を解くアルゴリズムとして作動するのかという「生きた接続」を体感する機会は極めて少ない．

一方，情報科学側から書かれた量子計算の入門書を開くと，抽象的な「量子ビット」や「論理ゲート」の配線図が天下り式に導入され，物理学科や理工系学生が培ってきた「ハミルトニアン」「シュレーディンガー描像・ハイゼンベルク描像」「時間発展演算子 $e^{-i H t / \hbar}$」「断熱定理」「スピン鎖の交換相互作用」といった強固な物理的直感と結びつかず，学習の断絶が生じてしまう．

本書の目的は，この\textbf{「物理学の量子力学」と「計算科学の量子アルゴリズム」の間に横たわる断絶を完全に埋める}ことである：
\begin{enumerate}\setlength{\itemsep}{2pt}
    \item \textbf{スピン系から量子ゲートへの完全な翻訳}: 量子ビット $\ket{0}, \ket{1}$ はスピン$1/2$粒子の $S_z$ 固有状態 $\ket{\uparrow}, \ket{\downarrow}$ そのものであり，1量子ビット回転ゲートは磁場中のラーモア歳差運動 $e^{-i \frac{\theta}{2} \bm{n}\cdot\bm{\sigma}}$，イジング相互作用ゲート $R_{ZZ}$ は2スピン結合にほかならない\cite{nielsen2010quantum,barenco1995elementary}．
    \item \textbf{断熱量子計算（AQC）と量子アニーリング（QA）の物理数理}: 時間依存ハミルトニアンの基底状態追従，Born-Fock-Katoの量子断熱定理\cite{born1928beweis,kato1950adiabatic}，反交差とLandau-Zener遷移\cite{landau1932theorie,zener1932non}，門脇・西森の量子アニーリング（QA\cite{kadowaki1998quantum}）との深層的関係（孤立系vs開放系，計算万能性\cite{aharonov2007adiabatic}），そしてリー・トロッター分解\cite{trotter1959product,suzuki1976generalized,lloyd1996universal}という物理の王道からQAOAの誕生（Farhi et al., 2014\cite{farhi2014quantum}）を必然として導く．
    \item \textbf{原論文のハイゼンベルク描像による厳密証明}: FarhiらのLight Cone解析は，大学3年生が習う「ハイゼンベルク方程式と演算子のユニタリ時間発展 $A(t) = U^\dagger A U$」そのものであり，代数計算を一切省略せずに3正則MaxCut近似比0.6924を完全に導出する．
    \item \textbf{凝縮系物理の知恵（XYミキサー）}: One-Hot制約付き最適化において，従来のペナルティ法（Standard QAOA\cite{lucas2014ising}）が解空間の指数爆発により破綻するのに対し，Hadfield et al. (2019)\cite{hadfield2019from}やWang et al. (2020)\cite{wang2020xy}のXYミキサーは，物性物理で親しまれている「XXZスピン鎖のフリップフロップ項 $\sigma_i^+ \sigma_j^- + \sigma_i^- \sigma_j^+$」であり，全スピン $S_z^{\text{tot}}$ 保存則（ハミング重み保存）によって無効解空間を物理排除する．
    \item \textbf{古典角度最適化と機械学習サロゲートの融合}: 変分エネルギー地形の非凸性，Barren Plateau現象\cite{mcclean2018barren,wang2021noise}，パラメータシフト則\cite{mitarai2018quantum,schuld2019evaluating}，そしてFactorization Machine（FM\cite{rendle2010factorization,kitai2020designing}）との調和までを一気通貫で解説する．
\end{enumerate}

すべての記述は一次文献（査読付き原著論文・プレプリント等）に基づいて厳密に検証されており，数式をごまかさず，豊富なTikZ回路図と物理描像を交えて論理を展開している．大学3年生の皆さんが，手元のノートと鉛筆で数式を追いながら，現代の量子アルゴリズム研究の最前線へと自信を持って踏み出せる一冊となれば幸いである．

\vspace{5mm}
\noindent
\textbf{本書の構成と学習フロー}：
\begin{itemize}
    \item \textbf{基礎理論コース（第1章〜第3章）}: スピン系と回路の対応，イジング模型への変換，AQC・量子アニーリングとトロッター分解の数理．
    \item \textbf{QAOA徹底解読・古典角度最適化コース（第4章）}: 量子回路，古典ハイブリッドループ，3大ボトルネック，主要オプティマイザ，INTERP・転移学習，Farhi (2014) のMaxCut近似比0.6924の厳密証明．
    \item \textbf{制約付き最適化とXYミキサーコース（第5章〜第7章）}: ペナルティ法の破綻，Hadfield (2019) の対称性保存則，Wang (2020) のNISQ実装論．
    \item \textbf{機械学習サロゲート融合コース（第8章）}: Factorization Machine（FM）との完全調和による実用BBOパイプライン．
    \item \textbf{実力定着コース（第9章〜第10章）}: 章末演習問題，詳細解答，総括，および全36編の参考文献．
\end{itemize}

\vspace{5mm}
\begin{flushright}
2026年9月\\
著者一同
\end{flushright}

\newpage

% -----------------------------------------------------------------------------
% 数学記号・量子力学の表記法一覧
% -----------------------------------------------------------------------------
\chapter*{記号と表記法一覧（Table of Notations）}
\markboth{記号と表記法一覧}{}
\addcontentsline{toc}{chapter}{記号と表記法一覧}

\begin{table}[htbp]
\centering
\begin{tabularx}{\textwidth}{l p{5cm} X}
\toprule
\textbf{記号} & \textbf{名称・用語} & \textbf{定義・物理的意味} \\
\midrule
$\ket{0}, \ket{1}$ & 計算基底（Computational Basis） & 1量子ビットの基底ベクトル $\ket{0} = \begin{pmatrix} 1 \\ 0 \end{pmatrix}, \ket{1} = \begin{pmatrix} 0 \\ 1 \end{pmatrix}$ \\
$\ket{\psi}$ & 状態ベクトル（Ket） & 規格化されたヒルベルト空間の純粋状態ベクトル \\
$\bra{\psi}$ & 双対状態ベクトル（Bra） & $\ket{\psi}^\dagger$（エルミート共役） \\
$\braket{\phi}{\psi}$ & 内積（Inner Product） & 2つの状態ベクトルの複素内積 \\
$X, Y, Z$ & パウリ行列（Pauli Matrices） & $X = \begin{pmatrix} 0 & 1 \\ 1 & 0 \end{pmatrix}, Y = \begin{pmatrix} 0 & -i \\ i & 0 \end{pmatrix}, Z = \begin{pmatrix} 1 & 0 \\ 0 & -1 \end{pmatrix}$ \\
$H$ & アダマールゲート & $H = \frac{1}{\sqrt{2}}\begin{pmatrix} 1 & 1 \\ 1 & -1 \end{pmatrix}$（均等重ね合わせ生成） \\
$R_Z(\theta)$ & $Z$軸回転ゲート & $e^{-i \frac{\theta}{2} Z} = \begin{pmatrix} e^{-i\theta/2} & 0 \\ 0 & e^{i\theta/2} \end{pmatrix}$ \\
$R_{ZZ}(\theta)$ & 2量子ビット相互作用ゲート & $e^{-i \frac{\theta}{2} Z_i Z_j}$（イジング相互作用項） \\
$\text{CNOT}$ & 制御NOTゲート & $\ket{c, t} \to \ket{c, t \oplus c}$（エンタングルメント生成） \\
$H_C$ & コストハミルトニアン & 最適化問題の目的関数を表現する対角演算子 \\
$H_M, B$ & ミキサーハミルトニアン & 探索状態を駆動する演算子（横磁場またはXY） \\
$p$ & QAOAのレイヤー数（深さ） & コスト層とミキサー層の適用回数 \\
$\vec{\gamma}, \vec{\beta}$ & 変分パラメータ & コスト層とミキサー層の回転角度系列 \\
$[A, B]$ & 交換子（Commutator） & $AB - BA$（非可換性の尺度） \\
$\ket{W_d}$ & W状態 & $d$ 個のビットに1つだけ1が立った均等重ね合わせ状態 \\
\bottomrule
\end{tabularx}
\end{table}

\newpage

% -----------------------------------------------------------------------------
% 目次（Table of Contents）
% -----------------------------------------------------------------------------
\markboth{目次}{}
\tableofcontents

\mainmatter

% -----------------------------------------------------------------------------
% 本文（Main Chapters）
% -----------------------------------------------------------------------------
\chapter{序論と全体ロードマップ}

\section{本教材の目的と構成}
組合せ最適化問題は，現代の産業・情報社会（物流，金融，材料探索，創薬，回路設計など）に遍在する極めて重要な問題群である．しかし，Karpの21のNP完全問題をはじめとするそれらの問題群の多くはNP困難に属し，古典コンピュータでの厳密解探索は変数の増加に伴って指数関数的な計算爆発（次元の呪い）に直面する\cite{lucas2014ising}．

本教材では，以下のステップに従ってQAOAの理論体系を基礎から最先端まで段階的に修得する：
\begin{enumerate}\setlength{\itemsep}{0pt}\setlength{\parskip}{0pt}
    \item \textbf{第2章}: \textbf{量子回路入門：スピン1/2系から量子ゲートへの物理的架け橋}．大学3年生の量子力学で親しんだスピン1/2系（$\ket{\uparrow}, \ket{\downarrow}$）とパウリ行列を基礎に，ラーモア歳差運動としての回転ゲート，CNOT，およびイジング相互作用（$R_{ZZ}$）のゲート分解\cite{nielsen2010quantum,barenco1995elementary}を学ぶ．
    \item \textbf{第3章}: 組合せ最適化を物理のスピン系（イジング模型/QUBO）へマッピングする数学的基盤\cite{lucas2014ising}と，QAOAの直接の先祖である断熱量子計算（AQC）の物理原理\cite{farhi2000quantum,farhi2001quantum,born1928beweis,kato1950adiabatic,jansen2007bounds}，門脇・西森の量子アニーリング（QA）\cite{kadowaki1998quantum}との深層的関係，およびトロッター分解\cite{trotter1959product,suzuki1976generalized,lloyd1996universal}を修得する．
    \item \textbf{第4章}: \textbf{QAOA原論文の完全徹底解読と古典角度最適化}．Farhi et al. (2014) の原論文\cite{farhi2014quantum}を徹底解読し，コスト演算子・横磁場ミキサー・変分ハイブリッドループから，古典角度最適化の3大ボトルネック，主要オプティマイザ（COBYLA vs SPSA vs パラメータシフト則），先端戦略（INTERP・転移学習），そして3正則グラフMaxCutの近似比導出までを一気通貫で修得する．
    \item \textbf{第5章}: 現実の産業問題で頻出するOne-Hot制約などの等式制約に対し，従来のペナルティ法（Standard QAOA）がなぜ解空間の指数消滅により壊滅的な破綻を迎えるのかを解明する\cite{lucas2014ising,choi2008different,hadfield2019from}．
    \item \textbf{第6章}: \textbf{Hadfield et al. (2019) の原著論文}\cite{hadfield2019from}を解読し，探索空間を有効解のみに閉じる「Quantum Alternating Operator Ansatz」とXYミキサーの理論，およびW状態生成\cite{bartschi2019deterministic}を学ぶ．
    \item \textbf{第7章}: \textbf{Wang et al. (2020)}\cite{wang2020xy}の解析結果に基づき，XYミキサーの回路分解，パラメータ周期性，偶奇スワップネットワーク\cite{kivlichan2018quantum}によるNISQコンパイルを学ぶ．
    \item \textbf{第8章}: Factorization Machine（FM\cite{rendle2010factorization}）サロゲートモデルとXYミキサーを融合させた実用ブラックボックス最適化（BBO\cite{kitai2020designing}）パイプラインへの接続を解説する．
    \item \textbf{第9章}: 理解度を定着させるための演習問題と詳細解答を示す．
    \item \textbf{第10章}: 本書の要点総括と今後の展望を示す．
\end{enumerate}

\chapter{量子回路入門：スピン1/2系から量子ゲートへの物理的架け橋}
\label{sec:quantum_circuits_intro}

大学の量子力学講義で学んだように，スピン$1/2$粒子（電子や中性子など）のヒルベルト空間は2次元複素ベクトル空間 $\mathbb{C}^2$ であり，$S_z$ 固有基底 $\{\ket{\uparrow_z}, \ket{\downarrow_z}\}$ によって張られる．
量子情報科学における「量子ビット（Qubit）」とは，まさにこのスピン$1/2$系を情報担体として捉え直した物理的実体そのものである\cite{nielsen2010quantum}．

本章では，大学3年生が親しんできた「スピンのハミルトニアン」「パウリ行列」「ユニタリ時間発展演算子 $e^{-i H t / \hbar}$」から出発し，それらがいかにして「量子ゲート」や「量子回路配線図」へと翻訳されるのか，その物理的対応関係を解き明かす\cite{barenco1995elementary}．

\section{量子ビット (Qubit) と重ね合わせの直感的イメージ}
古典コンピュータのビットは $0$ または $1$ のどちらか一方の状態しかとれない（コインの表か裏）．
これに対し，\textbf{量子ビット（Qubit）}は，$0$ と $1$ の両方の可能性を同時にあわせ持つことができる（\textbf{量子力学的重ね合わせ: Superposition}）．
数学的には，規格化された2次元複素ベクトルとして表される：
\begin{equation}
    \ket{\psi} = \alpha \ket{0} + \beta \ket{1} = \begin{pmatrix} \alpha \\ \beta \end{pmatrix}, \quad |\alpha|^2 + |\beta|^2 = 1
\end{equation}
ここで $\ket{0} = \begin{pmatrix} 1 \\ 0 \end{pmatrix}, \ket{1} = \begin{pmatrix} 0 \\ 1 \end{pmatrix}$ は基準となる基底（計算基底）である．
この量子ビットを観測（$Z$ 基底での射影測定）すると，$|\alpha|^2$ の確率で $0$ が得られ，$|\beta|^2$ の確率で $1$ が得られる（ボルンの規則）．

\begin{insightbox}[量子力学（スピン$1/2$系）との物理的対応関係の完全整理]
大学の量子力学講義で学ぶスピン$1/2$粒子（電子や中性子など）の物理と，量子情報科学における量子ビットは，以下の通り数理的・物理的に完全に1対1対応している：

\begin{enumerate}
    \item \textbf{状態空間とブロッホ球の幾何学}:
    スピン$1/2$粒子の状態空間は2次元ヒルベルト空間 $\mathbb{C}^2$ であり，$S_z$ 固有基底 $\{\ket{\uparrow_z}, \ket{\downarrow_z}\}$ が量子ビットの計算基底 $\{\ket{0}, \ket{1}\}$ に対応する：
    \begin{equation}
        \ket{0} \equiv \ket{\uparrow_z} = \begin{pmatrix} 1 \\ 0 \end{pmatrix}, \quad \ket{1} \equiv \ket{\downarrow_z} = \begin{pmatrix} 0 \\ 1 \end{pmatrix}
    \end{equation}
    任意の一粒子純粋状態は，規格化条件と全位相の自由度を除くと，極角 $\theta \in [0, \pi]$ と方位角 $\phi \in [0, 2\pi)$ を用いて：
    \begin{equation}
        \ket{\psi} = \cos\left(\frac{\theta}{2}\right)\ket{0} + e^{i\phi}\sin\left(\frac{\theta}{2}\right)\ket{1}
    \end{equation}
    と表される．この状態におけるパウリ行列 $\bm{\sigma} = (X, Y, Z)$ の期待値ベクトル（スピン偏極ベクトル）を計算すると：
    \begin{equation}
        \bm{P} = \langle \bm{\sigma} \rangle = (\mel{\psi}{X}{\psi}, \; \mel{\psi}{Y}{\psi}, \; \mel{\psi}{Z}{\psi})^\top = (\sin\theta \cos\phi, \; \sin\theta \sin\phi, \; \cos\theta)^\top
    \end{equation}
    となり，これはまさに三次元実空間の単位球面（\textbf{ブロッホ球: Bloch Sphere}）上の方向余弦ベクトルそのものである！

    \item \textbf{ゼーマン相互作用とラーモア歳差運動（1量子ビットゲートの物理）}:
    スピン角運動量演算子 $\bm{S} = \frac{\hbar}{2}\bm{\sigma}$（磁気回転比 $\gamma$）を持つスピンに対し，外部静磁場 $\bm{B} = B \bm{n}$（$\bm{n}$ は単位ベクトル）を印加したときのゼーマン相互作用ハミルトニアンは：
    \begin{equation}
        \mathcal{H} = -\bm{\mu}\cdot\bm{B} = -\gamma B \bm{n}\cdot\bm{S} = \frac{\hbar \omega_L}{2} \bm{n}\cdot\bm{\sigma} \quad (\omega_L \equiv -\gamma B \text{ はラーモア周波数})
    \end{equation}
    である．ハイゼンベルク方程式 $\frac{d}{dt}\bm{S}(t) = \frac{i}{\hbar}[\mathcal{H}, \bm{S}(t)]$ を計算すると：
    \begin{equation}
        \frac{d}{dt}\bm{S}(t) = \gamma \bm{S}(t) \times \bm{B}
    \end{equation}
    が得られ，スピンベクトルが磁場軸 $\bm{n}$ の周りを角周波数 $\omega_L$ で歳差運動（ラーモア歳差運動）することが示される．
    シュレーディンガー描像における時間発展演算子は，$(\bm{n}\cdot\bm{\sigma})^2 = I$ というパウリ行列の反交換関係 $\{ \sigma_a, \sigma_b \} = 2\delta_{ab} I$ に着目してテイラー展開の偶数次・奇数次を分離することで：
    \begin{align}
        U(t) &= \exp\left(-\frac{i}{\hbar} \mathcal{H} t\right) = \exp\left(-i \frac{\omega_L t}{2} \bm{n}\cdot\bm{\sigma}\right) \nonumber \\
        &= \sum_{m=0}^\infty \frac{(-1)^m (\omega_L t / 2)^{2m}}{(2m)!} I - i \sum_{m=0}^\infty \frac{(-1)^m (\omega_L t / 2)^{2m+1}}{(2m+1)!} (\bm{n}\cdot\bm{\sigma}) \nonumber \\
        &= \cos\left(\frac{\omega_L t}{2}\right) I - i \sin\left(\frac{\omega_L t}{2}\right) (\bm{n}\cdot\bm{\sigma})
    \end{align}
    と厳密に求まる．回転角を $\theta = \omega_L t$ と置けば，これはまさに量子計算における\textbf{「1量子ビット回転ゲート $R_{\bm{n}}(\theta)$」}そのものである！
    すなわち，\textbf{量子回路における1量子ビットゲートとは，電磁場パルスによってスピンを所望の時間だけラーモア歳差運動させる物理過程}にほかならない．

    \item \textbf{2スピン結合とイジングゲート $R_{ZZ}(\theta)$（QAOAコスト層の物理）}:
    2つの空間的に近接したスピン$1/2$粒子間に働くスピン結合ハミルトニアン（一軸異方性イジング相互作用）は：
    \begin{equation}
        \mathcal{H}_{\text{Ising}} = J S_{1z} S_{2z} = \frac{\hbar J}{4} Z_1 Z_2
    \end{equation}
    である．このハミルトニアンに従う自然な時間発展演算子は $U(t) = \exp\left(-i \frac{J t}{4} Z_1 Z_2\right)$ であり，回転角 $\theta = \frac{J t}{2}$ と同定すれば，まさに\textbf{QAOAのコスト層の主役である2量子ビット相互作用ゲート $R_{ZZ}(\theta) = e^{-i \frac{\theta}{2} Z_1 Z_2}$ そのもの}となる．
\end{enumerate}
\end{insightbox}

\section{量子回路図の読み方}
量子回路図は，時間の流れとともに量子ビットの状態がどう加工されていくかを表す配線図である：
\begin{itemize}
    \item \textbf{水平方向の線（ワイヤー）}: 1つの量子ビットを表す．時間は\textbf{左から右へ}進む．
    \item \textbf{四角い箱（ゲート）}: 量子ビットの状態を回転・変換する操作（ユニタリ行列 $U$）．
    \item \textbf{縦に繋がる線（2量子ビットゲート）}: 2つの量子ビットの間にもつれ（エンタングルメント）や相互作用を発生させる操作．
    \item \textbf{メーター記号（測定）}: 量子重ね合わせを観測し，古典的な二値情報（0 または 1）を取り出す操作．
\end{itemize}

\section{QAOAを理解するための必須ゲート一覧}

\subsection{1. 1量子ビット基本ゲート}
\begin{figure}[htbp]
\centering
\begin{tikzpicture}[scale=0.95, >=latex]
    % X Gate
    \node[left] at (0.2, 2.5) {$\ket{\psi}$};
    \draw[thick] (0.2, 2.5) -- (2.8, 2.5) node[right] {$X\ket{\psi}$};
    \draw[fill=blue!12, draw=blue!75!black, thick] (1.0, 2.05) rectangle (2.0, 2.95) node[midway, font=\bfseries\large] {$X$};
    \node[below, font=\footnotesize, text width=3.2cm, align=center] at (1.5, 1.8) {NOTゲート\\($\ket{0} \leftrightarrow \ket{1}$ 反転)};

    % H Gate
    \node[left] at (5.4, 2.5) {$\ket{0}$};
    \draw[thick] (5.4, 2.5) -- (7.8, 2.5) node[right] {$\ket{+}$};
    \draw[fill=green!12, draw=green!65!black, thick] (6.1, 2.05) rectangle (7.1, 2.95) node[midway, font=\bfseries\large] {$H$};
    \node[below, font=\footnotesize, text width=3.4cm, align=center] at (6.6, 1.8) {アダマールゲート\\($\ket{+} = \frac{\ket{0}+\ket{1}}{\sqrt{2}}$ 生成)};

    % Rz Gate
    \node[left] at (10.6, 2.5) {$\ket{\psi}$};
    \draw[thick] (10.6, 2.5) -- (13.0, 2.5) node[right] {$R_Z(\theta)\ket{\psi}$};
    \draw[fill=orange!12, draw=orange!75!black, thick] (11.3, 2.05) rectangle (12.3, 2.95) node[midway, font=\bfseries\small] {$R_Z(\theta)$};
    \node[below, font=\footnotesize, text width=3.4cm, align=center] at (11.8, 1.8) {Z回転ゲート\\(位相差 $e^{-i\theta/2}$ 付加)};
\end{tikzpicture}
\caption{代表的な1量子ビットゲート．$H$ は重ね合わせを作り，$X$ はビットを反転させ，$R_Z, R_X$ はブロッホ球上で状態を角度 $\theta$ だけ連続回転させる．}
\label{fig:single_qubit_gates}
\end{figure}

\subsection{2. 制御NOTゲート (CNOT / CX)}
2量子ビットゲートの代表格が \textbf{CNOT（Controlled-NOT）ゲート} である．
「制御ビット（黒丸）」が $1$ のときだけ，「標的ビット（十字丸）」のビットを反転（$X$ を作用）させる：
\begin{figure}[htbp]
\centering
\begin{tikzpicture}[scale=1.0, >=latex]
    % Control wire
    \node[left] at (-0.2, 1) {$\ket{c}$};
    \draw[thick] (-0.2, 1) -- (2.6, 1) node[right] {$\ket{c}$};
    % Target wire
    \node[left] at (-0.2, 0) {$\ket{t}$};
    \draw[thick] (-0.2, 0) -- (2.6, 0) node[right] {$\ket{t \oplus c}$};
    
    % Control dot
    \filldraw[black] (1.3, 1) circle (3pt);
    % Vertical line
    \draw[thick] (1.3, 1) -- (1.3, 0);
    % Target circle with cross
    \draw[thick, fill=white] (1.3, 0) circle (6pt);
    \draw[thick] (1.3, -0.21) -- (1.3, 0.21);
    \draw[thick] (1.09, 0) -- (1.51, 0);
    
    % Explanatory Truth Table
    \node[right, font=\small, text width=8cm] at (4.6, 0.5) {
        $\ket{00} \to \ket{00}$ \quad (制御=0: 標的不変) \\
        $\ket{01} \to \ket{01}$ \quad (制御=0: 標的不変) \\
        $\ket{10} \to \ket{11}$ \quad (\textbf{制御=1: 標的反転}) \\
        $\ket{11} \to \ket{10}$ \quad (\textbf{制御=1: 標的反転})
    };
\end{tikzpicture}
\caption{制御NOT（CNOT）ゲートの構造と作用．2つの量子ビット間に量子力学的なエンタングルメント（もつれ）を生み出す最重要ゲート．}
\label{fig:cnot_gate}
\end{figure}

\section{コスト層の主役：$R_{ZZ}(\theta)$ 相互作用ゲートの分解}
QAOAのコスト層では，イジング相互作用項 $e^{-i \frac{\theta}{2} Z_i Z_j}$（2つの変数の積 $x_i x_j$ に相当）が大量に登場する．
物理的な量子コンピュータの多くは2量子ビットの直接結合としてCNOTゲートしか持たないが，\textbf{「CNOT で挟んで間に $R_Z(\theta)$ を挟む」} だけで，このイジング相互作用が完璧に実現できる（図\ref{fig:rzz_decomposition}）．

\begin{figure}[htbp]
\centering
\begin{tikzpicture}[scale=1.0, >=latex]
    % Left side: Composite Rzz gate
    \node at (-0.8, 1) {$\ket{q_i}$};
    \node at (-0.8, 0) {$\ket{q_j}$};
    \draw[thick] (-0.3, 1) -- (2.5, 1);
    \draw[thick] (-0.3, 0) -- (2.5, 0);
    \draw[fill=purple!10, draw=purple!70!black, thick] (0.3, -0.3) rectangle (1.9, 1.3) node[midway, font=\bfseries] {$R_{ZZ}(\theta)$};

    % Equals sign
    \node[font=\Huge] at (3.2, 0.5) {$=$};

    % Right side: Decomposed circuit
    \node at (3.8, 1) {$\ket{q_i}$};
    \node at (3.8, 0) {$\ket{q_j}$};
    \draw[thick] (4.2, 1) -- (9.8, 1);
    \draw[thick] (4.2, 0) -- (9.8, 0);

    % CNOT 1
    \filldraw[black] (5.2, 1) circle (3pt);
    \draw[thick] (5.2, 1) -- (5.2, 0);
    \draw[thick, fill=white] (5.2, 0) circle (5pt);
    \draw[thick] (5.2, -0.18) -- (5.2, 0.18);
    \draw[thick] (5.02, 0) -- (5.38, 0);

    % Rz gate
    \draw[fill=orange!10, draw=orange!70!black, thick] (6.4, -0.4) rectangle (7.6, 0.4) node[midway, font=\small\bfseries] {$R_Z(\theta)$};

    % CNOT 2
    \filldraw[black] (8.8, 1) circle (3pt);
    \draw[thick] (8.8, 1) -- (8.8, 0);
    \draw[thick, fill=white] (8.8, 0) circle (5pt);
    \draw[thick] (8.8, -0.18) -- (8.8, 0.18);
    \draw[thick] (8.62, 0) -- (8.98, 0);
\end{tikzpicture}
\caption{2量子ビットイジング相互作用ゲート $R_{ZZ}(\theta) = \exp\left(-i \frac{\theta}{2} Z_i Z_j\right)$ の標準的ゲート分解．最初のCNOTでパリティ（2ビットの排他的論理和）を標的に集約し，$R_Z$ で位相を回転させた後，2つ目のCNOTで標的の状態を元に戻す．}
\label{fig:rzz_decomposition}
\end{figure}

この基本ブロックを理解していれば，後述するQAOAの回路全体が一目で「何をしているか」手に取るように読めるようになる．

\chapter{数学的基礎：組合せ最適化・イジング模型・断熱量子計算}

\section{QUBOからイジング模型への変換}
二値変数 $x_i \in \{0, 1\}$ を決定変数とする二次制約なし二値最適化（Quadratic Unconstrained Binary Optimization: QUBO）は以下で表される\cite{lucas2014ising}：
\begin{equation}
    f(x) = \sum_{i=1}^N q_i x_i + \sum_{i < j}^N Q_{ij} x_i x_j = x^\top Q x
\end{equation}
古典二値変数 $x_i \in \{0, 1\}$ とスピン変数 $s_i \in \{+1, -1\}$ の間には，以下の全単射なアフィン変換が存在する：
\begin{equation}
    x_i = \frac{1 - s_i}{2} \quad \Longleftrightarrow \quad s_i = 1 - 2x_i
\end{equation}
これを代入すると，物理学における古典イジング模型のハミルトニアンが得られる：
\begin{equation}
    H_{\text{Ising}}(s) = \sum_{i < j} J_{ij} s_i s_j + \sum_i h_i s_i + C_0
\end{equation}
ここで結合定数は $J_{ij} = \frac{1}{4} Q_{ij}$ であり，局所磁場は $h_i = -\frac{1}{2} q_i - \frac{1}{4}\sum_{j \neq i} Q_{ij}$ である．

\section{パウリ行列と量子ハミルトニアン}
量子系では，古典変数 $s_i \in \{+1, -1\}$ を1量子ビットのエルミート演算子であるパウリ $Z$ 行列 $Z_i$ に昇格させる：
\begin{equation}
    Z = \begin{pmatrix} 1 & 0 \\ 0 & -1 \end{pmatrix}, \quad 
    X = \begin{pmatrix} 0 & 1 \\ 1 & 0 \end{pmatrix}, \quad 
    Y = \begin{pmatrix} 0 & -i \\ i & 0 \end{pmatrix}
\end{equation}
計算基底 $\ket{0}, \ket{1}$ は $Z$ の固有状態である：$Z\ket{0} = +1\ket{0}, \; Z\ket{1} = -1\ket{1}$．
系全体のコストハミルトニアン $H_C$ は以下のように定義される：
\begin{equation}
    H_C = \sum_{i < j} J_{ij} Z_i Z_j + \sum_i h_i Z_i + C_0 I
\end{equation}
$H_C$ は計算基底において完全に対角行列であり，任意の古典ビット列 $x \in \{0, 1\}^N$ に対応する基底状態 $\ket{x} = \ket{x_1 x_2 \dots x_N}$ に対し，
\begin{equation}
    H_C \ket{x} = f(x) \ket{x}
\end{equation}
が厳密に成立する．したがって，\textbf{目的関数 $f(x)$ を最小化することは，$H_C$ の基底状態（最小エネルギー固有状態）を求めることと完全に同値}である．

\section{断熱量子計算 (AQC) の背景・物理原理・限界}
\label{sec:aqc_detailed}

QAOAを真に深く理解するためには，その直接の先祖である \textbf{断熱量子計算（Adiabatic Quantum Computation: AQC）} の誕生背景と，物理学的な作動原理，そして直面した壁を理解することが不可欠である．

\subsection{1. なぜ AQC が考案されたのか？（歴史的背景）}
1994年にショア（Peter Shor）が素因数分解の多項式時間量子アルゴリズムを，1996年にグローバー（Lov Grover）が非構造化データ探索の2次加速アルゴリズムを発表し，量子コンピュータの驚異的な潜在能力が世界に衝撃を与えた．
しかし，これらは「周期性の検出（量子フーリエ変換）」や「ブラックボックス判定オラクル」という極めて特殊な代数的・数学的構造を持つ問題に特化していた．
産業・社会に無数に存在する，巡回セールスマン問題（TSP）やグラフ彩色，充足可能性問題（SAT）のような\textbf{「一般的なNP困難な組合せ最適化問題」を解く普遍的な量子アプローチは存在しなかった}．

この膠着状態を打破すべく，MITの理論物理学者 Edward Farhi, Jeffrey Goldstone, Sam Gutmann, Michael Sipser（2000）が提唱したのが断熱量子計算（Adiabatic Quantum Computation: AQC）である\cite{farhi2000quantum,farhi2001quantum}．彼らの発想の原点は極めて直感的であった：
\begin{quote}
\textit{「解きたい問題の基底状態（大域的最適解）をゼロから直接作るのは絶望的に難しい．\\
しかし，\textbf{自明なハミルトニアンの基底状態を作るのは極めて簡単}である．\\
系を自明な基底状態からスタートさせ，外場をゆっくり連続的に変化させて解きたい問題へと変形させれば，物理法則（断熱定理）により，\textbf{量子系は勝手に真の最適解へと運ばれる}のではないか？」}
\end{quote}

\subsection{2. 量子断熱定理（Quantum Adiabatic Theorem）の物理と数学}
時間に依存するハミルトニアン $H(t)$ に従う系の時間発展は，時間依存シュレーディンガー方程式で記述される：
\begin{equation}
    i \hbar \frac{d}{dt} \ket{\psi(t)} = H(t) \ket{\psi(t)}
\end{equation}
各時刻 $t$ における瞬時固有状態方程式を $H(t)\ket{n(t)} = E_n(t)\ket{n(t)}$ とする（$E_0(t) \le E_1(t) \le \dots$）．
ボルンとフォック（Born \& Fock, 1928）\cite{born1928beweis}および加藤敏夫（Kato, 1950）\cite{kato1950adiabatic}によって厳密化された量子断熱定理は以下を主張する：

\begin{theorem}[量子断熱定理（Born \& Fock, 1928; Kato, 1950）]
系が時刻 $t=0$ において初期ハミルトニアン $H(0)$ の基底状態 $\ket{0(0)}$ に準備されているとする．
全実行時間 $T$ に対し，ハミルトニアン $H(t)$ の時間変化が十分にゆっくりであれば，終状態 $\ket{\psi(T)}$ は時刻 $T$ における目的ハミルトニアン $H(T)$ の基底状態 $\ket{0(T)}$ に極めて高い確率（漸近的に確率1）で留まり続ける．
\end{theorem}

量子力学や断熱量子計算を学ぶ読者にとって，\textbf{「なぜゆっくり変化させるだけで基底状態に留まるのか？」「なぜ断熱条件の分母にエネルギーギャップの2乗 $\Delta_{\min}^2$ が現れるのか？」}という疑問は，物理的・数学的にも最も探究心をそそる重要論点である．
多くの教科書では結果の不等式のみが天下り的に紹介されることが多いが，ここでは時間依存シュレーディンガー方程式から出発し，\textbf{一切の論理ジャンプなしにステップ・バイ・ステップで量子断熱条件を厳密に導出}する．

\begin{insightbox}[【大学3年生向け徹底導出】シュレーディンガー方程式からの量子断熱条件の完全導出]
\begin{enumerate}[label=\textbf{Step \arabic* :}, leftmargin=*]
    \item \textbf{瞬時固有基底の導入と動的位相の分離}:\\
    各時刻 $t \in [0, T]$ における瞬時ハミルトニアン $H(t)$ の固有方程式を：
    \begin{equation}
        H(t) \ket{n(t)} = E_n(t) \ket{n(t)} \quad (E_0(t) < E_1(t) \le E_2(t) \le \dots)
    \end{equation}
    とする．全発展路において基底状態は孤立（エネルギースペクトルギャップ $\Delta(t) = E_1(t) - E_0(t) > 0$）しているものとする．固有状態は正規直交完全系（$\braket{m(t)}{n(t)} = \delta_{mn}$）を張る．
    一般の状態ベクトル $\ket{\psi(t)}$ を瞬時固有基底で展開する際，静的ハミルトニアンの時間発展因子 $e^{-i E_n t / \hbar}$ の自然な拡張として，\textbf{動的位相（Dynamic Phase）}：
    \begin{equation}
        \theta_n(t) \equiv -\frac{1}{\hbar} \int_0^t E_n(t') dt'
    \end{equation}
    をあらかじめ位相因子として括り出し，以下のように展開する：
    \begin{equation}
        \ket{\psi(t)} = \sum_n c_n(t) e^{i \theta_n(t)} \ket{n(t)}
    \end{equation}
    初期条件は，時刻 $t=0$ で系が基底状態 $\ket{0(0)}$ にあるため，$c_0(0) = 1$ かつ $n \ge 1$ に対し $c_n(0) = 0$ である．

    \item \textbf{シュレーディンガー方程式への代入と動的位相の厳密相殺}:\\
    時間依存シュレーディンガー方程式 $i\hbar \frac{d}{dt}\ket{\psi(t)} = H(t)\ket{\psi(t)}$ の左辺に対し，積の微分法を適用すると：
    \begin{align}
        i\hbar \frac{d}{dt}\ket{\psi(t)} &= i\hbar \sum_n \Bigl[ \dot{c}_n(t) e^{i\theta_n(t)} \ket{n(t)} + c_n(t) (i \dot{\theta}_n(t)) e^{i\theta_n(t)} \ket{n(t)} \nonumber \\
        &\qquad\qquad + c_n(t) e^{i\theta_n(t)} \ket{\dot{n}(t)} \Bigr] \nonumber \\
        &= i\hbar \sum_n \dot{c}_n(t) e^{i\theta_n(t)} \ket{n(t)} + \sum_n E_n(t) c_n(t) e^{i\theta_n(t)} \ket{n(t)} \nonumber \\
        &\quad + i\hbar \sum_n c_n(t) e^{i\theta_n(t)} \ket{\dot{n}(t)}
    \end{align}
    となる（ここで $\dot{\theta}_n(t) = -\frac{E_n(t)}{\hbar}$ を代入した）．
    一方，右辺はハミルトニアンの固有方程式より：
    \begin{equation}
        H(t) \ket{\psi(t)} = \sum_n c_n(t) e^{i\theta_n(t)} H(t)\ket{n(t)} = \sum_n E_n(t) c_n(t) e^{i\theta_n(t)} \ket{n(t)}
    \end{equation}
    である．両辺の第2項（エネルギー固有値に起因する項）が完全に一致して相殺するため，以下の方程式が厳密に得られる：
    \begin{equation}
        i\hbar \sum_n \left[ \dot{c}_n(t) \ket{n(t)} + c_n(t) \ket{\dot{n}(t)} \right] e^{i\theta_n(t)} = 0
    \end{equation}

    \item \textbf{固有ブラ $\bra{m(t)}$ への射影とベリー位相}:\\
    上式の両辺に左から任意の固有状態 $\bra{m(t)}$ を掛け，基底の直交性 $\braket{m(t)}{n(t)} = \delta_{mn}$ を適用する：
    \begin{equation}
        \dot{c}_m(t) e^{i\theta_m(t)} + \sum_n c_n(t) e^{i\theta_n(t)} \braket{m(t)}{\dot{n}(t)} = 0
    \end{equation}
    両辺に $e^{-i\theta_m(t)}$ を掛けて $\dot{c}_m(t)$ について整理すると：
    \begin{equation}
        \dot{c}_m(t) = - c_m(t) \braket{m(t)}{\dot{m}(t)} - \sum_{n \neq m} c_n(t) e^{i(\theta_n(t) - \theta_m(t))} \braket{m(t)}{\dot{n}(t)}
    \end{equation}
    ここで対角項 $\braket{m(t)}{\dot{m}(t)}$ は，規格化条件 $\braket{m(t)}{m(t)}=1$ の時間微分 $\frac{d}{dt}\braket{m}{m} = 2\text{Re}[\braket{m}{\dot{m}}] = 0$ より純虚数である．
    この対角項は，固有基底の局所位相再定義（ゲージ変換 $\ket{m(t)} \to e^{i\gamma_m(t)}\ket{m(t)}$）によって\textbf{ベリー位相（Berry Phase）}\cite{berry1984quantal} $\gamma_m(t) = i \int_0^t \braket{m(t')}{\dot{m}(t')}dt'$ として完全に吸収できる．
    したがって，状態間の非断熱遷移（励起）を司る本質は，\textbf{非対角項 $\braket{m(t)}{\dot{n}(t)}$ ($n \neq m$)} のみである．

    \item \textbf{基本恒等式 $\braket{m}{\dot{n}} = \frac{\mel{m}{\dot{H}}{n}}{E_n - E_m}$ の厳密な代数証明}:\\
    非対角内積 $\braket{m(t)}{\dot{n}(t)}$ とハミルトニアンの時間変化率 $\dot{H}(t)$ の関係を明らかにするため，瞬時固有方程式 $H(t)\ket{n(t)} = E_n(t)\ket{n(t)}$ の両辺を時間 $t$ で微分する：
    \begin{equation}
        \dot{H}(t)\ket{n(t)} + H(t)\ket{\dot{n}(t)} = \dot{E}_n(t)\ket{n(t)} + E_n(t)\ket{\dot{n}(t)}
    \end{equation}
    左から $m \neq n$ なる固有状態 $\bra{m(t)}$ を掛けると，$H(t)$ のエルミート性 $\mel{m(t)}{H(t)}{\dot{n}(t)} = E_m(t)\braket{m(t)}{\dot{n}(t)}$ および直交性 $\braket{m(t)}{n(t)}=0$ より：
    \begin{equation}
        \mel{m(t)}{\dot{H}(t)}{n(t)} + E_m(t)\braket{m(t)}{\dot{n}(t)} = E_n(t)\braket{m(t)}{\dot{n}(t)}
    \end{equation}
    同類項をまとめると $(E_n(t) - E_m(t))\braket{m(t)}{\dot{n}(t)} = \mel{m(t)}{\dot{H}(t)}{n(t)}$ となる．エネルギー準位の縮退がない（$E_n(t) \neq E_m(t)$）ため，両辺をエネルギー差で割ることで，\textbf{一切の近似なしに厳密に}以下の基本恒等式が得られる：
    \begin{equation}
        \braket{m(t)}{\dot{n}(t)} = \frac{\mel{m(t)}{\dot{H}(t)}{n(t)}}{E_n(t) - E_m(t)} \quad (m \neq n)
    \end{equation}

    \item \textbf{部分積分と「ギャップの2乗 $\Delta(t)^2$」の数学的必然性}:\\
    上記恒等式を係数の運動方程式に代入し，断熱極限における最低次の摂動論（Born近似 $c_0(t) \approx 1, c_{m\ge 1}(t) \approx 0$）を適用すると，基底状態から第1励起状態への遷移振幅 $c_1(T)$（$m=1, n=0$）は：
    \begin{equation}
        c_1(T) \approx \int_0^T \frac{\mel{1(t)}{\dot{H}(t)}{0(t)}}{\Delta(t)} \exp\left(\frac{i}{\hbar}\int_0^t \Delta(t') dt'\right) dt
    \end{equation}
    と表される（ここで $\Delta(t) = E_1(t) - E_0(t) > 0$）．
    ここで被積分関数の量子振動因子に対し，$\frac{d}{dt}\left[\exp\left(\frac{i}{\hbar}\int_0^t \Delta dt'\right)\right] = \frac{i}{\hbar}\Delta(t) \exp\left(\frac{i}{\hbar}\int_0^t \Delta dt'\right)$ を用いて\textbf{部分積分（Integration by Parts）}を実行する：
    \begin{align}
        c_1(T) &\approx \int_0^T \frac{\mel{1(t)}{\dot{H}(t)}{0(t)}}{\Delta(t)} \cdot \left[ \frac{\hbar}{i\Delta(t)} \frac{d}{dt} e^{\frac{i}{\hbar}\int_0^t \Delta(t') dt'} \right] dt \nonumber \\
        &= \left[ \frac{\hbar}{i} \frac{\mel{1(t)}{\dot{H}(t)}{0(t)}}{\Delta(t)^2} e^{\frac{i}{\hbar}\int_0^t \Delta(t') dt'} \right]_0^T \nonumber \\
        &\quad - \frac{\hbar}{i}\int_0^T \frac{d}{dt}\left( \frac{\mel{1(t)}{\dot{H}(t)}{0(t)}}{\Delta(t)^2} \right) e^{\frac{i}{\hbar}\int_0^t \Delta(t') dt'} dt
    \end{align}
    \textbf{ここにギャップの2乗 $\Delta(t)^2$ が現れる数学的必然性がある！}
    ハミルトニアンの変化に伴う固有状態の混合率（$\sim 1/\Delta$）に対し，シュレーディンガー方程式特有の高速位相振動を時間積分する際にさらに位相微分（$\sim \hbar/\Delta$）で割られるため，分母に必然的に2乗が現れるのである．

    \item \textbf{断熱タイムスケール $T$ と量子断熱条件の導出}:\\
    全実行時間を $T$ とし，正規化無次元時間 $s = t/T \in [0, 1]$ を導入すると，連鎖律より $\dot{H}(t) = \frac{1}{T}\frac{dH}{ds}$ とスケールする．
    したがって遷移振幅 $c_1(T)$ の漸近的上限オーダーは：
    \begin{equation}
        |c_1(T)| \le \mathcal{O}\left( \frac{\hbar}{T} \frac{\max_{s \in [0, 1]} \left| \mel{1(s)}{\frac{dH}{ds}}{0(s)} \right|}{\min_{s \in [0, 1]} \Delta(s)^2} \right)
    \end{equation}
    と評価される．
    第1励起状態への遷移確率 $P_{0 \to 1} = |c_1(T)|^2 \ll 1$ に抑えて断熱性を保つための条件は，振幅 $|c_1(T)| \ll 1$ である．
    これより，全実行時間 $T$ に対する\textbf{量子断熱条件（Jansen et al., 2007\cite{jansen2007bounds}）}：
    \begin{equation}
        T \gg \frac{\hbar \max_{s \in [0, 1]} \left| \mel{1(s)}{\frac{dH}{ds}}{0(s)} \right|}{\Delta_{\min}^2}
    \end{equation}
    が数学的飛躍なく完全に導出される．ここで $\Delta_{\min} = \min_{s \in [0, 1]} \Delta(s)$ は発展路上の最小エネルギーギャップである．
\end{enumerate}
\end{insightbox}

\subsection{3. 反交差（Avoided Crossing）とランダウ・ツェナー遷移}
なぜ分母にギャップの2乗 $\Delta_{\min}^2$ が現れるのか？その物理的実態は，量子力学における準位の反交差（Avoided Level Crossing）とランダウ・ツェナー遷移（Landau-Zener Transition）\cite{landau1932theorie,zener1932non}にある（図\ref{fig:avoided_crossing}）．

\begin{insightbox}[【大学3年生向け徹底解説】2準位系ランダウ・ツェナー模型の物理と数理解析]
反交差近傍における量子ダイナミクスは，大学3年生の量子力学講義や演習問題で頻出する標準的な\textbf{「時間依存2準位系ランダウ・ツェナー（Landau-Zener: LZ）模型」}\cite{landau1932theorie,zener1932non}として完全に定式化できる．本項では，透熱基底と断熱基底の対比からツェナー公式の物理的意味までを余すところなく解読する．

\begin{enumerate}[label=\textbf{(\arabic*)}, leftmargin=*]
    \item \textbf{模型のハミルトニアンと透熱基底（Diabatic Basis）}:\\
    局所的な2つの直交状態基底（透熱基底）$\{\ket{0}, \ket{1}\}$ において，時間依存ハミルトニアンを行列表現すると：
    \begin{equation}
        \mathcal{H}_{\text{LZ}}(t) = \frac{1}{2}\begin{pmatrix} -v t & \Delta_{\min} \\ \Delta_{\min} & v t \end{pmatrix} = -\frac{v t}{2} Z + \frac{\Delta_{\min}}{2} X
    \end{equation}
    と表される．ここで $v > 0$ は透熱エネルギーの掃引速度（傾き），$\Delta_{\min} > 0$ は2状態間の量子力学的結合（トンネル分裂 / オフダイアゴナル相互作用）である．
    仮に結合が存在しない（$\Delta_{\min} = 0$）場合，透熱エネルギー固有値は対角成分 $\epsilon_0(t) = -\frac{vt}{2}$，$\epsilon_1(t) = +\frac{vt}{2}$ となり，時刻 $t=0$ で完全に交差する直線を描く（エネルギー交差）．

    \newpage
    \item \textbf{断熱準位（Adiabatic Levels）の厳密解と反交差（Avoided Crossing）の幾何学}:\\
    非対角結合 $\Delta_{\min} > 0$ が存在する場合，各時刻 $t$ における瞬時固有値 $E$ は，特性方程式：
    \begin{equation}
        \det\left(\mathcal{H}_{\text{LZ}}(t) - E I\right) = E^2 - \left(\frac{v t}{2}\right)^2 - \left(\frac{\Delta_{\min}}{2}\right)^2 = 0
    \end{equation}
    を解くことで，一切の近似なく厳密に求まる：
    \begin{equation}
        E_\pm(t) = \pm \frac{1}{2} \sqrt{(v t)^2 + \Delta_{\min}^2}
    \end{equation}
    この瞬時固有エネルギー $E_-(t)$（基底状態）および $E_+(t)$（励起状態）を\textbf{断熱準位}と呼ぶ．
    \begin{itemize}
        \item 遠方時刻 $|t| \gg \Delta_{\min}/v$ では，$E_\pm(t) \approx \pm \frac{v|t|}{2}$ となり，透熱エネルギーの直線に漸近する．
        \item しかし交差中心 $t=0$ では，エネルギー差が $E_+(0) - E_-(0) = \Delta_{\min} > 0$ となり，\textbf{2つの準位は決して交差せず，互いに反発し合って離れる（反交差: Avoided Crossing）}！
    \end{itemize}
    このとき，断熱基底状態 $\ket{E_-(t)}$ のスピン状態は，時刻 $t \to -\infty$ の $\ket{0}$ から，反交差近傍 $t=0$ での均等重ね合わせ $\frac{\ket{0}-\ket{1}}{\sqrt{2}}$ を経て，時刻 $t \to +\infty$ では $\ket{1}$ へと連続的・滑らかに反転する．

    \item \textbf{ツェナーの遷移確率公式（Zener Formula, 1932）}:\\
    時刻 $t \to -\infty$ において基底状態 $\ket{E_-(-\infty)} = \ket{0}$ に準備された系が，時刻 $t \to +\infty$ まで発展したとき，目的の基底状態 $\ket{E_+(+\infty)} = \ket{1}$ に留まれず，\textbf{透熱準位のレールに沿って直進して励起状態 $\ket{0}$ へと脱落（非断熱遷移）してしまう確率}は，ツェナー（1932）によって厳密に解かれた：
    \begin{equation}
        P_{\text{LZ}} = \exp\left(-\frac{\pi \Delta_{\min}^2}{2\hbar v}\right)
    \end{equation}

    \item \textbf{なぜ $\frac{\Delta_{\min}^2}{\hbar v}$ なのか？――時定数の比による物理的直観}:\\
    指数の引数に現れる無次元パラメータ $\gamma_{\text{LZ}} = \frac{\Delta_{\min}^2}{\hbar v}$ は，以下の2つの特徴的時定数の比として極めて直観的に理解できる：
    \begin{enumerate}[label=(\alph*)]
        \item \textbf{反交差領域の通過時間 $\tau_{\text{cross}}$}:
        準位が非対角結合によって大きく曲がる時間幅は，$|vt| \lesssim \Delta_{\min}$ の範囲であるから，$\tau_{\text{cross}} \sim \frac{\Delta_{\min}}{v}$ である．
        \item \textbf{量子力学的応答時間（ボーア周期）$\tau_{\text{Bohr}}$}:
        最小エネルギーギャップ $\Delta_{\min}$ に対応する系の量子位相の応答周期は，不確定性関係およびボーア周波数 $\omega = \Delta_{\min}/\hbar$ より，$\tau_{\text{Bohr}} \sim \frac{\hbar}{\Delta_{\min}}$ である．
    \end{enumerate}
    両者の比を取ると：
    \begin{equation}
        \frac{\tau_{\text{cross}}}{\tau_{\text{Bohr}}} \sim \frac{\Delta_{\min} / v}{\hbar / \Delta_{\min}} = \frac{\Delta_{\min}^2}{\hbar v}
    \end{equation}
    が現れる！
    \begin{itemize}
        \item \textbf{断熱通過（$\tau_{\text{cross}} \gg \tau_{\text{Bohr}}$ すなわち $\frac{\Delta_{\min}^2}{\hbar v} \gg 1$）}:\\
        ボトルネックを通過する間に量子位相が何周期も高速回転するため，系はハミルトニアンの変化に完全に追従し，断熱基底状態に留まり続ける（$P_{\text{LZ}} = e^{-\text{大}} \approx 0$）．
        \item \textbf{急変脱落（$\tau_{\text{cross}} \ll \tau_{\text{Bohr}}$ すなわち $\frac{\Delta_{\min}^2}{\hbar v} \ll 1$）}:\\
        通過が速すぎて量子系が応答できず，波束は元の状態のまま直進して励起状態へ跳び移る（$P_{\text{LZ}} \approx 1$）．
    \end{itemize}

    \item \textbf{AQC / 量子アニーリングへの直接的要請}:\\
    全実行時間を $T$ とすると，掃引速度は $v \propto 1/T$ である．断熱遷移確率を抑制して最適解を得る（$P_{\text{LZ}} \ll 1$）ためには，$\frac{\pi \Delta_{\min}^2}{2\hbar v} \gg 1$ すなわち：
    \begin{equation}
        T \gg \frac{\hbar}{\Delta_{\min}^2}
    \end{equation}
    が必然的に要求される．これは前項で一般的な時間依存摂動論から導いた量子断熱条件と完全に調和する．
\end{enumerate}
\end{insightbox}

\begin{figure}[htbp]
\centering
\begin{tikzpicture}[scale=1.05, >=latex]
    % Axes
    \draw[->, thick] (0, 0) -- (9.0, 0) node[right, font=\small] {時間進行 $t / T$ ($0 \to 1$)};
    \draw[->, thick] (0, 0) -- (0, 5.2) node[above, font=\small] {エネルギー $E(t)$};

    % Energy levels with avoided crossing
    % Excited state curve
    \draw[very thick, red!80!black] (0.8, 4.4) .. controls (3.0, 3.8) and (4.0, 2.9) .. (4.5, 2.8) .. controls (5.0, 2.9) and (6.0, 3.8) .. (8.2, 4.5);
    \node[above right, font=\small\bfseries, text=red!80!black] at (8.2, 4.4) {第1励起状態 $E_1(t)$};

    % Ground state curve
    \draw[very thick, blue!80!black] (0.8, 0.9) .. controls (3.0, 1.2) and (4.0, 1.9) .. (4.5, 2.0) .. controls (5.0, 1.9) and (6.0, 1.1) .. (8.2, 0.8);
    \node[above right, font=\small\bfseries, text=blue!80!black] at (8.2, 0.8) {基底状態 $E_0(t)$ (最適解 $\ket{x^*}$)};

    % Initial state label with helper line
    \node[above, font=\footnotesize\bfseries, text=blue!75!black] at (1.2, 2.2) {初期基底状態 $\ket{s} = \ket{+}^{\otimes n}$};
    \draw[->, thin, blue!70!black] (1.2, 2.1) -- (1.2, 1.1);

    % Phase transition point dashed line
    \draw[dashed, gray!80] (4.5, 0) -- (4.5, 2.0);
    \node[below, font=\footnotesize] at (4.5, 0) {$t^*$ (相転移点)};

    % Minimum gap arrow and label
    \draw[<->, very thick, purple] (4.5, 2.05) -- (4.5, 2.75);
    \node[right, font=\footnotesize\bfseries, text=purple] at (4.65, 2.4) {最小ギャップ $\Delta_{\min}$};

    % Landau-Zener non-adiabatic transition arrow
    \draw[->, very thick, dashed, red!70!black] (3.7, 1.7) -- (5.3, 3.1)
        node[pos=0.75, above left, font=\scriptsize\bfseries, text=red!75!black, fill=white, inner sep=1pt] {非断熱遷移（局所解へ脱落）};

    % Bottom warning note
    \node[below, font=\footnotesize, text width=8cm, align=center, text=black] at (4.5, -0.6) {
        $T$ が不足すると狭いボトルネック $\Delta_{\min}$ で励起状態へ跳び移り探索が失敗する
    };
\end{tikzpicture}
\caption{断熱量子計算における準位の反交差（Avoided Level Crossing）とランダウ・ツェナー遷移．横磁場とコスト関数が拮抗する相転移点 $t^*$ においてエネルギーギャップ $\Delta_{\min}$ が最小となり，ここを乗り越えるために膨大な時間 $T \propto 1/\Delta_{\min}^2$ が要求される．}
\label{fig:avoided_crossing}
\end{figure}

異なるエネルギー準位は，摂動（量子ゆらぎ・横磁場）が存在すると交差することができず，互いに反発して「反交差」を起こす．
最小ギャップ $\Delta_{\min}$ の極小領域を通過する際，ハミルトニアンの変化速度が速すぎると，波束は基底状態のレールから外れ，励起状態へとトンネリング脱落してしまう（ランダウ・ツェナー遷移確率 $P_{\text{LZ}} \approx \exp\left(-\frac{\pi \Delta_{\min}^2}{2 v}\right)$）．
特に，NP困難な問題では\textbf{量子一次相転移（First-Order Quantum Phase Transition）}やスピングラス相におけるアンダーソン局在\cite{altshuler2010anderson}が発生し，量子ビット数 $N$ に対してギャップが指数関数的に極小化（$\Delta_{\min} \propto e^{-\alpha N}$）する．
その結果，断熱性を維持するために必要な時間 $T$ は $T \propto e^{2\alpha N}$ へと指数爆発し，\textbf{「原理的には解けるが，現実的な時間内では解けない」というAQCの根本的限界}が明らかとなった．

\subsection{4. 量子アニーリング（Quantum Annealing: QA）との深層的関係と決定的な相違点}
ここで，多くの読者（特に物理学や情報科学を学ぶ学生）が抱く極めて自然な疑問がある：
\begin{quote}
\textit{「横磁場から始めてイジング模型の基底状態を求めるアナログ発展といえば，門脇・西森（1998）が提唱した\textbf{『量子アニーリング（Quantum Annealing: QA）』}や，カナダのD-Wave Systems社が実用化した量子アニーラーと何が違うのか？AQCとQAは同一のものなのか？」}
\end{quote}
結論から言えば，\textbf{AQCと量子アニーリングは「時間依存ハミルトニアンによる基底状態探索」という数理的枠組みを共有しつつも，物理的動作環境（孤立系か開放系か），数学的目標（計算万能性か組合せ最適化ヒューリスティクスか），そして断熱定理に対する依存度において決定的な3大相違点を持つ}．本項では，一次文献\cite{kadowaki1998quantum,farhi2000quantum,aharonov2007adiabatic,albash2018adiabatic}に基づき，両者の共通点と相違点を厳密に解き明かす．

\subsubsection*{1. 歴史的起源と着想の対比}
\begin{itemize}
    \item \textbf{量子アニーリング (QA: Kadowaki \& Nishimori, 1998\cite{kadowaki1998quantum})}:
    東京工業大学の門脇正史と西森秀稔によって提唱された．その起源は，統計物理学におけるスピングラスの基底状態探索および古典の\textbf{シミュレーテッド・アニーリング（Simulated Annealing: SA\cite{kirkpatrick1983optimization}）}の量子拡張にある．
    SAでは，熱揺らぎ（ボルツマン因子 $\propto e^{-\Delta E / k_B T}$）を利用して古典エネルギー障壁を乗り越えるため，細く高いポテンシャル障壁に遭遇すると脱出確率が指数関数的に極小化する．これに対し門脇と西森は，横磁場 $\Gamma(t) \sum_i X_i$ を導入し，\textbf{「量子トンネル効果（量子揺らぎ）」によって障壁をすり抜けるメタヒューリスティクス}としてQAを考案した．WKB近似によれば，幅 $w$，高さ $V$ の障壁に対するトンネル確率は $P_{\text{tunnel}} \propto \exp\left(-\frac{2}{\hbar}\int \sqrt{2m(V - E)}dx\right)$ であり，幅 $w$ が細ければどれほど高い障壁であっても瞬時に貫通できるという物理的強みを持つ．
    
    \item \textbf{断熱量子計算 (AQC: Farhi et al., 2000\cite{farhi2000quantum}, 2001\cite{farhi2001quantum})}:
    MITの理論物理学者・計算機科学者である Edward Farhi らによって提唱された．彼らの動機は，Shorの素因数分解やGroverの探索に続く「NP困難な問題を解く普遍的な量子計算モデル」の構築にあった．すなわち，メタヒューリスティクスにとどまらず，量子断熱定理\cite{born1928beweis,kato1950adiabatic}を計算原理とする厳密なアルゴリズムとして考案された．
\end{itemize}

\subsubsection*{2. 数理形式の共通点}
両者は，全系の時間発展を以下の時間依存ハミルトニアンで記述する点で完全に一致している：
\begin{equation}
    H(t) = A(t) H_{\text{init}} + B(t) H_C = -A(t)\sum_{i=1}^n X_i + B(t)\left(\sum_{i} h_i Z_i + \sum_{i<j} J_{ij} Z_i Z_j\right)
\end{equation}
ここで時間スケジュール関数は，$t=0$ で $A(0) = 1, B(0) = 0$（自明な一様重ね合わせ状態 $\ket{+}^{\otimes n}$ が厳密な基底状態）であり，$t=T$ で $A(T) = 0, B(T) = 1$（目的のイジング模型 $H_C$）へと連続的に変化する．

\subsubsection*{3. 決定的な3大相違点}
形式的な数式が酷似しているにもかかわらず，物理学および計算機科学の立場からは以下の決定的な違いが存在する（表\ref{tab:aqc_vs_qa}）：

\begin{enumerate}
    \item \textbf{孤立系のユニタリ発展（AQC） vs 開放系の散逸ダイナミクス（QA）}:
    \begin{itemize}
        \item \textbf{AQC}: 絶対零度（$T=0$）の完全孤立系を前提とする．シュレーディンガー方程式に従う純粋なユニタリ発展であり，断熱条件 $T \gg \max |\mel{1}{\dot{H}}{0}| / \Delta_{\min}^2$ が破れて第1励起状態 $\ket{1(t)}$ へ跳び移った瞬間，その励起エネルギーを逃がす術はなく，計算は即座に失敗（誤った局所解）となる．
        \item \textbf{QA（現実の量子アニーラー）}: 希釈冷凍機（$T \sim 12\text{--}15\text{ mK}$）で動作するが，超伝導基板の格子振動（フォノン）や周囲の電磁場環境（熱浴）と不可避に結合した\textbf{開放量子系（Open Quantum System）}である．そのダイナミクスはユニタリ発展ではなく，リンドブラッド（Lindblad）型マスター方程式に従う．
        ここで決定的に重要な物理的差異は，\textbf{「環境への熱緩和（Thermal Relaxation / Incoherent Tunneling）」}である\cite{albash2018adiabatic}．アニーリング途中でランダウ・ツェナー遷移や熱揺らぎによって励起状態へ跳び上がったとしても，\textbf{系が環境熱浴へフォノンや光子を放出してエネルギーを散逸させることで，再び基底状態へと緩和（クーリング）して戻ってくる機構}が働く．このため，QAは厳密な断熱条件に縛られず，有限の高速掃引でも高い確率で最適解に到達する実用的探索能力を発揮する．
    \end{itemize}

    \item \textbf{計算万能性（Universality）の有無}:
    \begin{itemize}
        \item 通常の横磁場イジング模型（$H(t) = -A(t)\sum X_i + B(t)H_C$）は，計算基底における非対角成分がすべて実数非正（$\le 0$）である「ストカスティック（stoquastic）ハミルトニアン」に分類される．この場合，量子モンテカルロ（QMC）法で負符号問題が発生しないため，古典計算機で多項式時間サンプリングできる領域が広く存在する．
        \item しかし，Aharonov et al. (2007)\cite{aharonov2007adiabatic}の金字塔的定理により，\textbf{非ストカスティック（non-stoquastic）な相互作用（例：$X_i X_j$, $Y_i Y_j$ や非対角複素位相項）を含む2局所ハミルトニアンを用いたAQCは，標準的な回路型万能量子計算（BQP）と多項式時間の換算で厳密に等価（万能断熱量子計算: Universal AQC）}であることが証明された．すなわち，AQCは単なる最適化手法にとどまらず，ゲート型量子コンピュータと同一の万能計算能力を秘めた普遍的計算モデルである．
    \end{itemize}

    \item \textbf{評価基準と思想（計算量理論 vs メタヒューリスティクス）}:
    \begin{itemize}
        \item AQCは理論計算機科学の厳密性を重視し，「最悪ケースにおいてギャップ $\Delta_{\min}$ が量子ビット数 $N$ に対して多項式的にしか縮まないか，指数関数的に閉じるか」という漸近的計算量を追求する．
        \item QAは実用アルゴリズムの観点から，相転移点近傍で非断熱遷移が生じることを織り込み済みで，多スタート探索（多数回アニーリングショット）と熱緩和を組み合わせ，現実的な時間内で「十分良質な近似解・厳密解」を高速にサンプリングすることを目指す．
    \end{itemize}
\end{enumerate}

\begin{table}[htbp]
\centering
\caption{断熱量子計算（AQC）と量子アニーリング（QA）の物理・数学・計算論的対比}
\label{tab:aqc_vs_qa}
\begin{tabularx}{\textwidth}{l X X}
\toprule
\textbf{比較項目} & \textbf{断熱量子計算 (AQC)} & \textbf{量子アニーリング (QA)} \\
\midrule
\textbf{提唱者・年代} & E.~Farhi et al. (2000, 2001)\cite{farhi2000quantum,farhi2001quantum} & 門脇正史，西森秀稔 (1998)\cite{kadowaki1998quantum} \\
\textbf{物理系} & 完全孤立系（Closed Quantum System） & 熱浴と結合した開放量子系（Open System） \\
\textbf{動作温度} & 絶対零度 $T = 0$（純粋状態） & 極低温有限温度 $T > 0$（熱混合状態） \\
\textbf{基礎方程式} & シュレーディンガー方程式（ユニタリ発展） & Lindbladマスター方程式（散逸ダイナミクス） \\
\textbf{断熱定理の扱い} & 厳格に成立必須（$T \gg 1/\Delta_{\min}^2$） & 破れても熱緩和による基底状態復帰を許容 \\
\textbf{励起への耐性} & 励起＝探索失敗（エネルギー散逸なし） & 熱浴へのエネルギー散逸により再緩和可能 \\
\textbf{計算万能性} & 非ストカスティック項で万能量子計算\cite{aharonov2007adiabatic} & 横磁場イジングではストカスティック（特化型） \\
\textbf{ハードウェア} & 理想的量子プロセッサ & 超伝導フラックス量子ビット（D-Wave等） \\
\bottomrule
\end{tabularx}
\end{table}

\subsubsection*{4. AQC / QA から QAOA への進化の必然性}
それでは，なぜこれら連続時間のアナログ方式（AQC / QA）が存在するにもかかわらず，ゲート型量子コンピュータ上で動作する\textbf{QAOA}が必要とされたのだろうか？そこには2つのハードウェア的・アルゴリズム的宿命があった：

\begin{enumerate}
    \item \textbf{結合トポロジーの制約とマイナー埋め込みの壁}:
    D-Waveなどの量子アニーラーは物理的な素子配線（ChimeraグラフやPegasusグラフなど）に隣接結合が固定されている．したがって，全結合グラフや高次結合を持つ実問題を解くためには，複数の物理量子ビットを強い強磁性結合（FM結合）で結びつけて1つの論理ビットとして扱う\textbf{「マイナー埋め込み（Minor Embedding\cite{choi2008different}）」}が不可欠となる．これにより，問題規模に対して必要な量子ビット数が数倍〜数十倍に膨れ上がり，強いFM結合が切断される「チェーンブレイク（Chain Break）」エラーがボトルネックとなる．
    
    \item \textbf{NISQゲート型ハードウェアにおけるトロッター化の壁}:
    汎用ゲート型量子コンピュータ（IBM，Google，富士通など）でAQCの連続時間発展を再現しようとすると，後述する\textbf{トロッター分解}によって離散ゲート列に変換する必要がある（次節\ref{sec:trotter_detailed}）．しかし，断熱条件を満たすための長時間 $T$ を離散化すると，数千〜数万段に及ぶ深い量子回路（$p \gg 10^3$）が必要となり，現代のNISQプロセッサのコヒーレンス時間（数十マイクロ秒）とゲートノイズの壁に衝突して完全に計算が崩壊する．
\end{enumerate}

この\textbf{「アナログアニーラーの結合制約」と「トロッター化AQCのNISQにおける破綻」を同時に打ち破る歴史的ブレークスルー}として考案されたのが，次章（第4章）で詳解する \textbf{QAOA（Farhi et al., 2014\cite{farhi2014quantum}）} なのである！

\section{トロッター分解 (Trotter-Suzuki Decomposition)：アナログからディジタルへ}
\label{sec:trotter_detailed}

AQCは本来，専用のアナログ量子デバイス（現在のD-Wave量子アニーラーなど）で連続的な物理時間をかけて実行される．
しかし，IBM，Google，Rigetti，富士通などが開発する汎用的な\textbf{ゲート型量子コンピュータ（Gate-Based QPU）}は，離散的な量子ゲート（ユニタリ行列）の系列としてしかアルゴリズムを実行できない．
ここで登場するのが，リー・トロッターの積公式\cite{trotter1959product}および鈴木増雄による高次分解公式\cite{suzuki1976generalized}，そしてLloyd (1996)\cite{lloyd1996universal}による量子シミュレーションへの適用である．

\subsection{1. 非可換演算子の壁と BCH 公式}
ハミルトニアンの時間発展演算子 $U(T) = \exp\left(-i \int_0^T H(t) dt\right)$ をゲート型回路に落とし込む際の最大の障壁は，\textbf{初期ハミルトニアン $H_M$（横磁場 $X$）とコストハミルトニアン $H_C$（イジング $Z$）が互いに可換でない（$[H_M, H_C] \neq 0$）}という点である．

通常のスカラー関数であれば $e^{A + B} = e^A e^B$ であるが，非可換な演算子（行列）に対してはこれは一般に成り立たない．
厳密な関係は，\textbf{ベイカー・キャンベル・ハウスドルフ（Baker-Campbell-Hausdorff: BCH）の公式}によって与えられる：
\begin{equation}
    e^A e^B = \exp\left( A + B + \frac{1}{2}[A, B] + \frac{1}{12}[A, [A, B]] - \frac{1}{12}[B, [A, B]] + \dots \right)
\end{equation}
交換子 $[A, B] = AB - BA \neq 0$ の存在により，指数を単純に分割すると高次の歪み（交換子誤差）が不可避に発生する．

\subsection{2. リー・トロッター・鈴木の積公式 (Lie-Trotter Product Formula)}
この非可換性の壁を突破するため，全実行時間 $T$ を $p$ 個の微小時間刻み $\Delta t = T / p$（$p \gg 1$）に細かくスライス（時間離散化）する．
微小時間 $\Delta t \to 0$ の極限では，BCH展開の高次交換子項は $\Delta t^2$ の微小量として抑え込める：
\begin{equation}
    e^{-i (H_M + H_C) \Delta t} = e^{-i H_M \Delta t} e^{-i H_C \Delta t} + O(\Delta t^2)
\end{equation}
これを $p$ ステップ積み重ねることで，全時間発展を交互のユニタリ積として近似する：
\begin{equation}
    U(T) \approx \prod_{k=1}^p e^{-i H_M \Delta t_k} e^{-i H_C \Delta t_k}
\end{equation}
各微小ステップで発生する局所誤差は $O(\Delta t^2) = O(T^2 / p^2)$ であるため，全 $p$ ステップでの累積誤差は：
\begin{equation}
    \text{全トロッター誤差} = p \times O\left(\frac{T^2}{p^2}\right) = O\left(\frac{T^2}{p}\right) \xrightarrow{p \to \infty} 0
\end{equation}
すなわち，\textbf{ステップ数 $p$ を十分に大きくすれば，アナログな連続時間発展 $U(T)$ を，ゲート型コンピュータ上の離散量子ゲート（コスト層 $e^{-i H_C \gamma}$ とミキサー層 $e^{-i H_M \beta}$ の交互適用）として厳密に再現できる}ことが数学的に証明された（図\ref{fig:aqc_to_trotter}）．

\begin{figure}[htbp]
\centering
\begin{tikzpicture}[scale=0.95, >=latex]
    % Continuous AQC
    \draw[fill=blue!5, draw=blue!60!black, thick, rounded corners=2mm] (0, 0) rectangle (6, 3.5);
    \node[font=\bfseries, text=blue!70!black] at (3, 3.0) {アナログ断熱量子計算 (AQC)};
    \node[font=\small, align=center] at (3, 1.8) {
        連続時間発展\\
        $H(t) = (1 - \frac{t}{T}) H_M + \frac{t}{T} H_C$\\
        $U(T) = \mathcal{T} \exp\left(-i \int_0^T H(t) dt\right)$
    };
    \node[font=\footnotesize, text=red!70!black] at (3, 0.6) {長時間 $T$ と連続外場制御が必要};

    % Arrow with Trotter
    \draw[->, very thick, purple] (6.2, 1.75) -- (8.3, 1.75) node[midway, above=2pt, font=\footnotesize\bfseries] {トロッター分解};
    \node[below=2pt, font=\scriptsize, text=purple] at (7.25, 1.6) {$p$ 個のスライス};

    % Discrete QAOA-like circuit
    \draw[fill=green!5, draw=green!60!black, thick, rounded corners=2mm] (8.5, 0) rectangle (14.5, 3.5);
    \node[font=\bfseries, text=green!60!black] at (11.5, 3.0) {離散量子回路（ディジタル化）};
    \node[font=\small, align=center] at (11.5, 1.8) {
        交互ユニタリ積\\
        $\prod_{k=1}^p e^{-i \beta_k H_M} e^{-i \gamma_k H_C}$\\
        コスト層とミキサー層の繰り返し
    };
    \node[font=\footnotesize, text=blue!70!black] at (11.5, 0.6) {ゲート型量子回路で実行可能！};
\end{tikzpicture}
\caption{アナログなAQCの連続時間発展から，トロッター分解によるディジタル量子回路への変換．非可換な $H_M$ と $H_C$ を微小時間幅で交互に適用することで，ゲート型ハードウェアへの実装路が切り拓かれた．}
\label{fig:aqc_to_trotter}
\end{figure}

\section{トロッター化AQCの挫折と「QAOAへのパラダイムシフト」}

しかし，この「AQCをそのままディジタル化する」というアイデアは，NISQ（Noisy Intermediate-Scale Quantum）時代の実機ハードウェアにおいて\textbf{絶望的な壁}に直面した．

\begin{alertbox}[トロッター化AQCを阻んだNISQの壁]
\begin{enumerate}\setlength{\itemsep}{0pt}\setlength{\parskip}{0pt}
    \item \textbf{回路深度の爆発}: トロッター誤差を抑えるには $p \gg T^2$ が必要である．前述の通りNP困難問題では $T$ 自体が指数関数的（$T \propto e^{2\alpha N}$）に長いため，必要なレイヤー数 $p$ は数千〜数百万段に膨れ上がる．
    \item \textbf{実機デコヒーレンス}: 現在のNISQ実機では，量子ビットのコヒーレンス時間（寿命）は数十〜数百マイクロ秒であり，エラー率（$0.1\% \sim 1\%$）を考慮すると，実行可能な回路深度はせいぜい $p \le 5 \sim 10$ 程度である．数千段の回路など到底動かない．
\end{enumerate}
\end{alertbox}

\subsection{Farhiらの歴史的大逆転の発想（2014年: QAOAの誕生）}
この八方塞がりの状況下で，Farhi, Goldstone, Gutmann (2014) が打ち出したのが，以下の\textbf{「思考の完全なコペルニクス的転回」}であった：

\begin{pointbox}[Farhiらの核心的洞察: 物理シミュレーションから変分最適化へ]
\textit{「我々の真の目的は，AQCの物理時間発展を寸分違わず忠実に再現することではない．\\
目的はあくまで\textbf{『最適解のビット列 $x^*$ を見つけること』}のはずだ．\\
ならば，$(\beta_k, \gamma_k)$ を断熱理論が決める微小量 $\Delta t$ に固定するのをやめよう．\\
角度パラメータを\textbf{自由な変分パラメータとして全開放}し，\textbf{古典オプティマイザに最適な角度を逆算}させればよいではないか！」}
\end{pointbox}

この発想の転換により：
\begin{itemize}\setlength{\itemsep}{2pt}
    \item 断熱的にゆっくり這う必要がなくなった．
    \item \textbf{量子力学的干渉（強め合い・弱め合い）とトンネリングを最大限に活かし，たった $p=1$ や $p=2$ の極浅い回路深度であっても，大域的最適解へとダイレクトにショートカット（近道）する状態}を作り出せるようになった．
    \item 回路が浅いため，現代のNISQ量子コンピュータ上でもノイズに負けずに実行可能となった．
\end{itemize}

これが，物理学の「断熱量子計算」から，情報科学の「量子近似最適化アルゴリズム（QAOA）」への歴史的進化の本質である．
次章では，この革新的な原論文（Farhi et al., 2014）の数式と理論を，一行一行解読していく．

\chapter{QAOA原論文（Farhi et al., 2014）の完全徹底解読}

\begin{paperbox}[原著論文データ]
\textbf{タイトル}: A Quantum Approximate Optimization Algorithm\\
\textbf{著者}: Edward Farhi, Jeffrey Goldstone, Sam Gutmann (MIT)\\
\textbf{公表}: arXiv:1411.4028 (2014)
\end{paperbox}

\section{アルゴリズムの定式化と変分アンサンツ}
Farhiらの決定的なブレークスルー\cite{farhi2014quantum}は，第3章で論じた「断熱量子計算（AQC）のアナログ時間発展を忠実にトロッター化する」という発想を根底から覆し，\textbf{「各ステップの回転角 $(\vec{\gamma}, \vec{\beta})$ を自由な変分パラメータとして開放し，浅い回路で古典オプティマイザと協調して最適解への量子ショートカットを実現する」}という変分量子アルゴリズム（VQA\cite{cerezo2021variational}）へと昇華させた点にある．

\subsection{1. 初期状態と変分アンサンツの定式化}
QAOAは，全 $n$ 量子ビットの一様重ね合わせ状態を初期状態として探索を開始する：
\begin{equation}
    \ket{s} = \ket{+}^{\otimes n} = \frac{1}{\sqrt{2^n}} \sum_{z \in \{0, 1\}^n} \ket{z} = H^{\otimes n} \ket{0}^{\otimes n}
\end{equation}
これはアダマールゲート $H$ を全ビットに並列適用することで，回路深さ $\mathcal{O}(1)$ で瞬時に準備できる．
この状態は，横磁場ハミルトニアン $B = \sum_{j=1}^n X_j$ の最大固有値（$+n$）に属する固有状態（基底状態）である．

深さ $p \ge 1$ に対し，$2p$ 個の実パラメータ $\vec{\gamma} = (\gamma_1, \dots, \gamma_p)$ および $\vec{\beta} = (\beta_1, \dots, \beta_p)$ を導入する．
問題の目的関数をエンコードしたコストハミルトニアン $C$ による「コスト時間発展」$U(C, \gamma) = e^{-i \gamma C}$ と，横磁場ハミルトニアン $B$ による「ミキサー時間発展」$U(B, \beta) = e^{-i \beta B}$ を交互に $p$ 回作用させることで，変分アンサンツ状態 $\ket{\vec{\gamma}, \vec{\beta}}$ を構成する：
\begin{equation}
    \ket{\vec{\gamma}, \vec{\beta}} = U(B, \beta_p) U(C, \gamma_p) \cdots U(B, \beta_1) U(C, \gamma_1) \ket{s}
\end{equation}
ここでミキサー演算子は各ビットへの1量子ビット回転ゲート $R_X(2\beta) = e^{-i \beta X}$ のテンソル積である：
\begin{equation}
    U(B, \beta) = \prod_{j=1}^n e^{-i \beta X_j} = \bigotimes_{j=1}^n (\cos\beta I - i \sin\beta X)
\end{equation}

\subsection{2. コスト時間発展とミキサー時間発展の物理描像と量子干渉（図\ref{fig:cost_mixer_evolution}）}
QAOAの心臓部を担う2つのユニタリ演算子 $U(C, \gamma)$ と $U(B, \beta)$ は，全く異なる物理的役割を担い，両者が交互に連動することで初めて大域最適解が抽出される（図\ref{fig:cost_mixer_evolution}）．

\begin{figure}[htbp]
\centering
\begin{tikzpicture}[scale=0.88, >=latex]
    % Panel (a): Cost Hamiltonian Evolution (Phase Imprinting)
    \draw[fill=purple!3, draw=purple!45, thick, rounded corners=3mm] (0.0, -0.6) rectangle (4.6, 5.2);
    \node[fill=purple!15, draw=purple!60, font=\scriptsize\bfseries, text=purple!90!black, rounded corners=1mm, inner sep=2.5pt] at (2.3, 4.85) {(a) コスト時間発展 $e^{-i\gamma C}$};

    % Energy axis on the left
    \draw[->, thick, gray!70] (0.55, 1.4) -- (0.55, 4.3);
    \node[font=\tiny\bfseries, text=gray!80!black, rotate=90] at (0.32, 2.85) {エネルギー $C(z)$};
    
    % Level 1: High cost (z_bad)
    \draw[thick, red!70!black] (0.7, 3.8) -- (1.8, 3.8);
    \node[above, font=\tiny\bfseries, text=red!80!black] at (1.25, 3.8) {高 $C(z_1)$};
    \draw[thick, gray!50] (2.5, 3.8) circle (0.35);
    \draw[->, very thick, red!80!black] (2.5, 3.8) -- (2.73, 4.03);
    \draw[->, thick, red!60!black] (2.76, 3.8) arc (0:120:0.26);
    \node[right, font=\tiny\bfseries, text=red!80!black] at (2.95, 3.8) {高速位相回転};

    % Level 2: Medium cost
    \draw[thick, orange!80!black] (0.7, 2.7) -- (1.8, 2.7);
    \node[above, font=\tiny\bfseries, text=orange!85!black] at (1.25, 2.7) {中 $C(z_2)$};
    \draw[thick, gray!50] (2.5, 2.7) circle (0.35);
    \draw[->, very thick, orange!85!black] (2.5, 2.7) -- (2.78, 2.84);
    \draw[->, thick, orange!60!black] (2.76, 2.7) arc (0:55:0.26);
    \node[right, font=\tiny\bfseries, text=orange!85!black] at (2.95, 2.7) {中速回転};

    % Level 3: Low cost / Optimal (z_opt)
    \draw[thick, teal!80!black] (0.7, 1.6) -- (1.8, 1.6);
    \node[above, font=\tiny\bfseries, text=teal!80!black] at (1.25, 1.6) {低 $C(z^*)$};
    \draw[thick, gray!50] (2.5, 1.6) circle (0.35);
    \draw[->, very thick, teal!80!black] (2.5, 1.6) -- (2.64, 1.3);
    \draw[->, thick, teal!60!black] (2.76, 1.6) arc (0:-65:0.26);
    \node[right, font=\tiny\bfseries, text=teal!80!black] at (2.95, 1.6) {基準・逆位相};

    % Text explanation (a)
    \draw[fill=white, draw=purple!25, rounded corners=1.5mm] (0.2, -0.45) rectangle (4.4, 0.95);
    \node[font=\tiny, align=left, text width=3.9cm] at (2.3, 0.25) {
        \textbf{物理描像：位相刻印（対角発展）}\\[1mm]
        ・固有状態：計算基底 $\ket{z}$\\
        ・エネルギー $C(z)$ に応じた周波数で位相回転\\
        ・\textbf{測定確率 $|c_z|^2$ は不変}（相対位相のみ蓄積）
    };

    % Panel (b): Mixer Hamiltonian Evolution (Quantum Tunneling)
    \draw[fill=blue!3, draw=blue!45, thick, rounded corners=3mm] (5.0, -0.6) rectangle (9.6, 5.2);
    \node[fill=blue!15, draw=blue!60, font=\scriptsize\bfseries, text=blue!90!black, rounded corners=1mm, inner sep=2.5pt] at (7.3, 4.85) {(b) ミキサー時間発展 $e^{-i\beta B}$};

    % Energy landscape with potential barrier
    \draw[thick, gray!70] (5.4, 3.2) .. controls (5.8, 1.5) and (6.2, 1.5) .. (6.6, 3.0)
                                    .. controls (7.0, 4.2) and (7.6, 4.2) .. (8.0, 3.0)
                                    .. controls (8.4, 1.5) and (8.8, 1.5) .. (9.2, 3.2);
    \node[font=\tiny\bfseries, text=gray!80!black] at (6.0, 1.3) {局所解 $\ket{z_{\text{local}}}$};
    \node[font=\tiny\bfseries, text=teal!90!black] at (8.6, 1.3) {大域解 $\ket{z^*}$};
    \node[font=\tiny\bfseries, text=red!80!black] at (7.3, 4.2) {エネルギー障壁};

    % Tunneling arrow through barrier
    \draw[<->, very thick, blue!80!black, dashed] (6.2, 2.4) -- (8.4, 2.4);
    \node[above, font=\tiny\bfseries, text=blue!90!black] at (7.3, 2.48) {量子トンネリング};
    \node[below, font=\fontsize{5}{6}\selectfont, text=blue!80!black] at (7.3, 2.32) {$X_j$ によるスピン反転};

    % Text explanation (b)
    \draw[fill=white, draw=blue!25, rounded corners=1.5mm] (5.2, -0.45) rectangle (9.4, 0.95);
    \node[font=\tiny, align=left, text width=3.9cm] at (7.3, 0.25) {
        \textbf{物理描像：振幅混合（非対角発展）}\\[1mm]
        ・非対角 $X$ 演算子がスピン反転を駆動\\
        ・古典的には越えられない障壁をすり抜ける\\
        ・\textbf{解空間全体に振幅を拡散・混合}
    };

    % Panel (c): Synergy / Trotterized Search
    \draw[fill=teal!3, draw=teal!45, thick, rounded corners=3mm] (10.0, -0.6) rectangle (14.6, 5.2);
    \node[fill=teal!15, draw=teal!60, font=\scriptsize\bfseries, text=teal!90!black, rounded corners=1mm, inner sep=2.5pt] at (12.3, 4.85) {(c) 交互時間発展のシナジー};

    % Stepped diagram of alternating evolution
    \draw[fill=gray!10, draw=gray!45, rounded corners=1.2mm] (10.25, 4.05) rectangle (14.35, 4.52);
    \node[font=\tiny\bfseries] at (12.3, 4.28) {初期状態：一様重ね合わせ $\ket{s} = \ket{+}^{\otimes n}$};

    \draw[->, thick, purple!75!black] (12.3, 4.05) -- (12.3, 3.65);

    \draw[fill=purple!10, draw=purple!50, rounded corners=1.2mm] (10.25, 3.12) rectangle (14.35, 3.65);
    \node[font=\tiny\bfseries, text=purple!90!black] at (12.3, 3.48) {コスト発展 $U(C, \gamma_1)$（位相差刻印）};
    \node[font=\fontsize{5}{6}\selectfont, text=purple!80!black] at (12.3, 3.25) {各ビット列に目的関数に応じた位相差を付与};

    \draw[->, thick, blue!75!black] (12.3, 3.12) -- (12.3, 2.72);

    \draw[fill=blue!10, draw=blue!50, rounded corners=1.2mm] (10.25, 2.18) rectangle (14.35, 2.72);
    \node[font=\tiny\bfseries, text=blue!90!black] at (12.3, 2.54) {ミキサー発展 $U(B, \beta_1)$（振幅混合）};
    \node[font=\fontsize{5}{6}\selectfont, text=blue!80!black] at (12.3, 2.31) {トンネリングにより位相差が波の干渉を起こす};

    \draw[->, thick, dashed, teal!80!black] (12.3, 2.18) -- (12.3, 1.70)
        node[midway, fill=teal!3, inner sep=1pt, font=\fontsize{5}{6}\selectfont\bfseries, text=teal!85!black] {$p$ 回交互反復（トロッター化）};

    \draw[fill=teal!18, draw=teal!65, thick, rounded corners=1.2mm] (10.25, 1.15) rectangle (14.35, 1.70);
    \node[font=\tiny\bfseries, text=teal!90!black] at (12.3, 1.51) {最適解 $\ket{z^*}$ に確率濃縮！};
    \node[font=\fontsize{5}{6}\selectfont, text=teal!80!black] at (12.3, 1.28) {建設的干渉により大域解の振幅のみが増幅};

    % Text explanation (c)
    \draw[fill=white, draw=teal!25, rounded corners=1.5mm] (10.2, -0.45) rectangle (14.4, 0.95);
    \node[font=\tiny, align=left, text width=3.9cm] at (12.3, 0.25) {
        \textbf{なぜ交互適用か？：$[C, B] \neq 0$}\\[1mm]
        ・可換なら干渉は起きない（$[C,B]=0$）\\
        ・位相刻印と振幅混合の相補的連動により，\\
        ・\textbf{断熱量子計算（AQC）の離散トロッター化}を実現
    };
\end{tikzpicture}
\caption{コスト時間発展とミキサー時間発展の物理描像と交互適用のダイナミクス．(a) コスト時間発展 $e^{-i\gamma C}$ は対角演算子であり，各計算基底 $\ket{z}$ の複素位相をエネルギー $C(z)$ に比例した周波数で回転させて「位相差」を刻印する（測定確率は不変）．(b) ミキサー時間発展 $e^{-i\beta B}$ は非対角なスピン反転を駆動し，古典的には越えられないエネルギー障壁を「量子トンネリング」で突き抜けて振幅を混合・拡散させる．(c) 非可換性 $[C, B] \neq 0$ に基づく両者の交互適用により，刻印された位相差が量子干渉を起こし，ステップを経るごとに最適解 $\ket{z^*}$ へと測定確率が集約される．}
\label{fig:cost_mixer_evolution}
\end{figure}

\subsubsection*{1. コスト時間発展 $U(C, \gamma) = e^{-i\gamma C}$【位相刻印エンジン（対角発展）】}
\begin{itemize}
    \item \textbf{数学的本質（対角行列）}: 計算基底 $\ket{z}$ は古典コストハミルトニアン $C$ の同時固有状態（$C\ket{z} = C(z)\ket{z}$）であるため，任意の重ね合わせ状態 $\sum_z c_z \ket{z}$ に作用させると各基底は：
    \begin{equation}
        e^{-i \gamma C} \ket{z} = e^{-i \gamma C(z)} \ket{z}
    \end{equation}
    へと変換される（図\ref{fig:cost_mixer_evolution}(a)）．
    \item \textbf{測定確率は不変}: ここで極めて重要なのは，ビット列 $z$ の測定確率 $P(z) = |c_z e^{-i \gamma C(z)}|^2 = |c_z|^2$ は\textbf{1ミリも変化していない}という点である．状態間の確率振幅の移動は一切生じない．
    \item \textbf{直感的イメージ（時計の針とランナーのアナロジー）}:
    変わったのは，複素平面上における\textbf{「位相の向き（角度 $\theta_z = -\gamma C(z)$）」だけ}である．
    スタートラインに並んだ全 $2^n$ 通りのビット列のランナーを想像しよう．各ランナーは胸に時計を持っている．
    アインシュタイン・ド・ブロイの基本関係（$E = \hbar \omega$）に従い，各ランナーは\textbf{自身のコスト値（エネルギー） $C(z)$ に比例した角速度で時計の針をぐるぐると回転}させる．
    コストの高い（悪い）ビット列は針を高速回転させ，コストの低い（優れた最適解）ビット列は針をゆっくり回転させる．
    これにより，ビット列間に\textbf{「相対位相差 $\Delta \phi = -\gamma(C(z) - C(z'))$」が厳密に記録・記憶}される．
\end{itemize}

\subsubsection*{2. ミキサー時間発展 $U(B, \beta) = e^{-i\beta B}$【量子トンネリング・波の拡散（非対角発展）】}
\begin{itemize}
    \item \textbf{数学的本質（非対角行列・スピン反転）}: 横磁場 $B = \sum_j X_j$ のパウリ $X$ 演算子は，ビットを反転させる非対角演算子（$X\ket{0}=\ket{1}, X\ket{1}=\ket{0}$）である．
    単一ビットに対して $e^{-i\beta X} = \cos\beta I - i\sin\beta X$ と作用し，ハミング距離 $\Delta d = 1$ の隣接状態間で確率振幅を遷移させる（図\ref{fig:cost_mixer_evolution}(b)）．
    \item \textbf{直感的イメージ（量子トンネル効果による山越え）}:
    古典探索（シミュレーテッドアニーリングなど）では，高いエネルギー障壁に囲まれた局所解から脱出するには，熱揺らぎで山を「よじ登って乗り越える」しかなく，山が高いと脱出が不可能になる．
    これに対し，量子ミキサー $X$ は，障壁の高さに関係なく，\textbf{「量子力学的トンネル効果」によって山を水平に突き抜け，隣接する谷へと確率の波を染み出させて解空間全体に拡散・混合}させる．
\end{itemize}

\subsubsection*{3. なぜ交互適用なのか？：非可換性 $[C, B] \neq 0$ と二重スリット型量子干渉}
\begin{itemize}
    \item \textbf{非可換性の必然性}: もしコスト $C$ とミキサー $B$ が可換（$[C, B]=0$）であれば，$e^{-i\beta B} e^{-i\gamma C} = e^{-i(\gamma C + \beta B)}$ となり，干渉は起きない．パウリ $Z$ と $X$ の反可換性（$\{X, Z\}=0$）に基づく非可換性：
    \begin{equation}
        [C, B] = \sum_{\langle j, k \rangle} J_{jk} [Z_j Z_k, X_j + X_k] = 2i \sum_{\langle j, k \rangle} J_{jk} (Y_j Z_k + Z_j Y_k) \neq 0
    \end{equation}
    があるからこそ，位相差が振幅の干渉へと結合される．
    \item \textbf{直感的イメージ（ヤングの二重スリット実験との完全一致）}:
    コスト層とミキサー層の交互適用は，物理学における\textbf{「二重スリット実験の干渉縞形成」}と完全に同値である：
    \begin{enumerate}
        \item コスト層 $e^{-i\gamma C}$ は，スリットを通過した光に\textbf{「光路差（位相差）」}を刻印する．
        \item ミキサー層 $e^{-i\beta B}$ は，スリットを通過した光の波をスクリーンに向けて広げ，\textbf{「重ね合わせて混合」}させる．
        \item その結果，目的関数が悪いビット列同士は位相がバラバラであるため，波の山と谷が打ち消し合って消滅する（\textbf{相殺的干渉 Destructive Interference}）．
        \item 一方，目的関数が優れている最適解 $\ket{z^*}$ では，あらゆる経路から届く波の位相がピタリと揃い，波の山と山が強め合って測定確率が劇的に増大する（\textbf{建設的干渉 Constructive Interference}）！
    \end{enumerate}
    深さ $p$ を重ねるごとに，この「位相刻印 $\to$ 振幅干渉」のサイクルが反復され，解空間に広がっていた波束が目的関数最小の最適解 $\ket{z^*}$ へと自己組織化的に濃縮されていく．
\end{itemize}

\begin{insightbox}[横磁場イジング模型・シュレーディンガー方程式・QAOAの進化と繋がり]
\centering
\footnotesize
\begin{tabular}{p{1.2cm}p{3.1cm}p{3.3cm}p{5.4cm}}
\toprule
\textbf{階層} & \textbf{概念・枠組み} & \textbf{数理表現} & \textbf{本質的役割・物理的描像} \\
\midrule
\textbf{物理系} & \textbf{横磁場イジング模型 (TFIM)} & $H(t) = - \Gamma(t) \sum X_j + C$ & 古典イジング（対角）に横磁場（非対角）で量子揺らぎを導入 \\
\textbf{基礎法則} & \textbf{時間依存シュレーディンガー方程式} & $i \frac{\partial}{\partial t}\ket{\psi(t)} = H(t)\ket{\psi(t)}$ & トンネル効果による障壁透過と断熱発展（QA/AQCの原理） \\
\textbf{ディジタル化} & \textbf{リー・トロッター分解} & $U(T) \approx \prod_{k=1}^p e^{-i \beta_k B} e^{-i \gamma_k C}$ & 連続アナログ発展をゲート型回路の交互ユニタリ層へ写像 \\
\textbf{アルゴリズム} & \textbf{QAOA（変分最適化）} & $\max_{\vec{\gamma}, \vec{\beta}} \bra{\vec{\gamma}, \vec{\beta}} C \ket{\vec{\gamma}, \vec{\beta}}$ & 時間刻みを変分パラメータ化し，浅い回路で解をショートカット \\
\bottomrule
\end{tabular}
\end{insightbox}

\subsection{3. アンサンツ量子回路の構成と内部ゲート配線}
QAOAの全体回路アーキテクチャを図\ref{fig:qaoa_circuit_diagram}に示す．

\begin{figure}[htbp]
\centering
\begin{tikzpicture}[scale=0.88, >=latex]
    % Background layers
    \draw[dashed, red!60!black, thick, fill=red!3, rounded corners=1.5mm] (-0.2, -0.2) rectangle (1.3, 3.8);
    \node[above, font=\small\bfseries, text=red!60!black] at (0.55, 3.85) {初期化 $\ket{s}$};

    \draw[dashed, purple!70!black, thick, fill=purple!3, rounded corners=1.5mm] (1.7, -0.2) rectangle (6.2, 3.8);
    \node[above, font=\small\bfseries, text=purple!70!black] at (3.95, 3.85) {第 1 レイヤー ($p=1$)};

    \draw[dashed, purple!70!black, thick, fill=purple!3, rounded corners=1.5mm] (7.7, -0.2) rectangle (12.2, 3.8);
    \node[above, font=\small\bfseries, text=purple!70!black] at (9.95, 3.85) {第 $p$ レイヤー};

    % Wires (Full length)
    \foreach \y/\lab in {3.2/{q_0}, 2.2/{q_1}, 0.6/{q_{n-1}}} {
        \node[left] at (-0.3, \y) {$\ket{0}_{\lab}$};
        \draw[thick] (-0.3, \y) -- (14.2, \y);
    }
    \node at (-0.7, 1.4) {$\vdots$};
    \node at (6.95, 1.4) {$\cdots$};
    \node at (6.95, 3.2) {$\cdots$};
    \node at (6.95, 0.6) {$\cdots$};
    \node at (12.6, 1.4) {$\vdots$};

    % Initial Hadamard gates
    \foreach \y in {3.2, 2.2, 0.6} {
        \draw[fill=green!18, draw=green!65!black, thick, rounded corners=1mm] (0.3, \y-0.35) rectangle (1.0, \y+0.35) node[midway, font=\bfseries] {$H$};
    }

    % Layer 1: Cost unitary
    \fill[white] (2.0, 0.2) rectangle (4.2, 3.6);
    \draw[fill=purple!15, draw=purple!75!black, thick, rounded corners=2mm] (2.0, 0.2) rectangle (4.2, 3.6);
    \node[font=\bfseries, align=center, text=purple!85!black] at (3.1, 1.9) {コスト層\\$U(C, \gamma_1)$\\[2pt]\scriptsize $\prod e^{-i\gamma_1 C_{jk}}$};
    
    % Layer 1: Mixer unitaries
    \foreach \y in {3.2, 2.2, 0.6} {
        \fill[white] (4.7, \y-0.35) rectangle (5.9, \y+0.35);
        \draw[fill=blue!15, draw=blue!75!black, thick, rounded corners=1mm] (4.7, \y-0.35) rectangle (5.9, \y+0.35) node[midway, font=\scriptsize\bfseries] {$R_X(2\beta_1)$};
    }

    % Layer p: Cost unitary
    \fill[white] (8.0, 0.2) rectangle (10.2, 3.6);
    \draw[fill=purple!15, draw=purple!75!black, thick, rounded corners=2mm] (8.0, 0.2) rectangle (10.2, 3.6);
    \node[font=\bfseries, align=center, text=purple!85!black] at (9.1, 1.9) {コスト層\\$U(C, \gamma_p)$\\[2pt]\scriptsize $\prod e^{-i\gamma_p C_{jk}}$};

    % Layer p: Mixer unitaries
    \foreach \y in {3.2, 2.2, 0.6} {
        \fill[white] (10.7, \y-0.35) rectangle (11.9, \y+0.35);
        \draw[fill=blue!15, draw=blue!75!black, thick, rounded corners=1mm] (10.7, \y-0.35) rectangle (11.9, \y+0.35) node[midway, font=\scriptsize\bfseries] {$R_X(2\beta_p)$};
    }

    % Measurement
    \foreach \y in {3.2, 2.2, 0.6} {
        \fill[white] (13.1, \y-0.35) rectangle (14.1, \y+0.35);
        \draw[fill=gray!20, draw=black, thick, rounded corners=1mm] (13.1, \y-0.35) rectangle (14.1, \y+0.35);
        \draw[thick] (13.3, \y-0.15) arc (180:0:3mm);
        \draw[thick, ->] (13.6, \y-0.15) -- (13.9, \y+0.2);
    }
    \node[above, font=\small\bfseries] at (13.6, 3.85) {測定};
\end{tikzpicture}
\caption{QAOAの全体量子回路アーキテクチャ．全ビットへのアダマールゲート $H$ により一様重ね合わせ状態 $\ket{s} = \ket{+}^{\otimes n}$ を準備した後，問題の目的関数をエンコードしたコスト層 $U(C, \gamma_k)$ と横磁場ミキサー層 $U(B, \beta_k)$ を深さ $p$ 回交互に作用させ，最後に計算基底でサンプリング測定を行う．}
\label{fig:qaoa_circuit_diagram}
\end{figure}

\subsubsection{コスト層 $U(C, \gamma)$ の内部ゲート分解とイジング配線}
図\ref{fig:qaoa_circuit_diagram}における紫色のブロック「コスト層 $U(C, \gamma)$」は，一見ブラックボックスに見えるが，その内部は標準的な基本量子ゲートへと厳密に分解される．
問題ハミルトニアン $C$ が局所節の和 $C = \sum_{\langle j, k \rangle \in E} C_{\langle j, k \rangle}$，あるいは一般的なイジング相互作用 $C = \sum_{\langle j, k \rangle} J_{jk} Z_j Z_k + \sum_j h_j Z_j$ で表される場合，すべての演算子は計算基底の対角行列であり，互いに可換（$[Z_j Z_k, Z_l Z_m] = 0$）である．したがって，時間発展演算子はトロッター誤差なく厳密に因数分解できる：
\begin{equation}
    e^{-i \gamma C} = \left( \prod_{\langle j, k \rangle \in E} e^{-i \gamma J_{jk} Z_j Z_k} \right) \left( \prod_{j=1}^n e^{-i \gamma h_j Z_j} \right)
\end{equation}
前述の図\ref{fig:rzz_decomposition}で示した通り，2量子ビットイジング相互作用 $e^{-i \gamma J_{jk} Z_j Z_k}$ は CNOT--$R_Z(2\gamma J_{jk})$--CNOT ゲートで実装され，1体局所磁場項 $e^{-i \gamma h_j Z_j}$ は単一量子ビット回転 $R_Z(2\gamma h_j)$ で実装される．
コスト層の内部ゲート配線構造を図\ref{fig:cost_layer_internal}に示す．

\begin{figure}[htbp]
\centering
\begin{tikzpicture}[scale=0.88, >=latex]
    % Wires (y = 3.6, 2.4, 1.2, 0.0)
    \foreach \y/\j in {3.6/0, 2.4/1, 1.2/2, 0.0/3} {
        \node[left, font=\normalsize] at (-0.4, \y) {$\ket{q_\j}$};
        \draw[thick] (-0.4, \y) -- (13.6, \y);
    }

    % Interaction layer outer box
    \draw[dashed, purple!70!black, thick, fill=purple!3, rounded corners=2.5mm] (0.4, -0.65) rectangle (9.6, 4.55);
    \node[above, font=\small\bfseries, text=purple!80!black] at (5.0, 4.6) {2体イジング相互作用：$\prod_{\langle j,k \rangle} e^{-i \gamma J_{jk} Z_j Z_k}$};

    % Edge (0, 1) badge & interaction
    \node[fill=purple!12, draw=purple!40, text=purple!90!black, font=\scriptsize\bfseries, inner sep=2.5pt, rounded corners=1mm] at (3.0, 4.08) {辺 $(0, 1)$：$e^{-i \gamma J_{01} Z_0 Z_1}$};
    % CNOT 1 (ctrl q0, target q1)
    \filldraw[black] (1.6, 3.6) circle (2.8pt);
    \draw[thick] (1.6, 3.6) -- (1.6, 2.4);
    \draw[thick, fill=white] (1.6, 2.4) circle (4.8pt);
    \draw[thick] (1.6, 2.23) -- (1.6, 2.57); \draw[thick] (1.43, 2.4) -- (1.77, 2.4);
    % Rz(2\gamma J_01) on q1
    \draw[fill=orange!15, draw=orange!80!black, thick, rounded corners=1.2mm] (2.2, 2.05) rectangle (3.8, 2.75) node[midway, font=\scriptsize\bfseries] {$R_Z(2\gamma J_{01})$};
    % CNOT 2 (ctrl q0, target q1)
    \filldraw[black] (4.4, 3.6) circle (2.8pt);
    \draw[thick] (4.4, 3.6) -- (4.4, 2.4);
    \draw[thick, fill=white] (4.4, 2.4) circle (4.8pt);
    \draw[thick] (4.4, 2.23) -- (4.4, 2.57); \draw[thick] (4.23, 2.4) -- (4.57, 2.4);

    % Edge (1, 2) badge & interaction
    \node[fill=purple!12, draw=purple!40, text=purple!90!black, font=\scriptsize\bfseries, inner sep=2.5pt, rounded corners=1mm] at (7.4, 3.0) {辺 $(1, 2)$：$e^{-i \gamma J_{12} Z_1 Z_2}$};
    % CNOT 1 (ctrl q1, target q2)
    \filldraw[black] (6.0, 2.4) circle (2.8pt);
    \draw[thick] (6.0, 2.4) -- (6.0, 1.2);
    \draw[thick, fill=white] (6.0, 1.2) circle (4.8pt);
    \draw[thick] (6.0, 1.03) -- (6.0, 1.37); \draw[thick] (5.83, 1.2) -- (6.17, 1.2);
    % Rz(2\gamma J_12) on q2
    \draw[fill=orange!15, draw=orange!80!black, thick, rounded corners=1.2mm] (6.6, 0.85) rectangle (8.2, 1.55) node[midway, font=\scriptsize\bfseries] {$R_Z(2\gamma J_{12})$};
    % CNOT 2 (ctrl q1, target q2)
    \filldraw[black] (8.8, 2.4) circle (2.8pt);
    \draw[thick] (8.8, 2.4) -- (8.8, 1.2);
    \draw[thick, fill=white] (8.8, 1.2) circle (4.8pt);
    \draw[thick] (8.8, 1.03) -- (8.8, 1.37); \draw[thick] (8.63, 1.2) -- (8.97, 1.2);

    % Local fields box
    \draw[dashed, blue!70!black, thick, fill=blue!3, rounded corners=2.5mm] (10.1, -0.65) rectangle (13.3, 4.55);
    \node[above, font=\small\bfseries, text=blue!80!black] at (11.7, 4.6) {1体局所磁場：$\prod_j e^{-i \gamma h_j Z_j}$};
    \foreach \y/\j in {3.6/0, 2.4/1, 1.2/2, 0.0/3} {
        \draw[fill=yellow!20, draw=orange!85!black, thick, rounded corners=1.2mm] (10.5, \y-0.35) rectangle (12.9, \y+0.35) node[midway, font=\scriptsize\bfseries] {$R_Z(2\gamma h_{\j})$};
    }
\end{tikzpicture}
\caption{コスト層 $U(C, \gamma) = e^{-i \gamma C}$ の基本ゲート分解と内部配線構造．問題ハミルトニアン $C = \sum_{\langle j, k \rangle} J_{jk} Z_j Z_k + \sum_j h_j Z_j$ の各2体相互作用項は CNOT--$R_Z$--CNOT ブロックとして各量子ビット対に配線され，1体局所磁場項は単一量子ビット $R_Z$ ゲートとして並列適用される．すべての項が対角行列（可換）であるため，ゲートの適用順序は結果に影響を与えない．}
\label{fig:cost_layer_internal}
\end{figure}

\subsubsection{具象例：4頂点グラフ上のMaxCutに対する完全展開QAOA回路 ($p=1$)}
アルゴリズムの全容を完全に具体化するため，4個の頂点を持つパスグラフ $G = (V, E)$（$V=\{0, 1, 2, 3\}$，辺集合 $E=\{(0,1), (1,2), (2,3)\}$）に対する深さ $p=1$ のQAOA回路の全展開を図\ref{fig:qaoa_maxcut4_circuit}に示す．

MaxCut問題のコスト関数は，カットされた辺の総数として以下で定義される：
\begin{equation}
    C = \sum_{\langle j, k \rangle \in E} \frac{I - Z_j Z_k}{2}
\end{equation}
定数項 $I/2$ は全体位相 $e^{-i \gamma / 2}$ を与えるのみで物理的測定確率に寄与しないため，正味の回転は各辺に対する $R_{ZZ}(-\gamma) = \exp\left(i \frac{\gamma}{2} Z_j Z_k\right)$ で与えられる．

\begin{figure}[htbp]
\centering
\begin{tikzpicture}[scale=0.74, >=latex]
    % Left side: Target Graph G
    \draw[fill=gray!4, draw=gray!40, thick, rounded corners=2.5mm] (-5.2, -0.9) rectangle (-0.8, 4.65);
    \node[font=\small\bfseries, text=gray!80!black] at (-3.0, 4.25) {対象グラフ $G=(V, E)$};
    
    % Vertices and edges
    \node[circle, fill=blue!20, draw=blue!80!black, thick, inner sep=2.5pt, font=\small\bfseries] (v0) at (-4.1, 3.0) {$v_0$};
    \node[circle, fill=blue!20, draw=blue!80!black, thick, inner sep=2.5pt, font=\small\bfseries] (v1) at (-1.9, 3.0) {$v_1$};
    \node[circle, fill=blue!20, draw=blue!80!black, thick, inner sep=2.5pt, font=\small\bfseries] (v2) at (-1.9, 1.4) {$v_2$};
    \node[circle, fill=blue!20, draw=blue!80!black, thick, inner sep=2.5pt, font=\small\bfseries] (v3) at (-4.1, 1.4) {$v_3$};
    
    \draw[very thick, red!70!black] (v0) -- node[above, font=\scriptsize\bfseries, text=red!80!black] {$e_{01}$} (v1);
    \draw[very thick, red!70!black] (v1) -- node[right, font=\scriptsize\bfseries, text=red!80!black] {$e_{12}$} (v2);
    \draw[very thick, red!70!black] (v2) -- node[below, font=\scriptsize\bfseries, text=red!80!black] {$e_{23}$} (v3);
    
    \node[font=\scriptsize, align=center, text=gray!85!black] at (-3.0, 0.05) {コスト関数：\\[2pt] {\tiny $\displaystyle C = \sum_{\langle j,k \rangle} \frac{I - Z_j Z_k}{2}$}};

    % Arrow from Graph to Circuit
    \draw[->, very thick, gray!60] (-0.65, 1.8) -- (-0.15, 1.8);

    % Right side: Quantum Circuit Wires (y = 3.6, 2.4, 1.2, 0.0)
    \foreach \y/\i in {3.6/0, 2.4/1, 1.2/2, 0.0/3} {
        \node[left, font=\small] at (0.0, \y) {$\ket{0}_{q_\i}$};
        \draw[thick] (0.0, \y) -- (15.9, \y);
    }

    % Initial Stage: H gates
    \draw[dashed, red!60!black, thick, fill=red!3, rounded corners=1.5mm] (0.4, -0.65) rectangle (1.7, 4.55);
    \node[above, font=\scriptsize\bfseries, text=red!75!black] at (1.05, 4.6) {初期化};
    \foreach \y in {3.6, 2.4, 1.2, 0.0} {
        \draw[fill=green!20, draw=green!70!black, thick, rounded corners=1mm] (0.7, \y-0.35) rectangle (1.4, \y+0.35) node[midway, font=\footnotesize\bfseries] {$H$};
    }

    % Cost Layer: 3 Edge Interactions
    \draw[dashed, purple!70!black, thick, fill=purple!3, rounded corners=1.5mm] (2.0, -0.65) rectangle (12.3, 4.55);
    \node[above, font=\scriptsize\bfseries, text=purple!80!black] at (7.15, 4.6) {コスト層 $U(C, \gamma)$：各辺 $e_{jk}$ の $R_{ZZ}(-\gamma)$};

    % Edge e_01 on (q0, q1)
    \node[fill=purple!12, draw=purple!40, text=purple!90!black, font=\tiny\bfseries, inner sep=2pt, rounded corners=0.8mm] at (3.9, 4.08) {辺 $e_{01}$};
    \filldraw[black] (2.8, 3.6) circle (2.5pt);
    \draw[thick] (2.8, 3.6) -- (2.8, 2.4);
    \draw[thick, fill=white] (2.8, 2.4) circle (4.5pt);
    \draw[thick] (2.8, 2.24) -- (2.8, 2.56); \draw[thick] (2.64, 2.4) -- (2.96, 2.4);
    \draw[fill=orange!20, draw=orange!80!black, thick, rounded corners=1mm] (3.2, 2.05) rectangle (4.6, 2.75) node[midway, font=\scriptsize\bfseries] {$R_Z(-\gamma)$};
    \filldraw[black] (5.0, 3.6) circle (2.5pt);
    \draw[thick] (5.0, 3.6) -- (5.0, 2.4);
    \draw[thick, fill=white] (5.0, 2.4) circle (4.5pt);
    \draw[thick] (5.0, 2.24) -- (5.0, 2.56); \draw[thick] (4.84, 2.4) -- (5.16, 2.4);

    % Edge e_12 on (q1, q2)
    \node[fill=purple!12, draw=purple!40, text=purple!90!black, font=\tiny\bfseries, inner sep=2pt, rounded corners=0.8mm] at (7.2, 3.0) {辺 $e_{12}$};
    \filldraw[black] (6.1, 2.4) circle (2.5pt);
    \draw[thick] (6.1, 2.4) -- (6.1, 1.2);
    \draw[thick, fill=white] (6.1, 1.2) circle (4.5pt);
    \draw[thick] (6.1, 1.04) -- (6.1, 1.36); \draw[thick] (5.94, 1.2) -- (6.26, 1.2);
    \draw[fill=orange!20, draw=orange!80!black, thick, rounded corners=1mm] (6.5, 0.85) rectangle (7.9, 1.55) node[midway, font=\scriptsize\bfseries] {$R_Z(-\gamma)$};
    \filldraw[black] (8.3, 2.4) circle (2.5pt);
    \draw[thick] (8.3, 2.4) -- (8.3, 1.2);
    \draw[thick, fill=white] (8.3, 1.2) circle (4.5pt);
    \draw[thick] (8.3, 1.04) -- (8.3, 1.36); \draw[thick] (8.14, 1.2) -- (8.46, 1.2);

    % Edge e_23 on (q2, q3)
    \node[fill=purple!12, draw=purple!40, text=purple!90!black, font=\tiny\bfseries, inner sep=2pt, rounded corners=0.8mm] at (10.5, 1.8) {辺 $e_{23}$};
    \filldraw[black] (9.4, 1.2) circle (2.5pt);
    \draw[thick] (9.4, 1.2) -- (9.4, 0.0);
    \draw[thick, fill=white] (9.4, 0.0) circle (4.5pt);
    \draw[thick] (9.4, -0.16) -- (9.4, 0.16); \draw[thick] (9.24, 0.0) -- (9.56, 0.0);
    \draw[fill=orange!20, draw=orange!80!black, thick, rounded corners=1mm] (9.8, -0.35) rectangle (11.2, 0.35) node[midway, font=\scriptsize\bfseries] {$R_Z(-\gamma)$};
    \filldraw[black] (11.6, 1.2) circle (2.5pt);
    \draw[thick] (11.6, 1.2) -- (11.6, 0.0);
    \draw[thick, fill=white] (11.6, 0.0) circle (4.5pt);
    \draw[thick] (11.6, -0.16) -- (11.6, 0.16); \draw[thick] (11.44, 0.0) -- (11.76, 0.0);

    % Mixer Layer: Rx gates
    \draw[dashed, blue!70!black, thick, fill=blue!3, rounded corners=1.5mm] (12.6, -0.65) rectangle (14.5, 4.55);
    \node[above, font=\scriptsize\bfseries, text=blue!80!black] at (13.55, 4.6) {ミキサー層};
    \foreach \y in {3.6, 2.4, 1.2, 0.0} {
        \draw[fill=blue!15, draw=blue!75!black, thick, rounded corners=1mm] (12.8, \y-0.35) rectangle (14.3, \y+0.35) node[midway, font=\scriptsize\bfseries] {$R_X(2\beta)$};
    }

    % Measurement
    \draw[dashed, gray!70!black, thick, fill=gray!3, rounded corners=1.5mm] (14.8, -0.65) rectangle (16.0, 4.55);
    \node[above, font=\scriptsize\bfseries, text=gray!80!black] at (15.4, 4.6) {測定};
    \foreach \y in {3.6, 2.4, 1.2, 0.0} {
        \draw[fill=gray!20, draw=black, thick, rounded corners=1mm] (15.0, \y-0.35) rectangle (15.8, \y+0.35);
        \draw[thick] (15.15, \y-0.12) arc (180:0:2.5mm);
        \draw[thick, ->] (15.4, \y-0.12) -- (15.65, \y+0.20);
    }
\end{tikzpicture}
\caption{4頂点パスグラフに対する深さ $p=1$ の完全展開QAOA量子回路．問題グラフの各辺 $e_{jk}$ が，回路上の対応する量子ビット対 $(q_j, q_k)$ 間のイジング相互作用ブロック（CNOT--$R_Z(-\gamma)$--CNOT）に1対1で変換される．これにより，アルゴリズムの全ゲート実行手順が物理ハードウェア上で完全に可視化される．}
\label{fig:qaoa_maxcut4_circuit}
\end{figure}

\subsection{4. QAOAにおける「回転角」$(\gamma, \beta)$ の正体：時間発展・ブロッホ球・干渉制御}
QAOAの文献において，変分パラメータ $\gamma$ および $\beta$ は「発展時間」や「重みパラメータ」だけでなく，しばしば「回転角（rotation angle）」と呼ばれる．初学者が抱きやすい疑問として，「なぜ物理的な時間発展が『角度』と呼ばれるのか？」「具体的に何を，どの空間で，何度回転させているのか？」という点が挙げられる．ここでは，その数学的・幾何学的・アルゴリズム的意味を図\ref{fig:qaoa_rotation_angles}とともに厳密かつ直観的に解明する．

\noindent\textbf{1. なぜ「発展時間」が「回転角」と呼ばれるのか？（代数的必然性）}\\[1.5mm]
一般に量子力学において，ハミルトニアン $H$ による時間 $t$ の時間発展はユニタリ演算子 $U(t) = e^{-i H t / \hbar}$（以下 $\hbar = 1$）で与えられる．
ここで，任意のパウリ行列 $P \in \{X, Y, Z\}$ はその自乗が単位行列となる性質（対合性：$P^2 = I$）を持つ．この性質を用いて指数関数 $e^{-i \theta P}$ をテイラー展開すると，偶数次項（$P^{2k} = I$）と奇数次項（$P^{2k+1} = P$）が綺麗に分離する：
\begin{equation}
    e^{-i \theta P} = \sum_{n=0}^\infty \frac{(-i\theta P)^n}{n!} = \left( \sum_{k=0}^\infty \frac{(-1)^k \theta^{2k}}{(2k)!} \right) I - i \left( \sum_{k=0}^\infty \frac{(-1)^k \theta^{2k+1}}{(2k+1)!} \right) P = \cos\theta \, I - i \sin\theta \, P
\end{equation}
これはオイラーの公式 $e^{-i\theta} = \cos\theta - i\sin\theta$ の演算子版そのものである．
一方，ブロッホ球において回転軸 $P$ の周りに状態ベクトルを角度 $\phi$ だけ右ネジ方向に回転させる量子ゲート $R_P(\phi)$ の定義式は以下で与えられる：
\begin{equation}
    R_P(\phi) \equiv \exp\left( -i \frac{\phi}{2} P \right) = \cos\frac{\phi}{2} I - i \sin\frac{\phi}{2} P
\end{equation}
両式を比較すると，ハミルトニアン発展時間 $t$ を持つユニタリ演算子は，回転角 $\phi = 2t$ を持つ幾何学的回転ゲートと完全に恒等であることがわかる：
\begin{equation}
    e^{-i \beta X} = R_X(2\beta), \qquad e^{-i \gamma Z} = R_Z(2\gamma)
\end{equation}
すなわち，QAOAにおいてハミルトニアンの作用時間 $\gamma$ や $\beta$ を変化させることは，ヒルベルト空間およびブロッホ球上において状態ベクトルを角度 $2\gamma$ や $2\beta$ だけ幾何学的に回転させることと完全に同義である．これが「発展時間」が「回転角」と呼ばれる数学的理由である．

\medskip
\noindent\textbf{2. コスト層の回転角 $\gamma$：目的関数値を「複素位相の回転角」に焼き付ける}\\[1.5mm]
コスト演算子 $C$ は計算基底 $\ket{z}$ の対角行列であり，$C \ket{z} = C(z) \ket{z}$ を満たす．したがって，コスト層 $U(C, \gamma) = e^{-i \gamma C}$ を一般の重ね合わせ状態に作用させると，各計算基底の複素確率振幅に対して以下の変換が生じる（図\ref{fig:qaoa_rotation_angles}(a)）：
\begin{equation}
    U(C, \gamma) \sum_{z \in \{0,1\}^n} c_z \ket{z} = \sum_{z \in \{0,1\}^n} \left( c_z \, e^{-i \gamma C(z)} \right) \ket{z}
\end{equation}
\begin{itemize}
    \item \textbf{何を回転させているか}：各候補解ビット列 $z$ の複素振幅の\textbf{位相角（Phase Angle）}を，複素平面の原点周りに $\Delta \theta_z = -\gamma C(z)$ だけ回転させている．ブロッホ球の描像では，各量子ビットを $Z$ 軸周りに角度 $2\gamma$（2体項ではパリティに応じた相対位相差）だけ回転（歳差運動）させることに対応する．
    \item \textbf{測定確率はどうなるか}：注目すべきは，各状態の出現確率が $|c_z e^{-i \gamma C(z)}|^2 = |c_z|^2$ であり，\textbf{この段階では測定確率は1ミリも変化していない}点である．位相は回転したが，状態の絶対値の二乗（確率）は初期の一様分布のままである．
    \item \textbf{直観的役割}：$\gamma$ は，問題の目的関数値 $C(z)$ の高低差を，量子状態の幾何学的な「相対位相差」へと忠実に写像（エンコード）するためのゲイン調整ツマミである．
\end{itemize}

\medskip
\noindent\textbf{3. ミキサー層の回転角 $\beta$：基底間の振幅を混合し「量子干渉」を起こす}\\[1.5mm]
ミキサー演算子 $B = \sum_{j=1}^n X_j$ によるユニタリ変換 $U(B, \beta) = \prod_{j=1}^n R_{X_j}(2\beta)$ は，各量子ビットをブロッホ球の \textbf{$X$ 軸周りに角度 $2\beta$ だけ物理回転}させる（図\ref{fig:qaoa_rotation_angles}(b)）：
\begin{equation}
    R_X(2\beta) \ket{0} = \cos\beta \ket{0} - i \sin\beta \ket{1}, \qquad R_X(2\beta) \ket{1} = -i \sin\beta \ket{0} + \cos\beta \ket{1}
\end{equation}
\begin{itemize}
    \item \textbf{何を回転させているか}：計算基底 $\ket{0}$ と $\ket{1}$ の間の振幅を物理的にシャッフル（回転）させ，状態間の量子トンネリング（状態遷移）を引き起こしている．$\beta = \pi/4$ のとき回転角は $2\beta = \pi/2$（90度回転）となり，$\ket{0}$ と $\ket{1}$ が等量に混合される．
    \item \textbf{直観的役割}：$\beta$ は，異なるビット列状態同士の間で確率振幅をどの程度強く混ぜ合わせるか（干渉させるか）をコントロールする混合ツマミである．
\end{itemize}

\medskip
\noindent\textbf{4. 2つの回転角が織りなす「干渉エンジン」の本質（位相差から振幅差への変換）}\\[1.5mm]
なぜコスト層（$Z$ 回転）とミキサー層（$X$ 回転）を交互に作用させると，最適解の確率が増大するのか？
その答えが\textbf{量子干渉（Quantum Interference）}である（図\ref{fig:qaoa_rotation_angles}(c)）：
\begin{enumerate}
    \item コスト層 $\gamma$ が，各ビット列 $z$ に目的関数値 $C(z)$ に比例した「位相の回転角」の差をつける．
    \item 直後のミキサー層 $\beta$ が $X$ 軸周りに回転させて異なるビット列の振幅を混合させると，\textbf{回転させた位相差に応じて波の干渉（強め合い・打ち消し合い）が発生する}．
    \item 最適な回転角 $(\gamma^*, \beta^*)$ を選ぶと，目的関数値が良い（低コストな）基底状態への遷移経路同士の位相が揃って\textbf{建設的干渉（Constructive Interference）}を起こして確率振幅が極大化し，逆に悪い基底状態は逆位相となって\textbf{相殺的干渉（Destructive Interference）}により消滅する．
\end{enumerate}
すなわち，QAOAにおける古典最適化とは，\textbf{「コスト値によって刻印された位相ベクトル群が，ミキサー回転を経た後に最適解の方向へ一斉に強め合って集束（共鳴）するような，絶妙な2種類の回転角の組み合わせ $(\vec{\gamma}, \vec{\beta})$ を探索する作業」}に他ならない．

\begin{figure}[htbp]
\centering
\begin{tikzpicture}[scale=0.88, >=latex]
    % =========================================================================
    % Panel (a): Cost layer gamma (Z rotation / Phase encoding)
    % =========================================================================
    \draw[fill=purple!3, draw=purple!45, thick, rounded corners=3mm] (0.0, -0.6) rectangle (4.6, 5.2);
    \node[fill=purple!15, draw=purple!60, font=\scriptsize\bfseries, text=purple!90!black, rounded corners=1mm, inner sep=2.5pt] at (2.3, 4.85) {(a) コスト層：$Z$ 軸位相回転};

    % Bloch Sphere (a) centered at (2.3, 2.75), R=1.15
    \draw[thick, gray!50] (2.3, 2.75) circle (1.15);
    \draw[dashed, gray!35] (2.3, 2.75) ellipse (1.15 and 0.35);
    \draw[thick, gray!45] (3.45, 2.75) arc (0:-180:1.15 and 0.35);

    % Z axis (stops cleanly with margin below banner)
    \draw[->, thick, gray!75] (2.3, 1.40) -- (2.3, 4.15) node[above=1pt, font=\tiny\bfseries, text=gray!80!black] {$Z$};
    \node[right=2pt, font=\tiny\bfseries, text=gray!70!black] at (2.3, 3.75) {$\ket{0}$};
    \node[right=2pt, font=\tiny\bfseries, text=gray!70!black] at (2.3, 1.75) {$\ket{1}$};

    % X axis (towards front-left)
    \draw[->, thick, gray!50] (2.3, 2.75) -- (1.40, 2.38) node[above left=-1pt, font=\tiny\bfseries, text=gray!70!black] {$X$};

    % Initial state |+> (on equator along +X axis)
    \draw[->, very thick, blue!85!black] (2.3, 2.75) -- (1.48, 2.41)
        node[below left=1pt, font=\scriptsize\bfseries, text=blue!90!black] {$\ket{+}$};

    % Rotated state e^{-i gamma C}|+> (rotated along front equator)
    \draw[->, very thick, red!80!black] (2.3, 2.75) -- (3.05, 2.50)
        node[below right=1pt, font=\scriptsize\bfseries, text=red!90!black] {$e^{-i\gamma C}\ket{+}$};

    % Rotation arc along front equator
    \draw[->, very thick, purple!85!black] (1.75, 2.50) to[bend right=25] (2.85, 2.54);
    \node[font=\tiny\bfseries, text=purple!90!black] at (2.30, 2.05) {回転角 $2\gamma$};

    % Text explanation (a)
    \draw[fill=white, draw=purple!25, rounded corners=1.5mm] (0.2, -0.45) rectangle (4.4, 0.95);
    \node[font=\tiny, align=left, text width=3.9cm] at (2.3, 0.25) {
        \textbf{幾何学的変換：} $\ket{z} \mapsto e^{-i \gamma C(z)} \ket{z}$\\[1mm]
        ・ブロッホ球の $Z$ 軸周りに角度 $2\gamma$ 位相回転\\
        ・\textbf{測定確率 $|c_z|^2$ は不変}（赤道面を周回）\\
        ・目的関数値 $C(z)$ を「位相差」として刻印
    };

    % =========================================================================
    % Panel (b): Mixer layer beta (X rotation / Amplitude tunneling)
    % =========================================================================
    \draw[fill=blue!3, draw=blue!45, thick, rounded corners=3mm] (5.0, -0.6) rectangle (9.6, 5.2);
    \node[fill=blue!15, draw=blue!60, font=\scriptsize\bfseries, text=blue!90!black, rounded corners=1mm, inner sep=2.5pt] at (7.3, 4.85) {(b) ミキサー層：$X$ 軸振幅回転};

    % Bloch Sphere (b) centered at (7.3, 2.75), R=1.15
    \draw[thick, gray!50] (7.3, 2.75) circle (1.15);
    \draw[dashed, gray!35] (7.3, 2.75) ellipse (1.15 and 0.35);

    % Z axis (stops cleanly with margin below banner)
    \draw[->, thick, gray!45] (7.3, 1.40) -- (7.3, 4.15) node[above=1pt, font=\tiny\bfseries, text=gray!80!black] {$Z$};
    \node[right=2pt, font=\tiny\bfseries, text=gray!70!black] at (7.3, 1.75) {$\ket{1}$};

    % X axis (rotation axis: securely inside boundary)
    \draw[very thick, blue!85!black] (7.7, 2.90) -- (6.35, 2.40);
    \draw[->, very thick, blue!85!black] (6.35, 2.40) -- (6.10, 2.30) node[below left=-1pt, font=\tiny\bfseries] {$X$};
    \node[font=\fontsize{5.5}{6.5}\selectfont\bfseries, text=blue!90!black] at (6.25, 2.75) {(回転軸)};

    % Initial state |0> (at north pole: cleanly positioned above left)
    \draw[->, very thick, blue!85!black] (7.3, 2.75) -- (7.3, 3.90);
    \node[above left=0pt, font=\scriptsize\bfseries, text=blue!90!black] at (7.3, 3.85) {$\ket{0}$};

    % Rotated state (tilted along Y-Z circle around X axis)
    \draw[->, very thick, red!80!black] (7.3, 2.75) -- (8.15, 3.50)
        node[above right=-1pt, font=\scriptsize\bfseries, text=red!90!black] {$e^{-i\beta X}\ket{0}$};

    % Rotation arc around X axis (along sphere contour)
    \draw[->, very thick, orange!85!black] (7.3, 3.75) arc (90:42:1.0);
    \node[font=\tiny\bfseries, text=orange!90!black] at (8.15, 3.95) {回転角 $2\beta$};

    % Text explanation (b)
    \draw[fill=white, draw=blue!25, rounded corners=1.5mm] (5.2, -0.45) rectangle (9.4, 0.95);
    \node[font=\tiny, align=left, text width=3.9cm] at (7.3, 0.25) {
        \textbf{幾何学的変換：} $R_X(2\beta) = e^{-i \beta X}$\\[1mm]
        ・ブロッホ球の $X$ 軸周りに角度 $2\beta$ 物理回転\\
        ・$\ket{0} \leftrightarrow \ket{1}$ 間の振幅を混合（トンネリング）\\
        ・基底状態間の\textbf{干渉（重なり合い）}を発生
    };

    % =========================================================================
    % Panel (c): Quantum Interference (Phase to Probability concentration)
    % =========================================================================
    \draw[fill=teal!3, draw=teal!45, thick, rounded corners=3mm] (10.0, -0.6) rectangle (14.6, 5.2);
    \node[fill=teal!15, draw=teal!60, font=\scriptsize\bfseries, text=teal!90!black, rounded corners=1mm, inner sep=2.5pt] at (12.3, 4.85) {(c) 量子干渉：位相差 $\to$ 確率濃縮};

    % Constructive interference (Optimal state) - perfectly spaced
    \node[anchor=west, font=\tiny\bfseries, text=teal!90!black] at (10.25, 4.40) {【最適解 $z^*$】：同位相で強め合う};
    % Vector 1 & 2
    \draw[->, thick, purple!80!black] (10.4, 3.78) -- (11.45, 3.78)
        node[midway, above=1pt, font=\fontsize{6}{7}\selectfont] {$\psi_1$};
    \draw[->, thick, blue!80!black] (11.45, 3.78) -- (12.6, 3.78)
        node[midway, above=1pt, font=\fontsize{6}{7}\selectfont] {$\psi_2$};
    % Result vector
    \draw[->, very thick, teal!90!black] (10.4, 3.32) -- (12.6, 3.32)
        node[right=2pt, font=\tiny\bfseries] {合成大！};
    \node[anchor=west, font=\fontsize{5.5}{6.5}\selectfont, text=teal!85!black] at (10.4, 2.95)
        {$\implies$ 建設的干渉：確率 $P(z^*)$ が極大化};

    % Destructive interference (Suboptimal state) - perfectly spaced
    \node[anchor=west, font=\tiny\bfseries, text=red!80!black] at (10.25, 2.45) {【非最適解 $z'$】：逆位相で相殺};
    % Vector 1 & 2
    \draw[->, thick, purple!80!black] (10.4, 1.95) -- (11.45, 1.95)
        node[midway, above=1pt, font=\fontsize{6}{7}\selectfont] {$\psi_1$};
    \draw[->, thick, blue!80!black] (11.45, 1.95) -- (10.6, 1.95)
        node[midway, below=1pt, font=\fontsize{5}{6}\selectfont] {$\psi_2$ (逆位相)};
    % Result vector
    \draw[->, very thick, red!80!black] (10.4, 1.45) -- (10.6, 1.45)
        node[right=2pt, font=\tiny\bfseries] {合成 $\approx 0$};
    \node[anchor=west, font=\fontsize{5.5}{6.5}\selectfont, text=red!75!black] at (10.4, 1.10)
        {$\implies$ 相殺的干渉：確率 $P(z')$ が消滅};

    % Text explanation (c)
    \draw[fill=white, draw=teal!25, rounded corners=1.5mm] (10.2, -0.45) rectangle (14.4, 0.95);
    \node[font=\tiny, align=left, text width=3.9cm] at (12.3, 0.25) {
        \textbf{QAOAの本質メカニズム：}\\[1mm]
        ・$\gamma$ の回転でつけた\textbf{位相差}を，$\beta$ の回転で\textbf{干渉}させて測定確率の偏りへと変換\\
        ・最適角 $(\gamma^*, \beta^*)$ は\textbf{最適解の共鳴条件}
    };
\end{tikzpicture}
\caption{QAOAにおける回転角 $(\gamma, \beta)$ の物理的・幾何学的意味と量子干渉メカニズム．(a) コスト層 $U(C, \gamma)$ は計算基底 $\ket{z}$ の複素位相を角度 $\theta_z = -\gamma C(z)$（ブロッホ球の $Z$ 軸周りに角度 $2\gamma$）回転させてコスト差を「位相差」に変換する（測定確率は不変）．(b) ミキサー層 $U(B, \beta)$ はブロッホ球の $X$ 軸周りに角度 $2\beta$ だけ物理的に回転させ，基底間の振幅を混合（トンネリング）させる．(c) 2つの回転角が協調することで，刻印された位相差が量子干渉を起こし，最適解（低コスト状態）の確率振幅が建設的干渉により増幅され，非最適解が相殺される．}
\label{fig:qaoa_rotation_angles}
\end{figure}

\subsection{5. 期待値と古典・量子協調ループ}
目的関数 $F_p(\vec{\gamma}, \vec{\beta})$ は，状態 $\ket{\vec{\gamma}, \vec{\beta}}$ におけるコスト演算子 $C$ の期待値として定義される：
\begin{equation}
    F_p(\vec{\gamma}, \vec{\beta}) = \mel{\vec{\gamma}, \vec{\beta}}{C}{\vec{\gamma}, \vec{\beta}}
\end{equation}
QAOAは図\ref{fig:vqa_loop}に示すハイブリッドループによって最適化を進める．

\begin{figure}[htbp]
\centering
\begin{tikzpicture}[scale=0.92, >=latex]
    % QPU Box (x: 0 to 5.8, y: 0 to 4.8)
    \draw[fill=cyan!6, draw=cyan!80!black, thick, rounded corners=3mm] (0, 0) rectangle (5.8, 4.8);
    % QPU Title bar
    \draw[fill=cyan!18, draw=cyan!80!black, thick, rounded corners=3mm] (0, 3.9) rectangle (5.8, 4.8)
        node[midway, font=\bfseries\large, text=cyan!90!black] {量子プロセッサ (QPU)};
    
    \node[font=\footnotesize, align=center, text width=5.4cm] at (2.9, 2.95) {
        \textbf{【量子回路の実行】}\\[1mm]
        $\ket{\vec{\gamma}, \vec{\beta}} = \prod_{k=1}^p U_B(\beta_k) U_C(\gamma_k) \ket{s}$
    };
    
    \draw[dashed, cyan!50!black] (0.4, 2.05) -- (5.4, 2.05);
    
    \node[font=\footnotesize, align=center, text width=5.4cm] at (2.9, 1.05) {
        \textbf{【計算基底サンプリング】}\\[1mm]
        ($1,000 \sim 10,000$ shots)\\
        ビット列サンプル $\{z^{(1)}, z^{(2)}, \dots\}$
    };

    % CPU Box (x: 8.6 to 14.4, y: 0 to 4.8)
    \draw[fill=orange!6, draw=orange!85!black, thick, rounded corners=3mm] (8.6, 0) rectangle (14.4, 4.8);
    % CPU Title bar
    \draw[fill=orange!18, draw=orange!85!black, thick, rounded corners=3mm] (8.6, 3.9) rectangle (14.4, 4.8)
        node[midway, font=\bfseries\large, text=orange!95!black] {古典コンピュータ (CPU)};
    
    \node[font=\footnotesize, align=center, text width=5.4cm] at (11.5, 2.95) {
        \textbf{【エネルギー期待値の推定】}\\[1mm]
        $F_p(\vec{\gamma}, \vec{\beta}) \approx \frac{1}{K}\sum_{k=1}^K C(z^{(k)})$
    };
    
    \draw[dashed, orange!50!black] (9.0, 2.05) -- (14.0, 2.05);
    
    \node[font=\footnotesize, align=center, text width=5.4cm] at (11.5, 1.05) {
        \textbf{【古典オプティマイザ】}\\[1mm]
        (COBYLA / SPSA / 勾配法)\\
        パラメータ更新: $(\vec{\gamma}, \vec{\beta}) \leftarrow (\vec{\gamma}', \vec{\beta}')$
    };

    % High clearance interactive arrows (Gap: 5.8 to 8.6 = 2.8cm)
    \draw[->, very thick, purple!80!black] (5.8, 3.0) -- (8.6, 3.0)
        node[midway, above, font=\scriptsize\bfseries, text=purple!85!black] {測定サンプル列}
        node[midway, below, font=\tiny, text=purple!75!black] {$\{z^{(k)}\}$};
        
    \draw[<-, very thick, blue!80!black] (5.8, 1.1) -- (8.6, 1.1)
        node[midway, above, font=\scriptsize\bfseries, text=blue!85!black] {新パラメータ}
        node[midway, below, font=\tiny, text=blue!75!black] {$(\vec{\gamma}, \vec{\beta})$};
\end{tikzpicture}
\caption{QAOAにおける量子・古典ハイブリッド最適化ループ（VQE/VQA型）．QPUが量子重ね合わせから高速にサンプルを生成し，古典CPUがエネルギーを評価してパラメータ $(\vec{\gamma}, \vec{\beta})$ を更新する．}
\label{fig:vqa_loop}
\end{figure}

\section{古典角度最適化：理論・オプティマイザ・最新戦略}
\label{sec:classical_optimization}

QAOAは「量子回路による状態生成」と「古典オプティマイザによるパラメータ更新」が車の両輪となって機能する量子・古典ハイブリッドアルゴリズム（VQA\cite{cerezo2021variational}）である．
量子回路がいかに優れた表現力（Expressibility）を持っていても，変分パラメータ $\vec{\gamma} = (\gamma_1, \dots, \gamma_p)$ および $\vec{\beta} = (\beta_1, \dots, \beta_p)$ を適切に最適化できなければ，アルゴリズムは真価を発揮できず，単なる無秩序な重ね合わせ状態を出力するに留まる．

深さ $p$ のQAOAにおいて，古典オプティマイザが担う数学的タスクは，以下の $2p$ 次元実数パラメータ空間における非線形連続最適化問題として定式化される：
\begin{equation}
    \min_{\vec{\gamma}, \vec{\beta} \in \mathbb{R}^{2p}} F_p(\vec{\gamma}, \vec{\beta}) = \mel{\vec{\gamma}, \vec{\beta}}{H_C}{\vec{\gamma}, \vec{\beta}}
\end{equation}
物理的な量子ハードウェア（QPU）上では，波動関数そのものを直接観測することはできないため，計算基底での射影測定を有限のショット数 $K$ 回実行し，ビット列サンプル $\{z^{(1)}, z^{(2)}, \dots, z^{(K)}\}$ を取得する．古典CPUは，このサンプルから計算される標本平均推定量：
\begin{equation}
    \hat{F}_p(\vec{\gamma}, \vec{\beta}) = \frac{1}{K} \sum_{k=1}^K C(z^{(k)})
\end{equation}
を目的関数値として受け取り，次の一手となるパラメータ $(\vec{\gamma}', \vec{\beta}')$ を決定してQPUへとフィードバックする．

\begin{insightbox}[サンプリング目的関数の発展：CVaR-QAOA \cite{barkoutsos2020improving}]
通常のQAOAでは「エネルギー期待値（平均値）」の最小化を目指すが，実社会の組合せ最適化において真に必要なのは「平均値が良いこと」ではなく，\textbf{「最良解（厳密解または超高精度近似解）のビット列が少なくとも1回サンプルされること」}である．
Barkoutsosら (2020)\cite{barkoutsos2020improving} は，金融工学のリスク管理指標である \textbf{CVaR（Conditional Value-at-Risk）} をQAOAの目的関数として導入した．
全 $K$ ショットの測定サンプルをコスト関数の昇順（良質な順）にソートした系列 $C(z_{(1)}) \le C(z_{(2)}) \le \dots \le C(z_{(K)})$ に対し，最良の上位 $\alpha$ 割（例: $\alpha = 0.1 \sim 0.25$）の標本平均：
\begin{equation}
    F_{\text{CVaR}}^{(\alpha)}(\vec{\gamma}, \vec{\beta}) = \frac{1}{\lceil \alpha K \rceil} \sum_{k=1}^{\lceil \alpha K \rceil} C(z_{(k)})
\end{equation}
を古典オプティマイザの最小化対象とする．
これにより，オプティマイザは全体の分布を無理にシフトさせるのではなく，\textbf{「基底状態近傍の確率密度の裾野を尖らせる」}方向へとパラメータを誘導する．浅い深さ $p$ においても基底状態の抽出確率が劇的に向上することが知られている．
\end{insightbox}

\subsection{1. 角度最適化を阻む3大ボトルネック}
目的関数 $F_p(\vec{\gamma}, \vec{\beta})$ を最適化する古典タスクは，機械学習や通常の非線形計画問題と比べて特有の困難を抱えている（図 \ref{fig:classical_bottlenecks}）：

\begin{figure}[htbp]
\centering
\begin{tikzpicture}[scale=0.88, >=latex]
    % Panel A: Non-Convex Landscape & Local Minima
    \begin{scope}[shift={(0,0)}]
        \draw[fill=gray!5, draw=gray!40, rounded corners=2.5mm] (-0.3, -0.6) rectangle (4.7, 4.4);
        \node[fill=gray!15, draw=gray!50, rounded corners=1mm, font=\scriptsize\bfseries, text=black, inner sep=2.2pt] at (2.2, 4.05) {(a) 複雑な非凸ランドスケープ};
        
        \draw[->, thick] (0, 0) -- (4.3, 0) node[right, font=\scriptsize] {$\theta$};
        \draw[->, thick] (0, 0) -- (0, 3.4);
        \node[left, font=\scriptsize] at (0, 3.2) {$F(\theta)$};
        
        % Non-convex curve
        \draw[very thick, blue!80!black, smooth] 
            plot coordinates {(0.2, 2.6) (0.8, 1.3) (1.4, 2.1) (2.0, 0.8) (2.7, 2.9) (3.5, 0.35) (4.1, 2.3)};
            
        % Local Minima and Global Minimum
        \filldraw[red!80!black] (0.8, 1.3) circle (2.2pt);
        \node[above, font=\tiny, text=red!80!black] at (0.8, 1.35) {局所解};
        
        \filldraw[red!80!black] (2.0, 0.8) circle (2.2pt);
        \node[above, font=\tiny, text=red!80!black] at (2.0, 0.85) {局所解};
        
        \filldraw[green!60!black] (3.5, 0.35) circle (2.8pt);
        \node[below, font=\tiny\bfseries, text=green!60!black] at (3.5, 0.25) {大域的最適解};
        
        % Trapped particle
        \draw[fill=orange!80, draw=orange!90!black] (0.5, 1.9) circle (2.8pt);
        \draw[->, thick, orange!90!black] (0.5, 1.8) to[bend right=15] (0.75, 1.35);
        \node[left, font=\tiny, text=orange!90!black] at (0.45, 2.1) {初期値};
        \node[font=\tiny\bfseries, align=center, text width=4.4cm, text=gray!80!black] at (2.2, -0.3) {多峰性による局所解トラップ};
    \end{scope}

    % Panel B: Shot Noise & Gradient Breakdown
    \begin{scope}[shift={(5.3,0)}]
        \draw[fill=gray!5, draw=gray!40, rounded corners=2.5mm] (-0.3, -0.6) rectangle (4.7, 4.4);
        \node[fill=gray!15, draw=gray!50, rounded corners=1mm, font=\scriptsize\bfseries, text=black, inner sep=2.2pt] at (2.2, 4.05) {(b) ショットノイズと微分破綻};
        
        \draw[->, thick] (0, 0) -- (4.3, 0) node[right, font=\scriptsize] {$\theta$};
        \draw[->, thick] (0, 0) -- (0, 3.4);
        \node[left, font=\scriptsize] at (0, 3.2) {$\hat{F}(\theta)$};
        
        % True function (dashed)
        \draw[thick, dashed, blue!50] (0.3, 2.8) to[out=-40, in=180] (2.4, 1.1) to[out=0, in=-140] (4.1, 2.6);
        \node[blue!70, font=\tiny] at (3.6, 2.9) {真の期待値};
        
        % Noisy points with error bars
        \foreach \x/\y in {0.6/2.4, 1.1/1.9, 1.6/1.35, 2.1/1.15, 2.6/1.05, 3.1/1.4, 3.6/1.9} {
            \draw[gray!70, thin] (\x, \y-0.22) -- (\x, \y+0.22);
            \draw[gray!70, thin] (\x-0.07, \y-0.22) -- (\x+0.07, \y-0.22);
            \draw[gray!70, thin] (\x-0.07, \y+0.22) -- (\x+0.07, \y+0.22);
            \filldraw[blue!80!black] (\x, \y) circle (1.6pt);
        }
        
        % Gradient explosion with small step
        \filldraw[red!80!black] (1.9, 1.35) circle (1.8pt) node[below left, font=\tiny, text=red!80!black] {$\theta-\epsilon$};
        \filldraw[red!80!black] (2.3, 0.95) circle (1.8pt) node[below right, font=\tiny, text=red!80!black] {$\theta+\epsilon$};
        \draw[thick, red!80!black] (1.6, 1.65) -- (2.6, 0.65);
        \node[right, font=\tiny, text=red!80!black, align=left] at (2.6, 0.75) {ノイズによる\\勾配の乱反射};
        
        \node[font=\tiny\bfseries, align=center, text width=4.4cm, text=gray!80!black] at (2.2, -0.3) {差分刻み $\epsilon \to 0$ でノイズ発散};
    \end{scope}

    % Panel C: Barren Plateau
    \begin{scope}[shift={(10.6,0)}]
        \draw[fill=gray!5, draw=gray!40, rounded corners=2.5mm] (-0.3, -0.6) rectangle (4.7, 4.4);
        \node[fill=gray!15, draw=gray!50, rounded corners=1mm, font=\scriptsize\bfseries, text=black, inner sep=2.2pt] at (2.2, 4.05) {(c) 勾配消失 (Barren Plateau)};
        
        \draw[->, thick] (0, 0) -- (4.3, 0) node[right, font=\scriptsize] {$\theta$};
        \draw[->, thick] (0, 0) -- (0, 3.4);
        \node[left, font=\scriptsize] at (0, 3.2) {$F(\theta)$};
        
        % Golf-course flat landscape
        \draw[very thick, purple!80!black] (0.3, 2.3) -- (2.0, 2.3) 
            to[out=0, in=90] (2.2, 0.45) 
            to[out=90, in=180] (2.4, 2.3) -- (4.1, 2.3);
            
        \node[font=\tiny, text=purple!90!black, align=center] at (1.1, 2.7) {平坦な砂漠\\$\text{Var}[\partial F] \sim 2^{-N}$};
        \node[font=\tiny, text=purple!90!black, align=center] at (3.3, 2.7) {勾配 $\approx 0$\\探索不能};
        
        \filldraw[green!60!black] (2.2, 0.45) circle (2.8pt);
        \node[below, font=\tiny\bfseries, text=green!60!black] at (2.2, 0.35) {針の穴の最適解};
        
        \node[font=\tiny\bfseries, align=center, text width=4.4cm, text=gray!80!black] at (2.2, -0.3) {深さ・ノイズによる勾配の指数消失};
    \end{scope}
\end{tikzpicture}
\caption{QAOAの古典角度最適化を阻む3大ボトルネック．(a) 周期性と非線形性に起因する波打つ非凸ランドスケープと浅い局所解へのトラップ．(b) 有限ショットサンプリングの統計誤差による有限差分勾配の乱反射と発散．(c) 回路の深層化やハードウェアノイズによってパラメータ空間全域が平坦化し，勾配が指数関数的に消滅するBarren Plateau現象．}
\label{fig:classical_bottlenecks}
\end{figure}

\begin{enumerate}
    \item \textbf{エネルギーランドスケープの高度な非凸性（Non-Convexity）と対称性・周期性}:
          $F_p$ はユニタリ回転演算子の積から生じる三角関数（$\sin, \cos$）の高次多項式として構成される．コストハミルトニアンの固有値が整数である場合，$\gamma_k$ は周期 $\pi$ または $2\pi$ を持ち，ミキサー演算子のパウリ $X$ の固有値から $\beta_k$ は周期 $\pi$ を持つ．さらに，系は時間反転対称性（$F_p(-\vec{\gamma}, -\vec{\beta}) = F_p(\vec{\gamma}, \vec{\beta})$）やパリティ対称性を備えるため，パラメータ空間には無数の局所解（Local Minima / Maxima）と鞍点（Saddle Points）が網の目のように密集した複雑な地形（図 \ref{fig:classical_bottlenecks}(a)）が出現する．ランダムな初期値から局所最適化を開始すると，大域的最適解とは程遠い浅い局所解に即座に捕捉される．
          
    \item \textbf{測定の統計的ゆらぎ（ショットノイズ: Shot Noise）と数値差分の破綻}:
          QPU上で期待値を評価する際，有限のサンプリング測定（ショット数 $K$）に伴い，中心極限定理から標準誤差 $\sigma \propto \sigma_C / \sqrt{K}$ の統計誤差が不可避に乗る（図 \ref{fig:classical_bottlenecks}(b)）．通常の連続最適化で多用される有限差分法：
          \begin{equation}
              g_j(\vec{\theta}) \approx \frac{\hat{F}(\vec{\theta} + \epsilon \vec{e}_j) - \hat{F}(\vec{\theta} - \epsilon \vec{e}_j)}{2\epsilon} = \frac{\partial F}{\partial \theta_j} + \mathcal{O}(\epsilon^2) + \frac{\epsilon_{\text{shot}}^+ - \epsilon_{\text{shot}}^-}{2\epsilon}
          \end{equation}
          を適用すると，分子の統計ノイズ $\Delta \epsilon_{\text{shot}} \sim \mathcal{O}(1/\sqrt{K})$ が微小刻み幅 $\epsilon$ で割られるため，\textbf{ノイズ項の分散が $\mathcal{O}\left(\frac{1}{\epsilon^2 K}\right)$ へと急激に爆発}する．刻み幅 $\epsilon \to 0$ の極限で勾配ベクトルの方向が完全に乱反射し，最適化ルーチンは即座に発散する．
          
    \item \textbf{勾配消失問題（Barren Plateau）とノイズ誘起平坦化}:
          McCleanら (2018)\cite{mcclean2018barren} は，量子回路が深くなる，あるいはハール測度に従うランダムユニタリ（2-デザイン）に近づくと，パラメータ空間のほぼ全域において期待値の勾配の分散が量子ビット数 $N$ に対して指数関数的に縮退することを示した：
          \begin{equation}
              \text{Var}\left[ \frac{\partial F}{\partial \theta} \right] \in \mathcal{O}(2^{-N})
          \end{equation}
          これは，探索空間がほぼ完全な「平坦な砂漠（ゴルフ場のようなランドスケープ）」と化し（図 \ref{fig:classical_bottlenecks}(c)），勾配情報に基づく探索が完全に機能停止することを意味する．さらにWangら (2021)\cite{wang2021noise} は，実機NISQプロセッサのデコヒーレンスや局所パウリノイズが存在する場合，浅い深さ $p$ の回路であってもノイズの蓄積によって最大混合状態へ向かい，\textbf{ノイズ誘起Barren Plateau (NIBP)} が不可避に発生することを厳密に証明した．
\end{enumerate}

\subsection{2. 主要な古典オプティマイザの分類と特性}
QAOAで用いられる古典オプティマイザは，勾配情報の取得様式とノイズ耐性により，大きく以下の3系統に分類される（表 \ref{tab:optimizers}，図 \ref{fig:classical_optimizers}）．

\begin{figure}[htbp]
\centering
\begin{tikzpicture}[scale=0.88, >=latex]
    % Panel A: COBYLA
    \begin{scope}[shift={(0,0)}]
        \draw[fill=gray!5, draw=gray!40, rounded corners=2.5mm] (-0.3, -0.6) rectangle (4.7, 4.4);
        \node[fill=gray!15, draw=gray!50, rounded corners=1mm, font=\scriptsize\bfseries, text=black, inner sep=2.2pt] at (2.2, 4.05) {(a) COBYLA（単体近似）};
        
        % Contours
        \draw[gray!60, thin] (2.2, 1.8) ellipse (1.7cm and 1.1cm);
        \draw[gray!70, thin] (2.2, 1.8) ellipse (1.1cm and 0.7cm);
        \draw[gray!80, thin] (2.2, 1.8) ellipse (0.55cm and 0.35cm);
        \filldraw[green!60!black] (2.2, 1.8) circle (2pt) node[below right, font=\tiny] {極小値};
        
        % Simplex (triangle)
        \draw[thick, blue!80!black, fill=blue!10] (0.9, 2.9) -- (1.9, 3.1) -- (1.2, 2.1) -- cycle;
        \filldraw[blue!80!black] (0.9, 2.9) circle (2.2pt) node[above left, font=\tiny] {$\vec{x}_1$};
        \filldraw[blue!80!black] (1.9, 3.1) circle (2.2pt) node[above right, font=\tiny] {$\vec{x}_2$};
        \filldraw[blue!80!black] (1.2, 2.1) circle (2.2pt) node[left, font=\tiny] {$\vec{x}_3$};
        
        % Trust region circle
        \draw[dashed, red!70!black] (1.33, 2.7) circle (0.8cm);
        \node[font=\tiny, text=red!70!black] at (1.0, 3.6) {信頼領域};
        
        % Reflected point
        \draw[->, thick, red!80!black] (1.9, 3.1) -- (1.7, 2.0);
        \filldraw[red!80!black] (1.7, 2.0) circle (2.2pt) node[right, font=\tiny] {反転点};
        
        \node[font=\tiny\bfseries, align=center, text width=4.4cm, text=gray!80!black] at (2.2, -0.3) {勾配不要・評価回数 $\le 2p+1$\\ノイズ環境では単体が歪み停止};
    \end{scope}

    % Panel B: SPSA
    \begin{scope}[shift={(5.3,0)}]
        \draw[fill=gray!5, draw=gray!40, rounded corners=2.5mm] (-0.3, -0.6) rectangle (4.7, 4.4);
        \node[fill=gray!15, draw=gray!50, rounded corners=1mm, font=\scriptsize\bfseries, text=black, inner sep=2.2pt] at (2.2, 4.05) {(b) SPSA（同時確率摂動）};
        
        % Contours
        \draw[gray!60, thin] (2.2, 1.8) ellipse (1.7cm and 1.1cm);
        \draw[gray!70, thin] (2.2, 1.8) ellipse (1.1cm and 0.7cm);
        \draw[gray!80, thin] (2.2, 1.8) ellipse (0.55cm and 0.35cm);
        \filldraw[green!60!black] (2.2, 1.8) circle (2pt);
        
        % Current point
        \filldraw[black] (1.0, 2.7) circle (2.5pt) node[above left, font=\tiny\bfseries] {$\vec{\theta}_k$};
        
        % Perturbation directions (+c_k Delta and -c_k Delta)
        \draw[dashed, thick, orange!85!black] (0.4, 3.3) -- (1.6, 2.1);
        \filldraw[orange!90!black] (1.6, 2.1) circle (2.2pt) node[right, font=\tiny] {$\vec{\theta}_k + c_k\vec{\Delta}_k$};
        \filldraw[orange!90!black] (0.4, 3.3) circle (2.2pt) node[left, font=\tiny] {$\vec{\theta}_k - c_k\vec{\Delta}_k$};
        
        % Update step
        \draw[->, very thick, blue!85!black] (1.0, 2.7) -- (1.7, 2.0)
            node[midway, above right, font=\tiny\bfseries] {$-a_k \hat{\vec{g}}_k$};
            
        \node[font=\tiny\bfseries, align=center, text width=4.4cm, text=gray!80!black] at (2.2, -0.3) {常に \textbf{評価2回} で全勾配推定\\確率的平均化で実機ノイズに最強};
    \end{scope}

    % Panel C: Parameter-Shift Rule
    \begin{scope}[shift={(10.6,0)}]
        \draw[fill=gray!5, draw=gray!40, rounded corners=2.5mm] (-0.3, -0.6) rectangle (4.7, 4.4);
        \node[fill=gray!15, draw=gray!50, rounded corners=1mm, font=\scriptsize\bfseries, text=black, inner sep=2.2pt] at (2.2, 4.05) {(c) パラメータシフト則};
        
        % Sine curve representing F(theta)
        \draw[->, thick] (0.2, 0.4) -- (4.3, 0.4) node[right, font=\scriptsize] {$\theta$};
        \draw[->, thick] (0.2, 0.4) -- (0.2, 3.4);
        \node[left, font=\scriptsize] at (0.2, 3.2) {$F(\theta)$};
        
        \draw[very thick, purple!80!black, domain=0.5:3.9, samples=50] 
            plot (\x, {1.8 + 1.0*sin(deg(\x - 1.2))});
            
        % Evaluation at theta
        \filldraw[black] (2.2, {1.8 + 1.0*sin(deg(2.2 - 1.2))}) circle (2.2pt) node[above, font=\tiny] {$\theta$};
        
        % Macro shift points (+pi/2 and -pi/2)
        \filldraw[blue!80!black] (1.1, {1.8 + 1.0*sin(deg(1.1 - 1.2))}) circle (2.2pt)
            node[left, font=\tiny] {$\theta - \frac{\pi}{2}$};
        \filldraw[blue!80!black] (3.3, {1.8 + 1.0*sin(deg(3.3 - 1.2))}) circle (2.2pt)
            node[right, font=\tiny] {$\theta + \frac{\pi}{2}$};
            
        % Secant matching exact tangent
        \draw[thick, dashed, red!80!black] (0.9, 1.2) -- (3.5, 2.9)
            node[above left, font=\tiny, text=red!80!black] {厳密な微分値};
            
        \node[font=\tiny\bfseries, align=center, text width=4.4cm, text=gray!80!black] at (2.2, -0.3) {マクロな角度差分で厳密勾配\\有限差分のようなノイズ爆発皆無};
    \end{scope}
\end{tikzpicture}
\caption{QAOAに適用される主要古典オプティマイザの探索挙動と数理的メカニズムの対比．(a) COBYLA：$2p+1$ 個の頂点からなる単体上で目的関数を線形近似し，信頼領域内で反転・収縮を繰り返す．(b) SPSA：全パラメータを同時にランダム摂動させ，1ステップあたりわずか2回の回路実行で確率的勾配を推定する．(c) パラメータシフト則：マクロな有限角度（$\pm \pi/2$ または $\pm \pi/4$）の測定値の差から，量子状態の幾何学的性質により真の解析的導関数を厳密に導出する．}
\label{fig:classical_optimizers}
\end{figure}

\begin{table}[htbp]
\centering
\caption{QAOAに適用される主要古典オプティマイザの比較}
\label{tab:optimizers}
\vspace{-2mm}
\resizebox{\textwidth}{!}{
\begin{tabular}{lcccc}
\toprule
\textbf{手法名} & \textbf{勾配情報の利用} & \textbf{ショットノイズ耐性} & \textbf{関数評価回数 / step} & \textbf{主な特徴・適用領域} \\
\midrule
\textbf{COBYLA} & 勾配不要 (シンプレックス) & 中程度 & 高々 $2p + 1$ 回 & 最も標準的．$p \le 3$ かつ状態ベクトルシミュレータで最速． \\
\textbf{Nelder-Mead} & 勾配不要 (シンプレックス) & 弱い & $1 \sim 2$ 回 & ノイズにより単体が歪んで早期停止しやすい．実機非推奨． \\
\textbf{SPSA} & 確率的疑似勾配 & \textbf{極めて高い} & \textbf{常に 2 回} & NISQ実機・ショットノイズ環境のゴールドスタンダード． \\
\textbf{Adam / L-BFGS} & 厳密勾配 (Shift-Rule) & 高い & $2 \times 2p$ 回 & パラメータシフト則と併用．評価コスト大だが正確に収束． \\
\textbf{Bayesian Opt (BO)} & 代理モデル (ガウス過程) & 高い & 1 回 (獲得関数) & $p=1, 2$ の低次元で評価回数を極小化しつつ大域最適解を捕捉． \\
\bottomrule
\end{tabular}
}
\end{table}

\subsubsection*{1. COBYLA (Constrained Optimization BY Linear Approximation) \cite{powell1994direct}}
Qiskit等の量子計算ライブラリにおいてデフォルトで広く採用されている勾配不要（Derivative-Free）最適化手法である．
探索空間の $2p+1$ 個の頂点からなる単体（Simplex）を構築し，各頂点での目的関数値から目的関数および制約条件を局所線形補間（平面近似）する．そして，現在の信頼領域（Trust Region）の球内で線形計画問題を解いて次の一手を決定する（図 \ref{fig:classical_optimizers}(a)）．

\textbf{長所}: 勾配を計算しないため反復ごとの評価回数が $2p+1$ 回以下と非常に少なく，状態ベクトル計算などのノイズのない理想シミュレータ環境では極めて素早く極小値へ収束する．\\
\textbf{短所}: 局所解からの脱出能力を持たず，初期値依存性が極めて高い．また，ショットノイズが存在すると頂点の高低関係が誤判定され，単体が不自然に収縮（縮退）して探索が途中で早期停止する致命的欠点を持つ．

\subsubsection*{2. SPSA (Simultaneous Perturbation Stochastic Approximation) \cite{spall1992multivariate}}
実機ハードウェアやショットノイズが支配的な環境において，圧倒的な堅牢性と実用性を誇る確率的最適化手法である．
通常の有限差分勾配法では $2p$ 個のパラメータの勾配を求めるために各座標軸を個別に揺らし，ステップあたり $2 \times 2p = 4p$ 回の測定を要する．これに対し，Spall (1992) のSPSAは，\textbf{全 $2p$ 個のパラメータを同時にランダムな独立二値符号 $\Delta_k \in \{+1, -1\}$ で微小摂動}させる（図 \ref{fig:classical_optimizers}(b)）：
\begin{equation}
    \vec{\theta}_k^\pm = \vec{\theta}_k \pm c_k \vec{\Delta}_k
\end{equation}
そして，勾配の第 $j$ 成分を以下の単一の差分式で推定する：
\begin{equation}
    \hat{g}_{k, j}(\vec{\theta}_k) = \frac{\hat{F}(\vec{\theta}_k + c_k \vec{\Delta}_k) - \hat{F}(\vec{\theta}_k - c_k \vec{\Delta}_k)}{2 c_k \Delta_{k, j}}
\end{equation}
ここで $\mathbb{E}[\Delta_{k, j} / \Delta_{k, l}] = \delta_{jl}$ という直交統計性により，全成分を同時に揺らしているにもかかわらず，期待値 $\mathbb{E}[\hat{\vec{g}}_k]$ は真の勾配 $\nabla F(\vec{\theta}_k)$ の不偏推定量（高次誤差 $\mathcal{O}(c_k^2)$ を除く）に厳密に一致する．

パラメータ更新は，学習率 $a_k$ を用いて確率的勾配降下法（SGD）形式で行われる：
\begin{equation}
    \vec{\theta}_{k+1} = \vec{\theta}_k - a_k \hat{\vec{g}}_k, \quad a_k = \frac{a}{(k + 1 + A)^\alpha}, \quad c_k = \frac{c}{(k + 1)^\gamma}
\end{equation}
SPSAの決定的な強みは以下の2点にある：
\begin{enumerate}
    \item \textbf{実行コストの次元独立性}: パラメータ数 $2p$ がどれだけ増大しても，\textbf{1反復あたりの量子回路実行回数が常に厳密に2回（$\vec{\theta}^+$ と $\vec{\theta}^-$）で固定}される．
    \item \textbf{ショットノイズの統計的平滑化}: ランダムな摂動と学習率減衰により，ショットノイズが確率的勾配探索の「自然な温度揺らぎ」として機能し，浅い局所解のトラップを自然にすり抜けて大域的解へと前進する．
\end{enumerate}

\subsubsection*{3. パラメータシフト則 (Parameter-Shift Rule) による厳密勾配 \cite{mitarai2018quantum,schuld2019evaluating}}
量子回路の期待値を解析的に微分する画期的な数理手法がパラメータシフト則である．
大学の微分積分学で学ぶ有限差分 $\lim_{\epsilon \to 0} \frac{F(\theta+\epsilon)-F(\theta-\epsilon)}{2\epsilon}$ とは異なり，\textbf{角度を $\pm \frac{\pi}{4}$（または $\pm \frac{\pi}{2}$）シフトさせたマクロな2点の測定値の差}から，近似誤差なく厳密な解析導関数が得られる（図 \ref{fig:classical_optimizers}(c)）．

\begin{insightbox}[【大学3年生向け徹底導出】パラメータシフト則の厳密代数証明]
生成子 $G = \frac{1}{2} P$（$P$ はパウリ行列 $X, Y, Z$ またはパウリ積，固有値 $\pm 1$ により $P^2 = I$）によって駆動される1パラメータユニタリゲート $U(\theta) = \exp(-i \frac{\theta}{2} P)$ を考える．
オイラーの公式と同様に，パウリ行列の冪乗性 $P^2 = I$ を用いてテイラー展開すると：
\begin{equation}
    U(\theta) = \sum_{m=0}^\infty \frac{(-i \frac{\theta}{2})^m}{m!} P^m = \cos\frac{\theta}{2} I - i \sin\frac{\theta}{2} P
\end{equation}
任意のエルミート観測量 $H$ に対し，ハイゼンベルク描像における時間発展演算子 $H(\theta) = U^\dagger(\theta) H U(\theta)$ を展開する：
\begin{align}
    H(\theta) &= \left( \cos\frac{\theta}{2} I + i \sin\frac{\theta}{2} P \right) H \left( \cos\frac{\theta}{2} I - i \sin\frac{\theta}{2} P \right) \nonumber \\
    &= \cos^2\frac{\theta}{2} H + \sin^2\frac{\theta}{2} P H P + i \sin\frac{\theta}{2}\cos\frac{\theta}{2} (P H - H P) \nonumber \\
    &= \frac{1 + \cos\theta}{2} H + \frac{1 - \cos\theta}{2} P H P + \frac{\sin\theta}{2} i [P, H] \nonumber \\
    &= \frac{H + P H P}{2} + \frac{H - P H P}{2} \cos\theta + \frac{i[P, H]}{2} \sin\theta
\end{align}
任意の初期状態 $\ket{\psi_0}$ に対する期待値 $F(\theta) = \mel{\psi_0}{H(\theta)}{\psi_0}$ は，$\theta$ に依存しない定数 $A, B, C$ を用いて厳密に以下のように単一の三角関数波として表される：
\begin{equation}
    F(\theta) = A + B \cos\theta + C \sin\theta
\end{equation}
この関数の $\theta$ に関する解析導関数を普通に微分して求めると：
\begin{equation}
    \frac{d F(\theta)}{d \theta} = -B \sin\theta + C \cos\theta
\end{equation}
ここで，マクロな角度シフト $\pm \frac{\pi}{2}$ を与えたときの三角関数の加法定理に着目する：
\begin{align}
    \cos\left(\theta \pm \frac{\pi}{2}\right) = \mp \sin\theta, \quad \sin\left(\theta \pm \frac{\pi}{2}\right) = \pm \cos\theta
\end{align}
したがって，シフトさせた2点の測定値の差を計算すると：
\begin{align}
    F\left(\theta + \frac{\pi}{2}\right) - F\left(\theta - \frac{\pi}{2}\right) &= B (-\sin\theta - \sin\theta) + C (\cos\theta - (-\cos\theta)) \nonumber \\
    &= 2 \left( -B \sin\theta + C \cos\theta \right) = 2 \frac{d F(\theta)}{d \theta}
\end{align}
両辺を 2 で割ることにより，一切の近似を含まない厳密等式：
\begin{equation}
    \frac{d F(\theta)}{d \theta} = \frac{F\left(\theta + \frac{\pi}{2}\right) - F\left(\theta - \frac{\pi}{2}\right)}{2}
\end{equation}
が完全に証明された！ （なお，ハミルトニアン時間発展が $e^{-i\theta P}$ と定義されている場合は，角度シフト幅が $\pm \frac{\pi}{4}$ となる．）
\end{insightbox}

有限差分と異なり，シフト角が $\frac{\pi}{2}$ という有限の大きさを持つため，分母が微小量にならずショットノイズによる勾配爆発が一切生じない．
各パラメータあたり2回の回路実行（$2p$ 個のパラメータ全体で $4p$ 回）を要するが，古典深層学習で実績のある \textbf{Adam} や \textbf{L-BFGS} などの高精度勾配法を量子系にそのまま接続できる．

\subsubsection*{4. ベイズ最適化 (Bayesian Optimization: BO)}
目的関数の評価回数を極限まで減らしたい場合に用いられるブラックボックス最適化手法である．
これまでの測定履歴 $\mathcal{D} = \{(\vec{\theta}_i, y_i)\}$ からガウス過程回帰（GPR）を用いて目的関数ランドスケープの平均 $\mu(\vec{\theta})$ と不確実性（分散）$\sigma^2(\vec{\theta})$ をモデル化する．
そして，未知領域の探索（Exploration）と既知領域の活用（Exploitation）のトレードオフを最適化する獲得関数（Expected Improvement: EI など）を最大化する点を次の評価角度として選定する．
$p=1, 2$ のような低次元空間において，わずか数十回のQPU呼び出しで大域的最適解を捉える極めて強力な手段となる．

\subsection{3. 角度最適化を成功させる3大先端戦略}
浅い層から深層へとQAOAを展開するにあたり，ランダム初期化に頼った局所探索は $p \ge 3$ 程度で容易に破綻する．現代のQAOA研究では，断熱物理と数理統計の性質を活用した以下の3大戦略が確立されている（図 \ref{fig:interp_strategy}）．

\begin{figure}[htbp]
\centering
\begin{tikzpicture}[scale=0.85, >=latex]
    % Left side: INTERP recursive flow
    \begin{scope}[shift={(0,0)}]
        \draw[fill=blue!3, draw=blue!35, rounded corners=3mm] (-0.2, -0.7) rectangle (6.6, 5.0);
        \node[font=\scriptsize\bfseries, fill=blue!15, text=blue!90!black, rounded corners=1mm, inner sep=3pt] at (3.2, 4.65) {INTERP 再帰的補間フロー};
        
        % Level p=1
        \node[draw=blue!60!black, fill=white, thick, rounded corners=2mm, text width=5.6cm, align=center, font=\scriptsize, inner sep=3pt] at (3.2, 3.75) {
            \textbf{深さ $p=1$}: 大域探索（BO / グリッド探索）\\
            $\to$ パラメータ $(\gamma_1^*, \beta_1^*)$ を高精度に同定
        };
        
        % Arrow to p=2
        \draw[->, very thick, purple!80!black] (3.2, 3.2) -- (3.2, 2.65) 
            node[midway, right=1mm, font=\tiny\bfseries, text=purple!80!black] {線形補間公式 (式 \ref{eq:interp})};
            
        % Level p=2
        \node[draw=blue!60!black, fill=white, thick, rounded corners=2mm, text width=5.6cm, align=center, font=\scriptsize, inner sep=3pt] at (3.2, 2.1) {
            \textbf{深さ $p=2$}: 補間初期値から局所最適化\\
            $(\vec{\gamma}^{(0)}, \vec{\beta}^{(0)}) \longrightarrow (\vec{\gamma}_{(2)}^*, \vec{\beta}_{(2)}^*)$
        };
        
        % Arrow to p+1
        \draw[->, very thick, purple!80!black] (3.2, 1.55) -- (3.2, 1.0) 
            node[midway, right=1mm, font=\tiny\bfseries, text=purple!80!black] {再帰的拡張 ($p \to p+1$)};
            
        % Level p=3
        \node[draw=teal!70!black, fill=teal!5, thick, rounded corners=2mm, text width=5.6cm, align=center, font=\scriptsize, inner sep=3pt] at (3.2, 0.45) {
            \textbf{深さ $p+1$}: 常に断熱引き込み領域（Basin）から開始\\
            \textcolor{teal!80!black}{\textbf{局所解トラップ・Barren Plateau を完全回避}}
        };
        
        \node[font=\tiny, align=center, text=blue!85!black] at (3.2, -0.4) {深層化に伴う局所探索の破綻を克服する現代標準戦略};
    \end{scope}

    % Right side: Smooth Adiabatic Angle Trajectories
    \begin{scope}[shift={(7.5,0)}]
        \draw[fill=gray!4, draw=gray!35, rounded corners=3mm] (-0.3, -0.7) rectangle (6.7, 5.0);
        \node[font=\scriptsize\bfseries, fill=gray!20, text=gray!90!black, rounded corners=1mm, inner sep=3pt] at (3.2, 4.65) {層進行に伴う断熱的パラメータ軌道};
        
        % Axes
        \draw[->, thick] (0.3, 0.6) -- (6.3, 0.6) node[right=1mm, font=\tiny] {層の進行 $k/p$};
        \draw[->, thick] (0.3, 0.6) -- (0.3, 4.2);
        \node[right, font=\tiny] at (0.3, 4.1) {角度の大きさ};
        
        % Ticks on x-axis
        \draw (0.3, 0.5) -- (0.3, 0.7) node[below=1.5mm, font=\tiny] {$0$};
        \draw (3.2, 0.5) -- (3.2, 0.7) node[below=1.5mm, font=\tiny] {$0.5$};
        \draw (6.1, 0.5) -- (6.1, 0.7) node[below=1.5mm, font=\tiny] {$1.0$};
        
        % Gamma trajectory (increasing)
        \draw[very thick, red!85!black, smooth] 
            plot coordinates {(0.3, 1.0) (1.6, 1.4) (3.2, 2.2) (4.7, 3.2) (6.1, 3.8)};
        \node[above left, font=\scriptsize\bfseries, text=red!85!black] at (6.1, 3.8) {$\gamma_k^*$ (コスト層)};
        
        % Beta trajectory (decreasing)
        \draw[very thick, blue!85!black, smooth] 
            plot coordinates {(0.3, 3.7) (1.6, 3.1) (3.2, 2.2) (4.7, 1.5) (6.1, 1.1)};
        \node[above left, font=\scriptsize\bfseries, text=blue!85!black] at (6.1, 1.15) {$\beta_k^*$ (ミキサー層)};
        
        % Physical interpretation annotations
        \draw[<->, dashed, gray!70] (0.7, 1.2) -- (0.7, 3.5);
        \node[right, font=\tiny, align=left, text=gray!80!black] at (0.7, 2.35) {初期: ミキサー優勢\\(量子揺らぎ大)};
        
        \draw[<->, dashed, gray!70] (5.7, 1.3) -- (5.7, 3.6);
        \node[left, font=\tiny, align=right, text=gray!80!black] at (5.7, 2.45) {終期: コスト優勢\\(基底状態へ凝縮)};
        
        \node[font=\tiny, align=center, text width=6.5cm, text=gray!85!black] at (3.2, -0.4) {AQCの時間発展スケジュール $A(t) = 1-t/T, B(t)=t/T$ と数理的一致};
    \end{scope}
\end{tikzpicture}
\caption{INTERP戦略による層漸増型最適化と断熱的パラメータ軌道．(左) 深さ $p$ の最適化解から線形補間によって深さ $p+1$ の初期パラメータを生成し，常に大域的最適解の近傍（引き込み領域）から局所探索を開始する再帰的アルゴリズム．(右) 深層QAOAにおける最適角度系列 $(\gamma_k^*, \beta_k^*)$ は層番号 $k/p$ に対して滑らかな単調曲線を描き，アナログ断熱量子計算（AQC）の時間発展スケジュールと驚くべき数理的一致を見せる．}
\label{fig:interp_strategy}
\end{figure}

\subsubsection*{1. 断熱スケジューリング初期化 (LINEAR Heuristic)}
ランダム初期化は Barren Plateau や局所解トラップの主因となる．
これを避ける最も素直で強力なアプローチは，断熱量子計算（AQC）の時間発展スケジュール $H(t) = (1 - t/T)H_M + (t/T)H_C$ の離散化軌道を模倣する線形スケジューリングである：
\begin{equation}
    \gamma_k^{(0)} = \frac{k}{p} \Delta t, \quad \beta_k^{(0)} = \left(1 - \frac{k}{p}\right) \Delta t \quad (k = 1, \dots, p)
\end{equation}
ここで $\Delta t$ は全体の時間幅を層数 $p$ で割った刻み幅に相当する定数である．
この初期化を用いることで，探索開始時点からパラメータベクトルが「断熱経路の引き込み領域（Basin of Attraction）」の内部に配置されるため，単純な局所オプティマイザであっても高確率で大域的最適解へとスムーズに案内される．

\subsubsection*{2. 補間ヒューリスティック (INTERP Strategy: Zhou et al., 2020) \cite{zhou2020quantum}}
Zhouら（PRX, 2020）が提唱した INTERP戦略は，深さ $p$ の最適化結果を用いて深さ $p+1$ の初期値を構築する再帰的アルゴリズムである（図 \ref{fig:interp_strategy}(左)）．
数値実験および理論解析により，最適化されたパラメータ系列 $(\vec{\gamma}_{(p)}^*, \vec{\beta}_{(p)}^*)$ は，層番号 $k$ の進行に対して滑らかな単調曲線を描くことが判明している（図 \ref{fig:interp_strategy}(右)）．
そこで，$p$ レイヤーの最適解を 1 次元線形補間（Interpolation）することにより，$p+1$ レイヤーの初期値を以下のように生成する：
\begin{align}
    \gamma_{k}^{(p+1)} &= \frac{k-1}{p} \gamma_{k-1}^{*(p)} + \frac{p+1-k}{p} \gamma_{k}^{*(p)} \quad (k = 1, \dots, p+1) \label{eq:interp} \\
    \beta_{k}^{(p+1)} &= \frac{k-1}{p} \beta_{k-1}^{*(p)} + \frac{p+1-k}{p} \beta_{k}^{*(p)} \quad (k = 1, \dots, p+1)
\end{align}
（端点条件: $\gamma_0^* = 0, \gamma_{p+1}^* = \gamma_p^*, \beta_0^* = \beta_1^*, \beta_{p+1}^* = 0$）．
これにより，$p=1 \to 2 \to 3 \to \dots$ と回路深度を段階的に深めていく際に，\textbf{常に大域的最適解のすぐ近傍から古典最適化を開始できるため，高次元空間での局所解トラップと勾配消失を完全に回避}できる．

\subsubsection*{3. パラメータ集中現象と転移学習 (Transferability) \cite{brandao2018fixed,galda2021transferability,egger2021warm}}
Brandaoら (2018)\cite{brandao2018fixed} およびGaldaら (2021)\cite{galda2021transferability} は，同一の問題クラス（例えばランダム3正則グラフのMaxCutなど）からサンプリングされた異なる問題インスタンス間において，\textbf{QAOAの最適角度 $(\vec{\gamma}^*, \vec{\beta}^*)$ が問題サイズ $N \to \infty$ の極限で特定の定数値に集中する現象（Concentration of Parameters）}を数学的に発見した．

この物理的理由は，第 4 章で詳解する「局所因果律（Light Cone）」にある．回路の深さ $p$ が有限である限り，各量子ビットから伝播する相関の届く範囲（Lieb-Robinson 光円錐）はグラフ距離 $p$ 以内の局所部分グラフに限られる．同等の局所木構造を持つグラフ群であれば，全体の頂点数 $N$ が 10 であろうと 10,000 であろうと，局所期待値関数はほぼ完全に同一の形状をとるのである．

この現象は，実機ハードウェアにおける最適化コストを劇的に削減する「転移学習（Parameter Transfer）」を可能にする：
\begin{enumerate}
    \item 小規模な問題（$N = 12 \sim 16$）において，シミュレータ上で古典オプティマイザを十分に回して最適角度 $(\vec{\gamma}^*, \vec{\beta}^*)$ を高精度に特定する．
    \item 大規模な実問題（$N = 1,000$ 超）を実機QPUで実行する際，\textbf{古典の反復最適化ループを一切スキップ}し，小規模系で得た最適角度を直接セットして単発（または少数）ショット測定を行う．
    \item これにより，実機通信遅延とサンプリング反復に起因する計算コストを $\mathcal{O}(1)$ に圧縮できる．
\end{enumerate}

\begin{pointbox}[BBOにおける転移学習 (Warm-Start) とFMサロゲートの連係]
本研究のブラックボックス最適化（BBO）では，能動学習イテレーション $t \to t+1$ において新規データ点に基づきFMサロゲートモデルが微小更新される．
ハミルトニアンの変動 $\delta H_C$ が小さい場合，最適角度の変動もリプシッツ連続性により微小有界（$\|\theta_{t+1}^* - \theta_t^*\| \le C \|\delta H_C\|$）となる．
したがって，前イテレーションの最適角度 $\theta_t^*$ を次ステップの初期値（Warm-Start）として転移することで，古典最適化の反復回数を劇的に削減し，リアルタイム探索パイプラインを実現できる．
\end{pointbox}

\subsection{4. 古典オプティマイザ選定の実践的ガイドライン}
実際の研究・開発環境において，どのオプティマイザを選択すべきかの実践的指針を以下にまとめる：
\begin{itemize}
    \item \textbf{理想シミュレータ環境（状態ベクトル，ショットノイズなし）}:
          \begin{itemize}
              \item $p \le 2$: グリッド探索，またはベイズ最適化（BO）で大域解を完全に網羅．
              \item $p \ge 3$: LINEAR初期化またはINTERP戦略で初期値を設定後，\textbf{COBYLA} または \textbf{L-BFGS（厳密勾配）} により高速収束．
          \end{itemize}
    \item \textbf{実機QPU / ショットシミュレータ（ショットノイズ支配環境）}:
          \begin{itemize}
              \item \textbf{SPSA} を第一選択とする．パラメータ数 $2p$ に依らず評価回数が常に2回で固定され，ショットノイズを確率的更新で克服できるため，ノイズ環境で最も堅牢である．
              \item 目的関数には平均エネルギーではなく \textbf{CVaR目的関数} を適用し，基底状態のサンプリング確率を尖らせる．
              \item 同一クラスの反復問題では，小規模インスタンスで同定した角度を \textbf{転移学習（Warm-Start）} させて実機実行回数を最小化する．
          \end{itemize}
\end{itemize}

\section{原論文における理論的定理}

\begin{theorem}[単調非減少性 (Monotonicity)]
任意のレイヤー数 $p \ge 1$ に対して，以下が成り立つ：
\begin{equation}
    M_p \ge M_{p-1}
\end{equation}
\end{theorem}
\begin{proof}
$p$ レイヤーのパラメータ集合において，第 $p$ 層のパラメータを $\gamma_p = 0, \beta_p = 0$ と設定すると，$U(C, 0) = I, U(B, 0) = I$ となり，状態 $\ket{\vec{\gamma}, \vec{\beta}}$ は厳密に $p-1$ レイヤーの状態と一致する．したがって $p$ レイヤーの探索空間は $p-1$ レイヤーの空間を包含しているため，最大値は広義単調増加する．
\end{proof}

\begin{theorem}[断熱極限における厳密収束性]
レイヤー数 $p \to \infty$ の極限において，以下が成り立つ：
\begin{equation}
    \lim_{p \to \infty} M_p = \max_{z \in \{0, 1\}^n} C(z)
\end{equation}
\end{theorem}
\begin{proof}
断熱量子計算において，実行時間 $T \to \infty$ とし，トロッター分割数 $p = T / \Delta t \to \infty$ と設定する．断熱経路に沿った適切な角度系列 $(\vec{\gamma}_{\text{ad}}, \vec{\beta}_{\text{ad}})$ を選べば，系の状態は $C$ の最大固有状態へ断熱的に遷移する．変分最適化による最大値 $M_p$ はこの断熱系列の性能以上であることが保証されるため，極限値は真の最大値に一致する．
\end{proof}

\section{MaxCut問題と局所因果律（Light Cone）の厳密解析}
Farhi原論文の最大の見所は，3正則グラフ（3-regular graph）のMaxCut問題に対し，$p=1$ のQAOAの性能（近似比）を純粋に数式展開のみで厳密に導出した点にある．

MaxCutの目的関数は各辺 $\langle u, v \rangle \in E$ のカットの和である：
\begin{equation}
    C = \sum_{\langle u, v \rangle \in E} C_{uv}, \quad C_{uv} = \frac{1}{2}(I - Z_u Z_v)
\end{equation}
$p=1$ における辺 $\langle u, v \rangle$ の期待値は，ハイゼンベルク描像を用いて以下のように書ける：
\begin{equation}
    \langle C_{uv} \rangle = \mel{s}{e^{i \gamma C} e^{i \beta B} C_{uv} e^{-i \beta B} e^{-i \gamma C}}{s}
\end{equation}
ここで演算子の交換関係に着目する：
\begin{enumerate}
    \item $e^{i \beta B} = \prod_j e^{i \beta X_j}$ は1量子ビットゲートの積であるため，局所的な台（サポート）を広げない．
    \item $e^{i \gamma C} = \prod_{e \in E} e^{i \gamma C_e}$ のうち，$u$ または $v$ を共有しない辺 $e = \langle w, k \rangle$ ($w, k \notin \{u, v\}$) は $C_{uv}$ および $Z_u, Z_v$ と可換（$[C_e, C_{uv}] = 0$）であり，ユニタリの左右で相殺して消滅する．
\end{enumerate}

\begin{pointbox}[原論文の核心的洞察: 局所因果律 (Light Cone)]
浅いレイヤー $p=1$ において，1つの辺 $\langle u, v \rangle$ の期待値は，\textbf{頂点 $u, v$ からグラフ距離1以内の局所サブグラフ（部分木）の構造のみに完全に依存し，グラフ全体の頂点数 $n$ には一切依存しない！}
\end{pointbox}

3正則グラフにおいて，距離1の近傍構造は以下の2通りに分類される：
\begin{itemize}
    \item \textbf{木構造（Tree）}: $u$ の他の2本の隣接頂点と，$v$ の他の2本の隣接頂点の間に辺が存在しない（ループ長 $\ge 4$）．
    \item \textbf{三角形を含む構造（Triangle）}: $u$ と $v$ が共通の隣接頂点を持つ（閉路長 3）．
\end{itemize}

木構造の場合，パウリ代数の展開により期待値は厳密に以下となる：
\begin{equation}
    \langle C_{uv} \rangle_{\text{tree}} = \frac{1}{2} + \frac{1}{2} \sin(2\beta) \sin\gamma \cos^2\gamma \left[ 1 + \cos(2\beta) \sin^2\gamma \right]
\end{equation}
これをパラメータ $(\gamma, \beta)$ について数値最適化すると：
\begin{equation}
    \gamma^* \approx 0.61, \quad \beta^* \approx 0.39 \quad \implies \quad \langle C_{uv} \rangle \approx 0.6924
\end{equation}
任意の3正則グラフのカット総数の上界は $|E|$ であるため，最悪ケースの木構造トポロジーに対しても近似比 $r = \frac{\langle C \rangle}{|E|} \ge \mathbf{0.6924}$ が保証される\cite{farhi2014quantum}．なお，古典アルゴリズムではGoemans-Williamson (1995)\cite{goemans1995improved}の半正定値計画（SDP）緩和が近似比 0.87856 を達成しているが，Farhiらは最も浅い $p=1$ の量子回路が定数理論保証を持つことを証明し，さらに $p \to \infty$ で厳密解へ収束することを示した．これは，多項式時間の量子アルゴリズムが非自明な理論保証を持つことを実証した歴史的瞬間であった．

\enlargethispage{2.5\baselineskip}
\begin{insightbox}[ハイゼンベルク描像がもたらす幾何学的奇跡]
シュレーディンガー描像では，波動関数 $\ket{\vec{\gamma}, \vec{\beta}}$ は全 $2^n$ 次元の巨大なヒルベルト空間全体に複雑に広がるため，局所期待値の計算は一見絶望的に思える．
しかし，大学の量子力学で学ぶ\textbf{ハイゼンベルク描像（Heisenberg Picture）}に立ち戻り，演算子側に時間発展を背負わせると（$C_{uv}(t) = U^\dagger C_{uv} U$），パウリ交換子 $[A, B]$ の局所的非可換性により，相互作用の伝播はグラフ距離 $p$ 以内の局所部分木（Light Cone）に完全に封じ込められる．これは量子多体系の「Lieb-Robinson限界」の離散版であり，物理学科の学生にとってハイゼンベルク描像の真価を実感できる最高峰の応用例である．
\end{insightbox}

\chapter{制約付き最適化の壁：ペナルティ法の破綻メカニズム}
\enlargethispage{5\baselineskip}
\vspace{-4mm}

\vspace{-2.5mm}
\section{実問題における制約条件}
MaxCutのような「任意のビット列 $z \in \{0, 1\}^n$ が物理的に意味を持つ」無制約問題とは異なり，材料探索，触媒設計，TSP，スケジューリングなどの実産業最適化では，変数群の中から排他的に1つを選択する \textbf{One-Hot制約（ハミング重み1）} が至る所に現れる：
\begin{equation}
    \sum_{j=1}^{d} x_j = 1 \quad (x_j \in \{0, 1\})
\end{equation}

\vspace{-2.5mm}
\section{従来のペナルティ法 (Standard QAOA)}
従来のQUBOやStandard QAOAでは，この等式制約を満たすために，目的関数に大きな正のペナルティ項を付加する\cite{lucas2014ising}：
\begin{equation}
    H_{\text{total}} = H_{\text{obj}} + \lambda H_{\text{penalty}}, \quad H_{\text{penalty}} = \left( \sum_{j=1}^d x_j - 1 \right)^2
\end{equation}

\vspace{-2.5mm}
\section{なぜペナルティ法は破綻するのか？}
しかし，このアプローチは浅いQAOAにおいて\textbf{数学的・原理的に破綻}する．その理由は以下の3大要因による：

\vspace{-2mm}
\begin{alertbox}[ペナルティ法が抱える3大欠陥]
\begin{enumerate}\setlength{\itemsep}{0pt}\setlength{\parskip}{0pt}
    \item \textbf{全状態空間に対する有効解比率の指数関数的消滅}:
          $M$ 個の $d$ 選択肢カテゴリ変数があるとき，総量子ビット数は $N = M \cdot d$，全ヒルベルト空間次元は $2^N = 2^{M d}$ である．
          しかし，制約を満たす有効解の総数は $d^M$ 個に過ぎない．有効解の割合は：
          \begin{equation}
              \text{有効解比率} = \frac{d^M}{2^{Md}} = \left( \frac{d}{2^d} \right)^M \xrightarrow{M, d \to \infty} 0
          \end{equation}
          例えば $d=4, M=4$（$N=16$）の場合，全 $65,536$ 状態に対し有効解は $256$ 個（$0.39\%$）しかない．
    \item \textbf{一様重ね合わせ初期化による無効空間への散逸}:
          Standard QAOAの初期状態 $\ket{s} = \ket{+}^{\otimes N}$ は，全空間を一様に重ね合わせるため，波動関数の初期重みの $99.6\%$ は無効解空間に存在する．
    \item \textbf{ペナルティ障壁による局所解トラップと $\lambda$ のジレンマ}:
          制約違反を防ぐために $\lambda$ を大きくすると，解空間に巨大なエネルギー障壁が聳え立ち，浅い $p=1, 2$ の回路では波動関数がトンネリングできず，局所解に閉じ込められる\cite{choi2008different,hadfield2019from}．逆に $\lambda$ を小さくすると制約を無視して無効解を出力するという深刻な「パラメータ設定のジレンマ」に陥る（実際，本研究のスケーリング実験でも，Standard QAOAは $N \ge 12$ で制約充足率が $1\%$ 前後に墜落し，$N=16$ 以上では探索アルゴリズムとして完全に機能停止することが実証された）．
\end{enumerate}
\end{alertbox}

\chapter{XYミキサーとQAO Ansatz（Hadfield et al., 2017/2019）の完全解読}

\begin{paperbox}[原著論文データ]
\textbf{タイトル}: From the Quantum Approximate Optimization Algorithm to a Quantum Alternating Operator Ansatz\\
\textbf{著者}: Stuart Hadfield, Zhihui Wang, Bryan O'Gorman, Eleanor G. Rieffel, Davide Venturelli, Rupak Biswas (NASA / USRA)\\
\textbf{公表}: Algorithms 2019, 12(2), 34; arXiv:1709.03489
\end{paperbox}

\section{思想の転換：部分空間探索 (Subspace Search)}
Hadfieldらの画期的な着想は，\textbf{「制約をエネルギーペナルティで外から縛るのではなく，量子回路の幾何学的対称性として埋め込み，そもそも無効解空間に一切立ち入れないようにする」}という点にある．
これにより，アルゴリズムの名称は Approximate（近似）から \textbf{Alternating Operator（交互作用素）Ansatz} へと再定義された．

\section{許容部分空間 (Feasible Subspace) の定義}
全ヒルベルト空間 $\mathcal{H} = (\mathbb{C}^2)^{\otimes N}$ のうち，制約条件を満たす計算基底のみが張る線形部分空間を \textbf{許容部分空間 $\mathcal{H}_{\text{feas}}$} と定義する：
\begin{equation}
    \mathcal{H}_{\text{feas}} = \text{span}\left\{ \ket{x} \;\middle|\; \sum_{j=1}^d x_j = 1 \right\} \subset \mathcal{H}
\end{equation}
次元は $\dim \mathcal{H}_{\text{feas}} = d \ll 2^d$ である．

\section{XYミキサーの数学的構造と対称性保存則}

\subsection{1. ハミルトニアンの定義}
2つの量子ビット $i, j$ 間に作用する XYミキサーハミルトニアン $H_{ij}^{XY}$ は以下で定義される：
\begin{equation}
    H_{ij}^{XY} = \frac{1}{2}\left( X_i X_j + Y_i Y_j \right)
\end{equation}
パウリの昇降演算子 $\sigma^+ = \ket{1}\bra{0} = \frac{X + iY}{2}, \; \sigma^- = \ket{0}\bra{1} = \frac{X - iY}{2}$ を用いると：
\begin{equation}
    H_{ij}^{XY} = \sigma_i^+ \sigma_j^- + \sigma_i^- \sigma_j^+ = \ket{01}\bra{10} + \ket{10}\bra{01}
\end{equation}

\subsection{2. 行列表示}
計算基底 $\{\ket{00}, \ket{01}, \ket{10}, \ket{11}\}$ における行列表示は以下の通りである：
\begin{equation}
    H_{ij}^{XY} = \begin{pmatrix}
        0 & 0 & 0 & 0 \\
        0 & 0 & 1 & 0 \\
        0 & 1 & 0 & 0 \\
        0 & 0 & 0 & 0
    \end{pmatrix}
\end{equation}
この演算子は，$\ket{01}$ と $\ket{10}$ の間だけで励起（「1」のビット）を交換（フリップフロップ遷移）させ，$\ket{00}$ や $\ket{11}$ には一切干渉しない（固有値0）．

\begin{insightbox}[物性物理学（XXZスピン模型）との完全な同値性]
物性物理学・凝縮系物理学において，1次元XXZハイゼンベルク模型のハミルトニアンは以下で知られている：
\begin{equation}
    \mathcal{H}_{\text{XXZ}} = J \sum_{i} \left( X_i X_{i+1} + Y_i Y_{i+1} + \Delta Z_i Z_{i+1} \right)
\end{equation}
この第一項・第二項こそがまさにHadfieldらのXYミキサー $H_{ij}^{XY} = \frac{1}{2}(X_i X_j + Y_i Y_j)$ である！
この相互作用はスピンのフリップフロップ（励起の交換 $\ket{\uparrow\downarrow} \leftrightarrow \ket{\downarrow\uparrow}$）を司り，全スピン角運動量の $z$ 成分 $S_z^{\text{tot}} = \frac{\hbar}{2}\sum_k Z_k$ と厳密に可換（$[H^{XY}, S_z^{\text{tot}}] = 0$）である．
すなわち，QAOAにおける「ハミング重み保存則」は，量子力学における「全角運動量 $S_z$ の保存則（連続対称性 $U(1)$ から生じるネーターの定理）」そのものにほかならない．
\end{insightbox}

\begin{theorem}[ハミング重み（局所磁化）の保存]
全スピン $Z$ 演算子（総ハミング重み） $C_{\text{cons}} = \sum_k Z_k$ に対し，以下が成り立つ：
\begin{equation}
    \left[ H_{ij}^{XY}, \; Z_i + Z_j \right] = 0
\end{equation}
\end{theorem}
\begin{proof}
パウリ行列の交換関係 $[X, Z] = -2iY, \; [Y, Z] = 2iX$ より，
\begin{align*}
    [X_i X_j, Z_i] &= -2i Y_i X_j, \quad [X_i X_j, Z_j] = -2i X_i Y_j \\
    [Y_i Y_j, Z_i] &= 2i X_i Y_j, \quad [Y_i Y_j, Z_j] = 2i Y_i X_j
\end{align*}
したがって，
\begin{align*}
    [X_i X_j + Y_i Y_j, \; Z_i + Z_j] &= (-2i Y_i X_j - 2i X_i Y_j) + (2i X_i Y_j + 2i Y_i X_j) = 0
\end{align*}
よって可換である．
\end{proof}

この定理により，$e^{-i \beta H_{ij}^{XY}}$ による時間発展は，\textbf{系のハミング重みを厳密に1のまま保存}する．

\section{初期状態：W状態の直積}
ハミング重み1の等重み重ね合わせ状態は，まさに量子もつれ状態の代表である \textbf{W状態}（Dicke状態 $D_1^d$） $\ket{W_d}$ である\cite{hadfield2019from,bartschi2019deterministic}：
\begin{equation}
    \ket{W_d} = \frac{1}{\sqrt{d}} \sum_{j=1}^d \ket{0 \dots 0 \underbrace{1}_{j\text{-th}} 0 \dots 0}
\end{equation}

2量子ビットの最小W状態 $\ket{W_2} = \frac{\ket{10} + \ket{01}}{\sqrt{2}}$ は，Bärtschi \& Eidenbenz (2019)\cite{bartschi2019deterministic}らの手法を単純化した，図\ref{fig:w2_circuit}に示す回路で決定論的に生成できる．

\begin{figure}[htbp]
\centering
\begin{tikzpicture}[scale=1.0, >=latex]
    \node at (-0.8, 1) {$\ket{0}_0$};
    \node at (-0.8, 0) {$\ket{0}_1$};
    \draw[thick] (-0.3, 1) -- (5.0, 1);
    \draw[thick] (-0.3, 0) -- (5.0, 0);

    % X on q0
    \draw[fill=blue!15, draw=blue!70!black, thick] (0.5, 0.65) rectangle (1.3, 1.35) node[midway, font=\bfseries] {$X$};
    % H on q1
    \draw[fill=green!15, draw=green!60!black, thick] (1.8, -0.35) rectangle (2.6, 0.35) node[midway, font=\bfseries] {$H$};

    % CNOT (control q1, target q0)
    \filldraw[black] (3.6, 0) circle (3pt);
    \draw[thick] (3.6, 0) -- (3.6, 1);
    \draw[thick, fill=white] (3.6, 1) circle (5pt);
    \draw[thick] (3.6, 0.82) -- (3.6, 1.18);
    \draw[thick] (3.42, 1) -- (3.78, 1);

    \node[right, font=\large] at (5.2, 0.5) {$\ket{W_2} = \dfrac{\ket{10} + \ket{01}}{\sqrt{2}}$};
\end{tikzpicture}
\caption{2量子ビットW状態 $\ket{W_2}$ の生成回路．$X$ で励起を作り，$H$ で重ね合わせ，CNOT で排他的励起状態を共有する．測定結果は必ず $\ket{10}$ または $\ket{01}$ となり，ハミング重み1（One-Hot制約）が完全に満たされる．}
\label{fig:w2_circuit}
\end{figure}

$M$ 個のカテゴリ変数がある場合，系全体の初期状態は各ブロックのW状態の直積（テンソル積）として準備される：
\begin{equation}
    \ket{\psi_0} = \bigotimes_{m=1}^M \ket{W_{d_m}} \;\in\; \mathcal{H}_{\text{feas}}
\end{equation}
初期状態が100\%有効解空間にあり，時間発展演算子も有効解空間を不変に保つため，\textbf{測定結果は数学的に100.0\%の確率で制約を厳密充足することが保証される（ペナルティ項 $\lambda = 0$）}．
図\ref{fig:subspace_comparison}に，従来のStandard QAOAと提案手法FM-XY-QAOAの探索空間の幾何学的相違を示す．

\begin{figure}[htbp]
\centering
\begin{tikzpicture}[scale=0.88, >=latex]
    % Left side: Standard QAOA (x: -0.1 to 6.3)
    \draw[fill=red!3, draw=red!70!black, thick, rounded corners=2mm] (-0.1, 0) rectangle (6.3, 5.2);
    % Title bar
    \draw[fill=red!15, draw=red!70!black, thick, rounded corners=2mm] (-0.1, 4.4) rectangle (6.3, 5.2)
        node[midway, font=\bfseries\small, text=red!85!black] {Standard QAOA (全空間探索)};
    
    % Gray dashed invalid space
    \draw[fill=gray!12, draw=gray!60, dashed] (0.5, 0.5) rectangle (5.7, 4.1);
    \node[font=\scriptsize\bfseries, text=gray!80!black] at (3.1, 3.75) {全状態空間 $2^N$ (指数爆発)};
    
    % Feasible solution dots
    \filldraw[blue!75!black] (1.8, 1.3) circle (3pt) node[below, font=\scriptsize] {$\ket{1001}$};
    \filldraw[blue!75!black] (4.4, 1.3) circle (3pt) node[below, font=\scriptsize] {$\ket{0110}$};
    \filldraw[blue!75!black] (1.8, 2.7) circle (3pt) node[above, font=\scriptsize] {$\ket{1010}$};
    \filldraw[red!85!black] (4.4, 2.7) circle (4pt) node[above, font=\scriptsize\bfseries] {$\ket{0101}^*$};

    % Wasteful dispersal arrows (Inside outer box)
    \draw[->, thick, gray!70, dashed] (3.1, 2.0) -- (1.5, 2.0);
    \node[above, font=\tiny\bfseries, text=red!75!black] at (1.5, 2.05) {違反 $\ket{0000}$};
    
    \draw[->, thick, gray!70, dashed] (3.1, 2.0) -- (4.7, 2.0);
    \node[above, font=\tiny\bfseries, text=red!75!black] at (4.7, 2.05) {違反 $\ket{1111}$};

    \node[below, font=\footnotesize, text width=6.0cm, align=center] at (3.1, -0.2) {
        ペナルティ障壁により波動関数が無効空間（$99\%$以上）に散逸
    };

    % Right side: FM-XY-QAOA (x: 7.1 to 13.5)
    \draw[fill=green!3, draw=green!65!black, thick, rounded corners=2mm] (7.1, 0) rectangle (13.5, 5.2);
    % Title bar
    \draw[fill=green!15, draw=green!65!black, thick, rounded corners=2mm] (7.1, 4.4) rectangle (13.5, 5.2)
        node[midway, font=\bfseries\small, text=green!75!black] {FM-XY-QAOA (部分空間探索)};

    % Subspace Ring
    \draw[thick, blue!65!black, dashed] (10.3, 2.1) circle (1.5cm);
    % Subspace label at center of circle with clear background
    \node[font=\scriptsize\bfseries, text=blue!85!black, fill=white, draw=blue!40, rounded corners=1mm, inner sep=2pt] at (10.3, 2.1) {許容部分空間 $\mathcal{H}_{\text{feas}}$};

    % Feasible dots along ring (R=1.5cm)
    \filldraw[blue!75!black] (8.8, 2.1) circle (3pt) node[left=2pt, font=\scriptsize] {$\ket{1001}$};
    \filldraw[blue!75!black] (11.8, 2.1) circle (3pt) node[right=2pt, font=\scriptsize] {$\ket{0110}$};
    \filldraw[blue!75!black] (10.3, 3.6) circle (3pt) node[above=2pt, font=\scriptsize] {$\ket{1010}$};
    \filldraw[red!85!black] (10.3, 0.6) circle (4pt) node[below=2pt, font=\scriptsize\bfseries] {$\ket{0101}^*$};

    % XY mixer rotation arrows along ring
    \draw[->, very thick, green!60!black] (9.2, 2.7) arc (145:115:1.5cm) node[midway, above left, font=\tiny\bfseries] {XY};
    \draw[->, very thick, green!60!black] (10.9, 3.5) arc (65:35:1.5cm);
    \draw[->, very thick, green!60!black] (11.4, 1.5) arc (-35:-65:1.5cm);
    \draw[->, very thick, green!60!black] (9.7, 0.7) arc (-115:-145:1.5cm);

    \node[below, font=\footnotesize, text width=6.0cm, align=center] at (10.3, -0.2) {
        無効空間を物理排除し，有効解の間だけをトンネリング遷移
    };
\end{tikzpicture}
\caption{探索空間の幾何学的対比．(左) Standard QAOAは広大な無効空間を含む全空間 $2^N$ を探索するため，大半の確率が無効解に散逸する．(右) 提案手法 FM-XY-QAOA は，初期状態とXYミキサーの対称性により許容部分空間 $\mathcal{H}_{\text{feas}}$ に閉じるため，制約充足率 100\% を維持しながら高速に最適解 $\ket{x^*}$ へ集中する．}
\label{fig:subspace_comparison}
\end{figure}

\chapter{XYミキサーの数理解析とNISQ実装論（Wang et al., 2020）}

\begin{paperbox}[原著論文データ]
\textbf{タイトル}: XY-mixers: Analytical and numerical results for QAOA\\
\textbf{著者}: Zhihui Wang, Nicholas C. Rubin, Jason M. Dominy, Eleanor G. Rieffel\\
\textbf{公表}: Physical Review A 101, 012320 (2020); arXiv:1904.09314
\end{paperbox}

\section{時間発展ユニタリ演算子の厳密形}
2量子ビットXYミキサーの時間発展 $U_{ij}^{XY}(\beta) = e^{-i \beta H_{ij}^{XY}}$ は，基底 $\{\ket{00}, \ket{01}, \ket{10}, \ket{11}\}$ において厳密に以下の行列形を持つ：
\begin{equation}
    U_{ij}^{XY}(\beta) = \begin{pmatrix}
        1 & 0 & 0 & 0 \\
        0 & \cos\beta & -i\sin\beta & 0 \\
        0 & -i\sin\beta & \cos\beta & 0 \\
        0 & 0 & 0 & 1
    \end{pmatrix}
\end{equation}
これは量子回路における \textbf{iSWAP ゲート族} や Google Sycamore の \textbf{fSim ゲート} に対応する．
$[X_i X_j, Y_i Y_j] = 0$ であるため，この時間発展演算子はトロッター誤差なく連続する相互作用回転ゲートへと厳密に分解される：
\begin{equation}
    U_{ij}^{XY}(\beta) = \exp\left(-i \frac{\beta}{2} X_i X_j\right) \exp\left(-i \frac{\beta}{2} Y_i Y_j\right) = R_{XX}(\beta) R_{YY}(\beta)
\end{equation}
さらに，NISQハードウェアの基本ゲートセット（CNOTゲートおよび単一量子ビット回転）を用いた分解を図\ref{fig:xy_mixer_circuit}に示す．

\begin{figure}[htbp]
\centering
\begin{tikzpicture}[scale=0.90, >=latex]
    % (a) Composite Unitary
    \node at (-0.8, 1) {$\ket{q_i}$};
    \node at (-0.8, 0) {$\ket{q_j}$};
    \draw[thick] (-0.4, 1) -- (2.0, 1);
    \draw[thick] (-0.4, 0) -- (2.0, 0);
    \draw[fill=teal!15, draw=teal!70!black, thick, rounded corners=1.5mm] (0.0, -0.3) rectangle (1.6, 1.3) node[midway, font=\small\bfseries] {$U_{ij}^{XY}(\beta)$};

    \node[font=\Huge] at (2.6, 0.5) {$=$};

    % (b) Rxx and Ryy
    \node at (3.2, 1) {$\ket{q_i}$};
    \node at (3.2, 0) {$\ket{q_j}$};
    \draw[thick] (3.6, 1) -- (7.2, 1);
    \draw[thick] (3.6, 0) -- (7.2, 0);
    \draw[fill=cyan!15, draw=cyan!70!black, thick, rounded corners=1mm] (4.0, -0.3) rectangle (5.2, 1.3) node[midway, font=\scriptsize\bfseries] {$R_{XX}(\beta)$};
    \draw[fill=blue!15, draw=blue!70!black, thick, rounded corners=1mm] (5.7, -0.3) rectangle (6.9, 1.3) node[midway, font=\scriptsize\bfseries] {$R_{YY}(\beta)$};

    \node[font=\Huge] at (7.8, 0.5) {$=$};

    % (c) CNOT decomposed version
    \node at (8.4, 1) {$\ket{q_i}$};
    \node at (8.4, 0) {$\ket{q_j}$};
    \draw[thick] (8.8, 1) -- (14.2, 1);
    \draw[thick] (8.8, 0) -- (14.2, 0);

    % CNOT 1
    \filldraw[black] (9.4, 1) circle (2.5pt);
    \draw[thick] (9.4, 1) -- (9.4, 0);
    \draw[thick, fill=white] (9.4, 0) circle (4.5pt);
    \draw[thick] (9.4, -0.16) -- (9.4, 0.16); \draw[thick] (9.24, 0) -- (9.56, 0);

    % Ry(-beta/2) on qi, Ry(-beta/2) on qj
    \draw[fill=orange!15, draw=orange!70!black, thick, rounded corners=0.8mm] (9.9, 0.7) rectangle (11.3, 1.3) node[midway, font=\tiny\bfseries] {$R_Y(-\frac{\beta}{2})$};
    \draw[fill=orange!15, draw=orange!70!black, thick, rounded corners=0.8mm] (9.9, -0.3) rectangle (11.3, 0.3) node[midway, font=\tiny\bfseries] {$R_Y(-\frac{\beta}{2})$};

    % CNOT 2
    \filldraw[black] (11.8, 1) circle (2.5pt);
    \draw[thick] (11.8, 1) -- (11.8, 0);
    \draw[thick, fill=white] (11.8, 0) circle (4.5pt);
    \draw[thick] (11.8, -0.16) -- (11.8, 0.16); \draw[thick] (11.64, 0) -- (11.96, 0);

    % Ry(beta/2) on qi, Ry(-beta/2) on qj
    \draw[fill=orange!15, draw=orange!70!black, thick, rounded corners=0.8mm] (12.3, 0.7) rectangle (13.7, 1.3) node[midway, font=\tiny\bfseries] {$R_Y(\frac{\beta}{2})$};
    \draw[fill=orange!15, draw=orange!70!black, thick, rounded corners=0.8mm] (12.3, -0.3) rectangle (13.7, 0.3) node[midway, font=\tiny\bfseries] {$R_Y(-\frac{\beta}{2})$};
\end{tikzpicture}
\caption{2量子ビットXYミキサー時間発展ゲート $U_{ij}^{XY}(\beta) = \exp\left(-i \frac{\beta}{2}(X_i X_j + Y_i Y_j)\right)$ の等価ゲート分解．左の複合ユニタリは，中央の $R_{XX}(\beta)$ と $R_{YY}(\beta)$ の連続回転，あるいは右の CNOT と単一量子ビット回転 $R_Y$ による標準回路へと等価変換される．このゲートは基底状態 $\ket{01}$ と $\ket{10}$ の間でのみ回転を起こし，$\ket{00}$ と $\ket{11}$ は不変に保つため，ハミング重み（One-Hot制約）が厳密に保存される．}
\label{fig:xy_mixer_circuit}
\end{figure}

\section{トポロジーの選択：完全グラフ（Clique） vs リング型（Ring）}
1つのブロック内の量子ビット集合 $V = \{1, \dots, d\}$ に対し，ミキサーをどう構成するか：
\begin{enumerate}
    \item \textbf{完全グラフ型ミキサー (Clique Mixer)}:
          全対全結合 $H_M = \sum_{i < j} H_{ij}^{XY}$．
          長所：対称性が高くパラメータ探索が平滑．
          短所：ゲート数が $O(d^2)$ であり，NISQ実機でのオーバーヘッドが大きい．
    \item \textbf{リング型ミキサー (Ring Mixer)}:
          最近接結合のみ $H_M = \sum_{j=1}^{d-1} H_{j, j+1}^{XY} + H_{d, 1}^{XY}$．
          長所：ゲート数がわずか $O(d)$ であり，偶奇順序（Odd-Even ordering）により回路深さ2で並列実行可能！
\end{enumerate}

Wangら\cite{wang2020xy}は，リング型XYミキサーであっても部分空間内の任意の基底間の遷移（エルゴード性・推移性）が完全に保証され，Clique型と同等の高い探索性能を発揮することを数理的・数値的に証明した．さらに，Kivlichanら\cite{kivlichan2018quantum}が提唱した偶奇スワップネットワーク（Odd-Even Transposition Network）を適用することで，線形トポロジーを持つNISQプロセッサ上でも深さ $O(N)$ で完全に並列コンパイル可能であることを示した．

\chapter{Factorization Machine（FM）サロゲートモデルとの完全調和}

\section{なぜFMとXY-QAOAなのか？}
本研究の真骨頂は，機械学習のサロゲートモデルとして \textbf{Factorization Machine (FM)} を選定した点にある．
FM（Rendle, 2010\cite{rendle2010factorization}）は，高次元疎なOne-Hot変数 $x \in \{0, 1\}^N$ に対する予測関数として定義される：
\begin{equation}
    \hat{f}(x) = w_0 + \sum_{i=1}^N w_i x_i + \sum_{i=1}^N \sum_{j=i+1}^N \langle v_i, v_j \rangle x_i x_j
\end{equation}
ここで $v_i \in \mathbb{R}^k$ は潜在因子ベクトルである．

このモデルには，他の機械学習モデル（深層学習やガウス過程）にはない2大特長が存在する：
\begin{enumerate}
    \item \textbf{厳密な2次多項式}: 構造が最初からQUBOそのものであり，\textbf{一切の近似誤差なく 1対1 で量子コストハミルトニアン $H_C$ にマッピング可能}．
    \item \textbf{One-Hotカテゴリカル表現の母体}: FMは元来One-Hot変数の相互作用を解明するために誕生した．
\end{enumerate}

したがって，\textbf{「One-Hot変数を捉えるFM」$\times$「One-Hot部分空間を厳密保存するXY-QAOA」} は，理論的に完璧な対称性の調和を形成する．
ペナルティ項 $\lambda$ が一切存在しないため，BBOの能動学習イテレーションごとにサロゲートモデルが更新されても，従来のFMQA（量子アニーリングによるFM最適化\cite{kitai2020designing}）で不可避であった煩雑なペナルティ再調整が完全に不要となり，\textbf{「完全無人自律探索（Self-Driving Lab）」} が達成される．

\chapter{理解度確認演習問題と解答}

\section{演習問題}

\subsubsection*{【問題 1】パウリ交換関係の証明}
パウリ行列の基本代数 $[X, Z] = -2iY, [Y, Z] = 2iX$ を用いて，
\begin{equation}
    \left[ \frac{1}{2}(X_1 X_2 + Y_1 Y_2), \; Z_1 + Z_2 \right] = 0
\end{equation}
を詳細に示せ．

\subsubsection*{【問題 2】2量子ビットXYユニタリの時間発展}
ハミルトニアン $H = \ket{01}\bra{10} + \ket{10}\bra{01}$ に対し，$H^2, H^3$ を計算し，テイラー展開を用いて $e^{-i \beta H}$ が以下の行列になることを導出せよ：
\begin{equation}
    e^{-i \beta H} = I + (\cos\beta - 1) H^2 - i \sin\beta H
\end{equation}

\subsubsection*{【問題 3】W状態の制約充足率}
3量子ビット系において，$X$ミキサーによる一様状態 $\ket{+}^{\otimes 3}$ と，W状態 $\ket{W_3} = \frac{1}{\sqrt{3}}(\ket{100} + \ket{010} + \ket{001})$ のそれぞれについて，測定したビット列のハミング重みが1となる確率（制約充足率）を計算せよ．

\section{解答と解説}

\subsubsection*{【問題 1 の解答】}
積の交換子公式 $[AB, C] = A[B, C] + [A, C]B$ を適用する．
\begin{align*}
    [X_1 X_2, Z_1 + Z_2] &= [X_1 X_2, Z_1] + [X_1 X_2, Z_2] = [X_1, Z_1] X_2 + X_1 [X_2, Z_2] \\
    &= (-2i Y_1) X_2 + X_1 (-2i Y_2) = -2i (Y_1 X_2 + X_1 Y_2)
\end{align*}
同様に，
\begin{align*}
    [Y_1 Y_2, Z_1 + Z_2] &= [Y_1, Z_1] Y_2 + Y_1 [Y_2, Z_2] \\
    &= (2i X_1) Y_2 + Y_1 (2i X_2) = 2i (X_1 Y_2 + Y_1 X_2)
\end{align*}
両者を足し合わせると：
\begin{equation}
    [X_1 X_2 + Y_1 Y_2, Z_1 + Z_2] = -2i(Y_1 X_2 + X_1 Y_2) + 2i(X_1 Y_2 + Y_1 X_2) = 0 \quad \text{(Q.E.D.)}
\end{equation}

\subsubsection*{【問題 2 の解答】}
$H = \ket{01}\bra{10} + \ket{10}\bra{01}$ を2乗すると：
\begin{align*}
    H^2 &= (\ket{01}\bra{10} + \ket{10}\bra{01})^2 \\
    &= \ket{01}\braket{10}{10}\bra{01} + \ket{10}\braket{01}{01}\bra{10} = \ket{01}\bra{01} + \ket{10}\bra{10}
\end{align*}
したがって，$H^3 = H^2 H = H, H^4 = H^2$ となり，奇数乗は $H$，偶数乗（$\ge 2$）は $H^2$ となる．
指数関数のテイラー展開を偶数項と奇数項に分けると：
\begin{align*}
    e^{-i \beta H} &= \sum_{k=0}^\infty \frac{(-i \beta)^k}{k!} H^k = I + \sum_{m=1}^\infty \frac{(-1)^m \beta^{2m}}{(2m)!} H^2 - i \sum_{m=0}^\infty \frac{(-1)^m \beta^{2m+1}}{(2m+1)!} H \\
    &= I + (\cos\beta - 1) H^2 - i \sin\beta H
\end{align*}
これを基底 $\{\ket{00}, \ket{01}, \ket{10}, \ket{11}\}$ の行列で書き下すと，第6節の行列と完全に一致する．\quad (Q.E.D.)

\subsubsection*{【問題 3 の解答】}
\begin{itemize}
    \item \textbf{一様状態 $\ket{+}^{\otimes 3}$}:
          $\ket{+}^{\otimes 3} = \frac{1}{\sqrt{8}} (\ket{000} + \ket{001} + \dots + \ket{111})$．全8状態のうち，ハミング重み1の基底は $\ket{100}, \ket{010}, \ket{001}$ の3個．
          したがって充足率は $3/8 = \mathbf{37.5\%}$．
    \item \textbf{W状態 $\ket{W_3}$}:
          状態ベクトルを構成する基底が最初からすべてハミング重み1であるため，測定確率は $\frac{1}{3} + \frac{1}{3} + \frac{1}{3} = \mathbf{100.0\%}$．
\end{itemize}

\chapter{まとめ}

本教材では，QAOAの数学的・物理的根底から，Farhiら（2014）\cite{farhi2014quantum}の先駆的理論，Hadfieldら（2019）\cite{hadfield2019from}およびWangら（2020）\cite{wang2020xy}による制約保存型XYミキサーの数理解析，そしてFactorization Machineを用いた実産業最適化への結合\cite{kitai2020designing}までを包括的に解説した．Preskill (2018)\cite{preskill2018quantum}が提唱したNISQ時代の制約を克服し，将来の誤り耐性量子コンピュータ（FTQC）へと繋がる強固な理論的礎石がここにある．

\begin{pointbox}[本教材の要点総括]
\begin{enumerate}\setlength{\itemsep}{0pt}\setlength{\parskip}{0pt}
    \item \textbf{QAOAの原理}: 断熱発展のトロッター分解を変分パラメータ化し，浅い量子回路と古典オプティマイザで最適化するVQA．
    \item \textbf{Farhi原論文の金字塔}: 局所因果律（Light Cone）を看破し，3正則MaxCutで近似比 0.6924 を厳密証明．
    \item \textbf{Standard QAOAの限界}: One-Hot制約付き問題では，スケール増大に伴い無効解空間に波動関数が散逸し破綻．
    \item \textbf{XYミキサーの革命}: フリップフロップ項 $\frac{1}{2}(XX+YY)$ によりハミング重みを厳密保存．ペナルティ完全ゼロ（$\lambda=0$）で制約充足率100\%を達成．
    \item \textbf{産業実装}: One-Hot変数の相互作用を捉えるFMサロゲートモデルとXY-QAOAが三位一体となり，完全無人自律の材料探索BBOを実現する．
\end{enumerate}
\end{pointbox}




\begin{thebibliography}{99}

\bibitem{farhi2014quantum}
E.~Farhi, J.~Goldstone, and S.~Gutmann,
\newblock ``A Quantum Approximate Optimization Algorithm,''
\newblock \textit{arXiv preprint arXiv:1411.4028}, 2014.

\bibitem{hadfield2019from}
S.~Hadfield, Z.~Wang, B.~O'Gorman, E.~G.~Rieffel, D.~Venturelli, and R.~Biswas,
\newblock ``From the Quantum Approximate Optimization Algorithm to a Quantum Alternating Operator Ansatz,''
\newblock \textit{Algorithms}, vol.~12, no.~2, p.~34, 2019. \texttt{arXiv:1709.03489}.

\bibitem{wang2020xy}
Z.~Wang, N.~C.~Rubin, J.~M.~Dominy, and E.~G.~Rieffel,
\newblock ``XY-mixers: Analytical and numerical results for the quantum alternating operator ansatz,''
\newblock \textit{Physical Review A}, vol.~101, no.~1, p.~012320, 2020. \texttt{arXiv:1904.09314}.

\bibitem{lucas2014ising}
A.~Lucas,
\newblock ``Ising formulations of many NP problems,''
\newblock \textit{Frontiers in Physics}, vol.~2, p.~5, 2014.

\bibitem{farhi2000quantum}
E.~Farhi, J.~Goldstone, S.~Gutmann, and M.~Sipser,
\newblock ``Quantum Computation by Adiabatic Evolution,''
\newblock \textit{arXiv preprint quant-ph/0001106}, 2000.

\bibitem{farhi2001quantum}
E.~Farhi, J.~Goldstone, S.~Gutmann, J.~Lapan, A.~Lundgren, and D.~Preda,
\newblock ``A Quantum Adiabatic Evolution Algorithm Applied to Random Instances of an NP-Complete Problem,''
\newblock \textit{Science}, vol.~292, no.~5516, pp.~472--475, 2001.

\bibitem{kadowaki1998quantum}
T.~Kadowaki and H.~Nishimori,
\newblock ``Quantum annealing in the transverse Ising model,''
\newblock \textit{Physical Review E}, vol.~58, no.~5, pp.~5355--5363, 1998.

\bibitem{kirkpatrick1983optimization}
S.~Kirkpatrick, C.~D.~Gelatt~Jr., and M.~P.~Vecchi,
\newblock ``Optimization by Simulated Annealing,''
\newblock \textit{Science}, vol.~220, no.~4598, pp.~671--680, 1983.

\bibitem{born1928beweis}
M.~Born and V.~Fock,
\newblock ``Beweis des Adiabatensatzes,''
\newblock \textit{Zeitschrift f{\"u}r Physik}, vol.~51, no.~3, pp.~165--180, 1928.

\bibitem{kato1950adiabatic}
T.~Kato,
\newblock ``On the Adiabatic Theorem of Quantum Mechanics,''
\newblock \textit{Journal of the Physical Society of Japan}, vol.~5, no.~6, pp.~435--439, 1950.

\bibitem{jansen2007bounds}
S.~Jansen, M.-B.~Ruskai, and R.~Seiler,
\newblock ``Bounds for the adiabatic approximation with applications to quantum computation,''
\newblock \textit{Journal of Mathematical Physics}, vol.~48, no.~10, p.~102111, 2007.

\bibitem{berry1984quantal}
M.~V.~Berry,
\newblock ``Quantal phase factors accompanying adiabatic changes,''
\newblock \textit{Proceedings of the Royal Society of London. Series A}, vol.~392, no.~1802, pp.~45--57, 1984.

\bibitem{landau1932theorie}
L.~D.~Landau,
\newblock ``Zur Theorie der Energie{\"u}bertragung. II,''
\newblock \textit{Physikalische Zeitschrift der Sowjetunion}, vol.~2, pp.~46--51, 1932.

\bibitem{zener1932non}
C.~Zener,
\newblock ``Non-adiabatic crossing of energy levels,''
\newblock \textit{Proceedings of the Royal Society of London. Series A}, vol.~137, no.~833, pp.~696--702, 1932.

\bibitem{altshuler2010anderson}
B.~Altshuler, H.~Krovi, and J.~Roland,
\newblock ``Anderson localization makes adiabatic quantum optimization fail,''
\newblock \textit{Proceedings of the National Academy of Sciences}, vol.~107, no.~28, pp.~12446--12450, 2010.

\bibitem{aharonov2007adiabatic}
D.~Aharonov, W.~van~Dam, J.~Kempe, Z.~Landau, S.~Lloyd, and O.~Regev,
\newblock ``Adiabatic Quantum Computation Is Equivalent to Standard Quantum Computation,''
\newblock \textit{SIAM Journal on Computing}, vol.~37, no.~1, pp.~166--194, 2007.

\bibitem{albash2018adiabatic}
T.~Albash and D.~A.~Lidar,
\newblock ``Adiabatic quantum computation,''
\newblock \textit{Reviews of Modern Physics}, vol.~90, no.~1, p.~015002, 2018.

\bibitem{trotter1959product}
H.~F.~Trotter,
\newblock ``On the product of semi-groups of operators,''
\newblock \textit{Proceedings of the American Mathematical Society}, vol.~10, no.~4, pp.~545--551, 1959.

\bibitem{suzuki1976generalized}
M.~Suzuki,
\newblock ``Generalized Trotter's formula and systematic approximants of exponential operators and inner derivations with applications to many-body problems,''
\newblock \textit{Communications in Mathematical Physics}, vol.~51, no.~2, pp.~183--190, 1976.

\bibitem{lloyd1996universal}
S.~Lloyd,
\newblock ``Universal Quantum Simulators,''
\newblock \textit{Science}, vol.~273, no.~5278, pp.~1073--1078, 1996.

\bibitem{nielsen2010quantum}
M.~A.~Nielsen and I.~L.~Chuang,
\newblock \textit{Quantum Computation and Quantum Information: 10th Anniversary Edition},
\newblock Cambridge University Press, 2010.

\bibitem{barenco1995elementary}
A.~Barenco, C.~H.~Bennett, R.~Cleve, D.~P.~DiVincenzo, N.~Margolus, P.~Shor, T.~Sleator, J.~A.~Smolin, and H.~Weinfurter,
\newblock ``Elementary gates for quantum computation,''
\newblock \textit{Physical Review A}, vol.~52, no.~5, pp.~3457--3467, 1995.

\bibitem{goemans1995improved}
M.~X.~Goemans and D.~P.~Williamson,
\newblock ``Improved approximation algorithms for maximum cut and satisfiability problems using semidefinite programming,''
\newblock \textit{Journal of the ACM}, vol.~42, no.~6, pp.~1115--1145, 1995.

\bibitem{choi2008different}
V.~Choi,
\newblock ``Minor-embedding in adiabatic quantum computation: I. The parameter setting problem,''
\newblock \textit{Quantum Information Processing}, vol.~7, no.~5, pp.~193--209, 2008.

\bibitem{bartschi2019deterministic}
A.~B{\"a}rtschi and S.~Eidenbenz,
\newblock ``Deterministic Preparation of Dicke and W States on SMS Architectures,''
\newblock in \textit{Fundamentals of Computation Theory (FCT 2019)}, Lecture Notes in Computer Science, vol.~11651, pp.~24--39, Springer, 2019. \texttt{arXiv:1904.07358}.

\bibitem{kivlichan2018quantum}
I.~D.~Kivlichan, J.~McClean, Nathan Wiebe, C.~Gidney, A.~Aspuru-Guzik, G.~K.-L.~Chan, and R.~Babbush,
\newblock ``Quantum Simulation of Electronic Structure with Linear Depth and Connectivity,''
\newblock \textit{Physical Review Letters}, vol.~120, no.~11, p.~110501, 2018.

\bibitem{mcclean2018barren}
J.~R.~McClean, S.~Boixo, V.~N.~Smelyanskiy, R.~Babbush, and H.~Neven,
\newblock ``Barren plateaus in quantum neural network training landscapes,''
\newblock \textit{Nature Communications}, vol.~9, p.~4812, 2018.

\bibitem{wang2021noise}
S.~Wang, E.~Fontana, M.~Cerezo, K.~Sharma, A.~Sone, L.~Cincio, and P.~J.~Coles,
\newblock ``Noise-induced barren plateaus in variational quantum algorithms,''
\newblock \textit{Nature Communications}, vol.~12, p.~6961, 2021.

\bibitem{powell1994direct}
M.~J.~D.~Powell,
\newblock ``A Direct Search Optimization Method That Models the Objective and Constraint Functions by Linear Interpolation,''
\newblock in \textit{Advances in Optimization and Numerical Analysis}, pp.~51--67, Springer, 1994.

\bibitem{spall1992multivariate}
J.~C.~Spall,
\newblock ``Multivariate stochastic approximation using a simultaneous perturbation gradient approximation,''
\newblock \textit{IEEE Transactions on Automatic Control}, vol.~37, no.~3, pp.~332--341, 1992.

\bibitem{mitarai2018quantum}
K.~Mitarai, M.~Negoro, M.~Kitagawa, and K.~Fujii,
\newblock ``Quantum circuit learning,''
\newblock \textit{Physical Review A}, vol.~98, no.~3, p.~032309, 2018.

\bibitem{schuld2019evaluating}
M.~Schuld, V.~Bergholm, C.~Gogolin, J.~Izaac, and N.~Killoran,
\newblock ``Evaluating analytic gradients on quantum hardware,''
\newblock \textit{Physical Review A}, vol.~99, no.~3, p.~032331, 2019.

\bibitem{zhou2020quantum}
L.~Zhou, S.-T.~Wang, S.~Choi, H.~Pichler, and M.~D.~Lukin,
\newblock ``Quantum Approximate Optimization Algorithm: Performance, Mechanism, and Implementation on Near-Term Devices,''
\newblock \textit{Physical Review X}, vol.~10, no.~2, p.~021067, 2020.

\bibitem{brandao2018fixed}
F.~G.~S.~L.~Brandao, M.~Broughton, E.~Farhi, S.~Gutmann, and H.~Neven,
\newblock ``For Fixed Control Parameters the Quantum Approximate Optimization Algorithm's Objective Function Value Concentrates for Large Trees,''
\newblock \textit{arXiv preprint arXiv:1812.04170}, 2018.

\bibitem{galda2021transferability}
A.~Galda, X.~Liu, D.~Lykov, Y.~Alexeev, and I.~Safro,
\newblock ``Transferability of Optimal QAOA Parameters in Combinatorial Optimization,''
\newblock \textit{IEEE Transactions on Quantum Engineering}, vol.~2, p.~3102611, 2021.

\bibitem{egger2021warm}
D.~J.~Egger, J.~Mare{\v{c}}ek, and S.~Woerner,
\newblock ``Warm-starting quantum optimization,''
\newblock \textit{Quantum}, vol.~5, p.~479, 2021.

\bibitem{rendle2010factorization}
S.~Rendle,
\newblock ``Factorization Machines,''
\newblock in \textit{2010 IEEE International Conference on Data Mining (ICDM)}, pp.~995--1000, IEEE, 2010.

\bibitem{kitai2020designing}
K.~Kitai, J.~Guo, S.~Ju, S.~Tanaka, K.~Tsuda, J.~Shiomi, and R.~Tamura,
\newblock ``Designing metamaterials with quantum annealing and factorization machines,''
\newblock \textit{Physical Review Research}, vol.~2, no.~1, p.~013319, 2020.

\bibitem{preskill2018quantum}
J.~Preskill,
\newblock ``Quantum Computing in the {NISQ} era and beyond,''
\newblock \textit{Quantum}, vol.~2, p.~79, 2018.

\bibitem{cerezo2021variational}
M.~Cerezo, A.~Arrasmith, R.~Babbush, S.~C.~Benjamin, S.~Endo, K.~Fujii, J.~R.~McClean, K.~Mitarai, X.~Yuan, L.~Cincio, and P.~J.~Coles,
\newblock ``Variational quantum algorithms,''
\newblock \textit{Nature Reviews Physics}, vol.~3, no.~9, pp.~625--644, 2021.

\bibitem{barkoutsos2020improving}
P.~K.~Barkoutsos, G.~Nannicini, A.~Robert, I.~Tavernelli, and S.~Woerner,
\newblock ``Improving Variational Quantum Optimization using {CVaR},''
\newblock \textit{Quantum}, vol.~4, p.~256, 2020.

\end{thebibliography}


\end{document}
